\chapter{Configuration spaces}
\label{chap:config-spaces}

The open configuration space $C_n(X)$ is not compact, and the
propagator $\eta_{ij}=d\log(z_i-z_j)$ does not extend as a smooth
form. The bar construction needs a weaker and sharper object: a
logarithmic extension with residues along the collision boundary. The
Fulton--MacPherson compactification $\overline C_n(X)$ supplies that
object. Its normal-crossing boundary carries holomorphic scale
coordinates for collisions, its affine tangent screens carry the
genus-$0$ Arnold relation, and its residue maps give the geometric
part of the bar coderivation. The chiral product supplies the
algebraic OPE coefficient; the compactification does not by itself
define the bar differential. For $X=\mathbb P^1$ the affine Arnold
presentation must be quotiented by the projective linear relations
coming from the $\mathrm{PSL}_2$ action; for $g(X)\geq1$ it must be
replaced by the surface classes and mixed period relations recorded in
Proposition~\ref{prop:arnold-higher-genus}.

The logarithmic 1-forms $\eta_{ij}=d\log(z_i-z_j)$ are the scalar
coefficients of the universal KZ connection one-form on an affine
genus-$0$ screen,
$\omega_{\mathrm{Arnold}} = \sum_{i<j} t_{ij}\,\eta_{ij}$ valued in
the infinitesimal braid Lie algebra $\mathfrak{t}_n$; the
universal flat connection is
$\nabla_{\mathrm{Arnold}} = d - \omega_{\mathrm{Arnold}}$ on the
trivial $U(\mathfrak t_n)$-bundle over
$\operatorname{Conf}_n(\mathbb C)$. On the Arnold-local ordered
genus-$0$ subcategory, the simple-pole part of the chiral bar
differential is the image of this flat connection under the KZ
specialisation:
\[
 d_{\mathrm{bar}}^{(0),\,\mathrm{simp}}
 \;=\;
 \mathrm{KZ}^{*}\bigl(\nabla_{\mathrm{Arnold}}\bigr)_{\mathrm{simp}} .
\]
The full differential is obtained by adding the finite higher OPE
modes, the Koszul signs from the ordered tensor coalgebra, and the
projective or period terms away from the affine screen. Six inputs
carry the comparison: $\overline{C}_n(X)$
(\S\ref{sec:FM-compactification}); the logarithmic propagator
(Definition~\ref{def:log-forms}); the Arnold relation
(Theorem~\ref{thm:arnold-relations}), which is the scalar form
component of the flatness of $\nabla_{\mathrm{Arnold}}$ on the affine
genus-$0$ screen; the infinitesimal braid relations, which supply the
Lie-algebra component; the
projective/period corrections of
Proposition~\ref{prop:arnold-higher-genus}, which govern the global
curve geometry away from that screen; and the residue maps
(Theorem~\ref{thm:residue-operations}), which extract the logarithmic
form factor while the OPE supplies the chiral operation. The KZ
comparison is proved at the climax
(Chapter~\ref{chap:climax-platonic}) in precisely this affine
simple-pole scope.

For the modular extension (Mok logarithmic FM
\S\ref{sec:log-fm-punctured}; bordered FM
\S\ref{sec:bordered-fm}; period coordinates at higher genus
\S~\ref{sec:normal-crossings-higher-genus}) the Arnold relation
remains the local tangent-screen identity, but the global
configuration-space package acquires projective and period corrections
(Proposition~\ref{prop:arnold-higher-genus};
Proposition~\ref{prop:elliptic-arnold-relations}; \cite{Fay73}
for prime forms): on the scalar modular lane the resulting deficit is
measured by the curvature $\kappa(\cA) \cdot \omega_g$, which the universal-conductor
chapter (Chapter~\ref{chap:universal-conductor}) reads as
$K(\cA) = -c_{\mathrm{ghost}}(\mathrm{BRST}(\cA))$ on the
shadow side.

The homotopy upgrade enters once one takes the boundary stratification
seriously. A codimension-$k$ stratum of $\overline{C}_n(X)$ records
not a pairwise collision but a nested sequence of $k$ simultaneous
degenerations, each carrying its own screen. The resulting operad of
boundary inclusions is the Fulton--MacPherson operad, and the
convolution dg Lie algebra
$\Convstr(\overline{C}_\bullet, \cA)$ is the strict model of an
$\Linf$~convolution algebra $\Convinf(\overline{C}_\bullet, \cA)$
whose Maurer--Cartan elements are twisting morphisms. The dg Lie
bracket on $\Convstr$ captures the leading term of a higher
composition sequence indexed by the full poset of boundary strata.
The strict level suffices for the bar differential and curvature
computation, but the transfer theorem
(Chapter~\ref{chap:bar-cobar-adjunction}) that carries
quasi-isomorphisms through bar--cobar is intrinsically an $\Linf$
statement, and $\overline{C}_n(X)$ is where it lives.

Examples calibrate the theory. The Heisenberg algebra sees only the
top stratum (pairwise collisions, genus~$0$ screens); the affine
algebra $\widehat{\mathfrak{g}}_k$ activates three-fold collisions
producing the cubic shadow; the $\beta\gamma$ system reaches the
quartic contact invariant through four-fold collisions on
$\overline{C}_4(X)$; the Virasoro algebra, whose shadow obstruction
tower is infinite, requires the full depth of the FM stratification.
In each case the boundary geometry (dimensions of strata,
intersection numbers of divisors, clutching maps to
$\overline{\mathcal{M}}_{g,n}$) determines the analytic content of
the algebraic invariants computed in the characteristic datum chapters. The abstract
convolution algebra $\Convstr(\overline{C}_\bullet, \cA)$ of
Chapter~\ref{ch:algebraic-foundations} acquires its geometric
realization here, and the modular characteristic $\kappa(\cA)$ of
Chapter~\ref{chap:higher-genus} appears as a curvature class on
$\overline{C}_2(\Sigma_g)$.

Three results organize the chapter. First, the
Fulton--MacPherson compactification $\overline{C}_n(X)$ is
constructed by iterated blow-up along diagonals
(Theorem~\ref{thm:FM}). A divisor $D_S$ for $|S|\geq2$ lies over the
partial diagonal
$\Delta_S=\{x_i=x_j:i,j\in S\}$; on its open part it is a fibration
over $\overline C_{n-|S|+1}(X)$ whose fibre is the genus-$0$ tangent
screen
\[
 \overline C_{|S|}(\mathbb A^1)/
 \mathrm{Aff}(\mathbb A^1),
\qquad
\mathrm{Aff}(\mathbb A^1)=\mathbb C\rtimes\mathbb C^*.
\]
The quotient removes translation and scale on the tangent line at the
collision point. Pair collisions on a curve have a point screen
$\mathbb P^0$; higher collisions carry genuine projective shape data.
Nested collisions produce screens inside screens, each of genus~$0$,
and the full boundary stratification is indexed by nested sets rather
than by set partitions alone. Second, the logarithmic
$1$-forms $\eta_{ij} = d\log(z_i - z_j)$ extend across the
boundary as sections of
$\Omega^1_{\overline{C}_n(X)}(\log D)$ and satisfy the
\emph{Arnold relation}
(Theorem~\ref{thm:arnold-relations}):
\begin{equation}\tag{$\star$}
\eta_{ij} \wedge \eta_{jk}
+ \eta_{jk} \wedge \eta_{ki}
+ \eta_{ki} \wedge \eta_{ij} = 0
\qquad (i,j,k \text{ distinct}).
\end{equation}
Third, this relation forces the affine simple-pole component of
$d_{\mathrm{residue}}^2$ to vanish
(Theorem~\ref{thm:bar-nilpotency-complete}): the
bar differential at degree~$2$ is a sum of double residues
at compatible codimension-two collision corners, and
the genus-$0$ Arnold relation is the precise local cancellation
identity ensuring that these contributions sum to zero on an affine
screen. The equivalence between the local Arnold relation and
nilpotency of the bar differential is
Theorem~\ref{thm:arnold-iff-nilpotent}; on the affine genus-$0$ screen
it is the only relation on the scalar propagator classes. On
$\mathbb{P}^1$ one must also quotient by the projective linear
relations, and on $\Sigma_g$ with $g \geq 1$ one must adjoin the
surface generators and mixed period relations of
Proposition~\ref{prop:arnold-higher-genus}. The further elliptic and
prime-form corrections
(Proposition~\ref{prop:elliptic-arnold-relations}) measure the
failure of strict nilpotency by the curvature
$\kappa(\cA) \cdot \omega_g$.

\section{Fulton--MacPherson compactification}
\label{sec:FM-compactification}

\index{Verdier duality!on configuration spaces}
The Fulton--MacPherson compactification $\overline{C}_n(X)$ is obtained
by iterated blow-up of $X^n$ along all diagonals, adding boundary
divisors that encode all collision patterns. It is the geometric
carrier of the bar complex $\bar{B}_X(\cA)$, a factorization
\emph{coalgebra} on $\operatorname{Ran}(X)$. Verdier duality on
$\operatorname{Ran}(X)$ exchanges coalgebras and algebras and
realizes the intertwining
\[
\mathbb{D}_{\operatorname{Ran}}\, \bar{B}_X(\cA) \;\simeq\; \bar{B}_X(\cA^!)
\;\simeq\; \cA^!_\infty
\]
of factorization algebras (Theorem~A,
Convention~\textup{\ref{conv:bar-coalgebra-identity}}): the Verdier
dual of the bar coalgebra of $\cA$ is identified with the bar of the
Koszul dual $\cA^!$, equivalently the homotopy Koszul dual algebra
$\cA^!_\infty$. This is \emph{not} the bar--cobar inversion
$\Omega(\bar{B}_X(\cA)) \simeq \cA$ of Theorem~B, which recovers $\cA$
itself; the two are distinct functors on $\bar{B}_X(\cA)$.

\subsection{Explicit construction}

\begin{definition}[Ordered and unordered configuration spaces at genus \texorpdfstring{$g$}{g}]\label{def:config-space-genus-g}
\index{configuration space!ordered|textbf}
\index{configuration space!unordered|textbf}
For a Riemann surface $\Sigma_g$ of genus $g$, the
\emph{ordered configuration space} of $n$~distinct points is
\[
C^{\mathrm{ord}}_n(\Sigma_g) := C_n(\Sigma_g)
\;=\;\{(x_1, \ldots, x_n) \in \Sigma_g^n \mid x_i \neq x_j \text{ for all } i \neq j\},
\]
an open dense subset of $\Sigma_g^n$, with complement the ``fat diagonal''
$\Delta = \bigcup_{i<j} \Delta_{ij}$. The symmetric group $\Sigma_n$ acts
freely on $C^{\mathrm{ord}}_n(\Sigma_g)$ by permutation of indices, and
the \emph{unordered configuration space} is the quotient
\[
C^{\Sigma}_n(\Sigma_g)
\;:=\;
C^{\mathrm{ord}}_n(\Sigma_g)\bigm/ \Sigma_n.
\]
The natural projection
$\pi^\Sigma\colon C^{\mathrm{ord}}_n(\Sigma_g) \to C^{\Sigma}_n(\Sigma_g)$
is a regular $\Sigma_n$-cover. Throughout this chapter the unadorned
symbol $C_n(\Sigma_g)$ refers to the ordered configuration space, in
agreement with the bar-complex conventions of
Chapter~\ref{chap:bar-cobar-adjunction}; the unordered variant is used
only when the symmetric-group quotient is explicitly required.
\end{definition}

\begin{remark}[Ordered vs.\ unordered: where each appears]
\label{rem:ordered-vs-unordered-config}
The two variants serve distinct purposes in the bar--cobar formalism.
\begin{itemize}
\item The \emph{ordered} space $C^{\mathrm{ord}}_n$ is the geometric carrier
 of the $T^c$-style bar complex
 $B^{\mathrm{ord}}(\cA)$ used for chiral algebras with non-symmetric
 structure (the open-string side, $E_1$-type operations).
\item The \emph{unordered} space $C^{\Sigma}_n$ is the geometric carrier
 of the $\mathrm{Sym}^c$-style bar complex $B^{\Sigma}(\cA)$ used for
 factorization \emph{coalgebras} (closed-string, $E_\infty$-type
 operations) where the multiplication is symmetric.
\item The projection $\pi^\Sigma\colon C^{\mathrm{ord}}_n \to C^{\Sigma}_n$
 induces the averaging morphism
 $\mathrm{av}\colon B^{\mathrm{ord}}(\cA) \to B^{\Sigma}(\cA)$
 after the ordered differential has been checked to be
 $\Sigma_n$-equivariant. Algebraically this is passage to
 coinvariants; geometrically it is quotient by relabelling. It is not
 an additional configuration-space symmetry and it does not erase the
 ordered proof of the Arnold cancellation before the quotient.
\end{itemize}
The ordered compactification carries a natural $\Sigma_n$-action, and
the unordered compactification used here is the finite quotient
\[
\overline C^\Sigma_n(\Sigma_g)
\;:=\;
\overline C^{\mathrm{ord}}_n(\Sigma_g)/\Sigma_n .
\]
This quotient is not a second labelled FM blow-up with the same boundary
labels: divisors and deeper strata are identified by relabelling, and
stabilisers appear along symmetric collision patterns. The averaging map
on bar complexes is the corresponding passage to coinvariants after the
ordered differential has been proved $\Sigma_n$-equivariant. Its kernel
is the ordered commutator and ordering data lost in the symmetric shadow.
\end{remark}

\subsection{The Fulton--MacPherson compactification across genera}

The configuration space $C_n(\Sigma_g)$ is non-compact. The Fulton--MacPherson compactification~\cite{FM94} resolves this by iterated blow-ups along diagonals, recording collision rates and angles in exceptional divisors. For a logarithmic variant on punctured and nodal curves, see Mok~\cite{Mok25}.

\subsubsection{Iterated blow-up construction}

\begin{theorem}[Fulton--MacPherson compactification at genus \texorpdfstring{$g$}{g} \cite{FM94}; \ClaimStatusProvedElsewhere]\label{thm:FM}
\index{Fulton--MacPherson compactification|textbf}
\index{blowup!along diagonals}
There exists a canonical smooth compactification $\overline{C}_n(\Sigma_g)$ constructed via iterated blow-ups. More precisely:

\begin{enumerate}
\item There is a natural open embedding
\[j: C_n(\Sigma_g) \hookrightarrow \overline{C}_n(\Sigma_g)\]
with dense image.

\item The compactification $\overline{C}_n(\Sigma_g)$ is smooth and proper over $\mathbb{C}$; as a blow-up of the smooth projective variety~$\Sigma_g^n$ along smooth centers, it is itself a smooth projective variety.

\item The complement $D = \overline{C}_n(\Sigma_g) \setminus C_n(\Sigma_g)$ is a \emph{normal crossing divisor}, i.e., locally analytically isomorphic to coordinate hyperplanes.

\item The boundary admits a natural stratification by nested sets:
\[
\partial\overline{C}_n(\Sigma_g)
=
\bigcup_{\mathcal N} D_{\mathcal N},
\qquad
D_{\mathcal N}:=\bigcap_{S\in\mathcal N}D_S ,
\]
where $\mathcal N$ ranges over collections of subsets
$S\subset\{1,\ldots,n\}$ with $|S|\geq2$ such that any two members are
nested or disjoint. The codimension of $D_{\mathcal N}$ is
$|\mathcal N|$. Set partitions record only the coarse collision pattern;
they do not remember the hierarchy of screens.

\item Each codimension-$1$ boundary divisor $D_S$ (for
$S \subset \{1,\ldots,n\}$, $|S| \geq 2$) has a dense open stratum
fibred over the configuration in which the labels in $S$ have been
replaced by their collision centre:
\[
D_S^\circ
\longrightarrow
\overline{C}_{n-|S|+1}(\Sigma_g),
\qquad
\text{fibre }
\overline{C}_{|S|}(\mathbb{A}^1)\bigm/
\mathrm{Aff}(\mathbb{A}^1),
\]
where the base records the collision center plus the remaining
$n-|S|$ points on $\Sigma_g$, and the fibre records relative positions
on the tangent-space screen (genus~$0$), \emph{modulo the affine group}
$\mathrm{Aff}(\mathbb{A}^1) = \mathbb{C} \rtimes \mathbb{C}^*$
of translations and dilations on the screen. The translation quotient
eliminates the choice of origin in $T_q\Sigma_g$, and the dilation
quotient eliminates the overall scale. For $|S|=2$ on a curve this
screen is a point; projective shape data first appears for
$|S|\geq3$. For deeper strata corresponding to nested collisions, this
description iterates, giving products of affine-quotiented genus-$0$
screens over the surviving configuration. Unlike the
Deligne--Mumford compactification $\overline{\mathcal{M}}_{g,n}$, the
FM compactification of a fixed curve involves no curve degeneration.

\item The construction is \emph{functorial} for maps that preserve
distinctness of labelled points, in particular for open embeddings and
automorphisms of curves:
$\overline{C}_n(\Sigma_g) \to \overline{C}_n(\Sigma_g')$ is then
compatible with the boundary stratification. A general smooth map need
not send configurations to configurations.
\end{enumerate}
\end{theorem}

\begin{proof}[Construction]
The construction proceeds in stages:

\medskip
\noindent\emph{Stage 0: Initial Space}

Begin with the smooth space $\Sigma_g^n$. The configuration space is the complement of the ``fat diagonal'':
\[C_n(\Sigma_g) = \Sigma_g^n \setminus \bigcup_{1 \leq i < j \leq n} \Delta_{ij}\]
where $\Delta_{ij} = \{(x_1,\ldots,x_n) \in \Sigma_g^n : x_i = x_j\}$ is a smooth divisor of codimension 1.

\medskip
\noindent\emph{Stage 1: Blow Up Diagonal}

First blow up the full diagonal $\Delta_n = \{x_1 = \cdots = x_n\}$ (codimension $n-1$):
\[\widetilde{\Sigma_g^n}_1 = \text{Bl}_{\Delta_n}(\Sigma_g^n)\]

\emph{Local coordinates near $\Delta_n$.} Choose a point $p \in \Delta_n$ and local coordinate $z$ on $\Sigma_g$ near $p$. Near $p$, the coordinates on $\Sigma_g^n$ are $(z_1,\ldots,z_n)$. The blow-up introduces:
\begin{itemize}
\item \emph{Center of mass:} $u = \frac{1}{n}\sum_{i=1}^n z_i$
\item \emph{Relative coordinates:} $\zeta_i = z_i - u$ for $i = 1,\ldots,n-1$ (with $\zeta_n = -\sum_{i=1}^{n-1}\zeta_i$)
\item \emph{Projective directions:} $[\zeta_1 : \cdots : \zeta_{n-1}] \in \mathbb{P}^{n-2}$
\end{itemize}

The exceptional divisor $E_n$ is a $\mathbb{P}^{n-2}$-bundle over $\Sigma_g$ (the projectivization of the normal bundle $N_{\Delta_n/\Sigma_g^n} \cong \bigoplus^{n-1} T\Sigma_g$), parametrizing the collision location (the $\Sigma_g$ factor) and the relative directions of approach (the $\mathbb{P}^{n-2}$ fiber).

\medskip
\noindent\emph{Stage 2: Blow Up Partial Diagonals}

Next, blow up the proper transform of each partial diagonal $\Delta_S$ for $S \subsetneq \{1,\ldots,n\}$ with $|S| \geq 2$, proceeding in \emph{decreasing order of $|S|$} (equivalently, decreasing order of codimension, since $\operatorname{codim}(\Delta_S) = |S|-1$).

For a subset $S = \{i_1, \ldots, i_k\}$ with $2 \leq k < n$:
\[\widetilde{\Sigma_g^n}_{S} = \text{Bl}_{\widetilde{\Delta_S}}(\widetilde{\Sigma_g^n}_{S_{\text{prev}}})\]
where $\widetilde{\Delta_S}$ is the proper transform of $\Delta_S$ from the previous blow-up stage.

The blow-ups are taken in decreasing codimension, which ensures:
\begin{enumerate}
\item All centers of blow-up are smooth
\item The final result is independent of choices within each codimension
\item Normal crossings are preserved at each stage
\end{enumerate}

\medskip
\noindent\emph{Stage 3: Final Compactification}

After all blow-ups:
\[\overline{C}_n(\Sigma_g) = \widetilde{\Sigma_g^n}_{\text{final}}\]

The boundary divisors $D_S$ are indexed by subsets $S$ with
$|S|\geq2$. For $|S|\geq3$ they are exceptional divisors of the
corresponding blow-ups; for $|S|=2$ on a curve the partial diagonal is
already Cartier, so the pair divisor is its proper transform with
logarithmic boundary structure rather than a new projective exceptional
fiber.

\medskip
\noindent\emph{Normal crossings.}

The divisor $D = \bigcup_S D_S$ has normal crossings by a local
coordinate calculation. Near a point in
$D_{S_1} \cap \cdots \cap D_{S_m}$, where
$S_1, \ldots, S_m$ form a nest, there are local analytic coordinates:
\[
(u,\;x_{S_1},\ldots,x_{S_m},\;q_{S_1},\ldots,q_{S_m},\;w_1,\ldots,w_k)
\]
where:
\begin{itemize}
\item $u \in \Sigma_g$ is the common collision point
\item $x_{S_j}$ is a holomorphic normal coordinate for the $j$-th
collision divisor
\item $q_{S_j}$ is the corresponding compactified screen coordinate
\item $w_1, \ldots, w_k$ parametrize points not involved in collisions
\end{itemize}

The divisors are locally:
\[D_{S_j} = \{x_{S_j} = 0\}\]
These are precisely coordinate hyperplanes, hence normal crossing.
Disjoint collision clusters contribute independent factors of the same
form.

\medskip
\noindent\emph{Functoriality.}

If $f: \Sigma_g \to \Sigma_{g'}$ is an open embedding, it induces
$f^{(n)}: \Sigma_g^n \to \Sigma_{g'}^n$ by
$(x_1,\ldots,x_n) \mapsto (f(x_1),\ldots,f(x_n))$. The map
$f^{(n)}$ preserves the configuration locus and carries diagonals to
diagonals:
\[f^{(n)}(\Delta_S) \subseteq \Delta_S\]
so it lifts canonically to the blow-ups, giving:
\[\overline{f^{(n)}}: \overline{C}_n(\Sigma_g) \to \overline{C}_n(\Sigma_{g'})\]
compatible with boundary stratification.
\end{proof}

\begin{remark}[Recording collision data]\label{rem:recording-collisions}
The complex FM blow-up records holomorphic normal variables $x_S$ and
compactified screens. A real oriented blow-up may further decompose
$x_S=r_Se^{i\vartheta_S}$, but the algebraic collision residue uses
$d\log x_S$. These are the data appearing in the OPE
$\phi_i(z)\phi_j(w)=\sum_m(\phi_i{}_{(m)}\phi_j)(w)/(z-w)^{m+1}
+\cdots$ and in its nested associativity constraints.
\end{remark}

\subsubsection{Boundary stratification and stable curves}

For a fixed smooth curve $\Sigma_g$, the FM boundary consists only of
point-collision strata. The additional stable-curve boundary appears in
the relative Deligne--Mumford compactification, where the curve itself
is allowed to degenerate.

\begin{theorem}[Boundary strata of \texorpdfstring{$\overline{\mathcal{M}}_{g,n}$}{M-bar(g,n)} {\cite{DeligneM69,Knudsen83}}; \ClaimStatusProvedElsewhere]\label{thm:boundary-higher-genus}
\index{Deligne--Mumford compactification}
\index{boundary stratum}
For the Deligne--Mumford--Knudsen moduli space
$\overline{\mathcal{M}}_{g,n}$, the boundary
$\partial\overline{\mathcal{M}}_{g,n}$ consists of:

\begin{enumerate}
\item \emph{Rational-tail collision strata:} stable curves in which
marked points indexed by $S$ lie on a rational bubble attached to the
main component; this is the stable-curve replacement for literal marked
point collision
\item \emph{Degeneration strata:} $D_{\Gamma}$ where the curve degenerates to a stable nodal curve of genus $g$ with dual graph $\Gamma$
\end{enumerate}

The Fulton--MacPherson compactification $\overline{C}_n(\Sigma_g)$ of a \emph{fixed} smooth curve $\Sigma_g$ has only collision strata; no curve degeneration occurs. The degeneration strata appear when one works with the universal family over $\overline{\mathcal{M}}_g$, i.e., when the underlying curve is allowed to vary and degenerate.
\end{theorem}

\begin{proof}
There is a natural forgetful map:
\[
\overline{\mathcal{M}}_{g,n} \xrightarrow{\;\pi\;} \overline{\mathcal{M}}_g
\]
that forgets the marked points and stabilizes. Over the interior $\mathcal{M}_g$, the fiber $\pi^{-1}([\Sigma_g])$ is the FM compactification $\overline{C}_n(\Sigma_g)$. The full boundary $\partial\overline{\mathcal{M}}_{g,n}$ decomposes according to two mechanisms.

\emph{(1) Collision strata.}
For a subset $S \subseteq \{1, \ldots, n\}$ with $|S| \geq 2$, the
relative FM construction applied fiberwise over the smooth locus blows
up the partial diagonal
$\Delta_S=\{z_i=z_j\text{ for all }i,j\in S\}$. On the open part of
the resulting divisor the local model is the collision center together
with the screen
$\overline C_{|S|}(\mathbb{A}^1)/\operatorname{Aff}(\mathbb{A}^1)$ and
the remaining marked points. In the stable-curve compactification this
same limiting data is represented by a rational tail carrying the
marked points in $S$, attached at the collision center. The resulting
strata have normal crossings (Theorem~\ref{thm:normal-crossings}).

\emph{(2) Degeneration strata.}
These arise from the boundary of $\overline{\mathcal{M}}_g$, not from point collisions. A boundary stratum of $\overline{\mathcal{M}}_{g,n}$ is indexed by a stable graph $\Gamma$ (Definition~\ref{def:stable-graph}) with $\sum_{v} g(v) + b_1(\Gamma) = g$, where $b_1(\Gamma) = |E(\Gamma)| - |V(\Gamma)| + 1$ is the first Betti number. The corresponding stratum is:
\[
D_\Gamma \cong \prod_{v \in V(\Gamma)} \overline{\mathcal{M}}_{g(v), n(v)} \bigg/ \operatorname{Aut}(\Gamma)
\]
where $n(v)$ counts the half-edges and tails at vertex $v$. Each edge $e \in E(\Gamma)$ corresponds to a node of the degenerate curve, with local model $xy = t_e$ for a smoothing parameter $t_e$. When the curve has a fixed complex structure with period matrix $\tau \in \mathbb{H}_g$, the degeneration strata $D_{\Gamma,\tau}$ record how $\tau$ approaches the boundary of the Satake compactification of $\mathbb{H}_g / \mathrm{Sp}_{2g}(\mathbb{Z})$.

\emph{Exhaustiveness.}
Every point of
$\partial\overline{\mathcal{M}}_{g,n}
= \overline{\mathcal{M}}_{g,n}\setminus \mathcal{M}_{g,n}$
belongs to at least one stratum: either some marked points have
collided on a rational bubble, or the underlying curve has degenerated,
or both. This follows from the properness of
$\overline{\mathcal{M}}_{g,n}$ (Deligne--Mumford--Knudsen
\cite{DeligneM69,Knudsen83}). For a fixed smooth curve
$\Sigma_g$, the FM boundary
$\partial\overline{C}_n(\Sigma_g)$ contains only the collision part;
the curve-degeneration part appears after passage to the relative
universal family.
\end{proof}

\begin{definition}[Stable graph]\label{def:stable-graph}
\index{stable graph|textbf}
A \emph{stable graph} $\Gamma$ of genus $g$ with $n$ marked points consists of:
\begin{itemize}
\item A connected graph with vertices $V(\Gamma)$ and edges $E(\Gamma)$
\item A genus function $g: V(\Gamma) \to \mathbb{Z}_{\geq 0}$
\item $n$ marked half-edges (tails) attached to vertices
\item \emph{Stability condition:} For each vertex $v$,
\[2g(v) - 2 + n(v) > 0\]
where $n(v) = \text{val}(v)$ is the valence (number of incident half-edges and tails)
\end{itemize}
with total genus:
\[g(\Gamma) = \sum_{v \in V(\Gamma)} g(v) + h^1(\Gamma) = g\]
where $h^1(\Gamma) = |E(\Gamma)| - |V(\Gamma)| + 1$ is the first Betti number.
\end{definition}

\begin{remark}[Stable-graph counts: numerical anchors]
\label{rem:stable-graph-counts-anchors}
\index{stable graph!counts at low genus}
The cardinalities of the boundary stratification index sets are
small but error-prone. The counts are:
\begin{itemize}
\item Genus~$1$, $n = 1$: 2 stable graphs (the smooth elliptic
curve; and the irreducible nodal cubic, a single genus-$0$ vertex
with a self-loop carrying the marked half-edge);
\item Genus~$1$, $n = 2$: 3 stable graphs (Example~\ref{ex:stable-graphs-g1-n2});
\item Genus~$2$, $n = 0$: \emph{7 stable graphs} (one smooth
of genus~$2$; one smooth of genus~$1$ with a self-loop; one
single-vertex with two self-loops; one with two genus-$1$
vertices joined by an edge; one with a genus-$1$ vertex joined by a
bridge to a genus-$0$ vertex carrying a self-loop; and two more on a
pair of genus-$0$ vertices, the theta graph with three connecting
edges and the barbell with one self-loop on each vertex and a
connecting bridge).
\emph{Not~$6$.}
\end{itemize}
The genus-$2$ count is the load-bearing combinatorial input
into the $C_{W_3}$ second-lane discovery and the iterated
Sugawara depth count of the Volume~II climax (3D quantum
gravity).
\end{remark}

\begin{example}[Stable graphs at genus 1, \texorpdfstring{$n=2$}{n=2}]\label{ex:stable-graphs-g1-n2}
For $\overline{\mathcal{M}}_{1,2}$ (genus 1, two marked points), the stable graphs indexing boundary strata are:

\begin{enumerate}
\item \emph{Interior:} Both points distinct on a smooth genus 1 curve
\[\Gamma_0: \text{one vertex with } g(v) = 1, n(v) = 2\]

\item \emph{Rational tail (bubble):} A genus-$0$ component carrying both marked points separates off from a genus-$1$ component
\[\Gamma_1: \text{two vertices } v_1, v_2 \text{ with } g(v_1) = 1,\; g(v_2) = 0,\; \text{one edge, both tails on } v_2\]
Stability: $2(1)-2+1 = 1 > 0$ for $v_1$ and $2(0)-2+3 = 1 > 0$ for $v_2$. This boundary divisor is isomorphic to $\overline{\mathcal{M}}_{1,1}$ (the attaching point on the genus-$1$ component).

\item \emph{Node formation:} The torus degenerates to a nodal curve (pinched cycle)
\[\Gamma_3: \text{one vertex with } g(v) = 0, \text{ one self-loop, two tails}\]
The self-loop contributes $h^1(\Gamma_3) = 1$, giving total genus $g(v) + h^1 = 0 + 1 = 1$.
This gives a divisor parametrizing nodal genus 1 curves with 2 marked points.
\end{enumerate}
\end{example}

\begin{remark}[Connection to moduli of stable curves]\label{rem:moduli-stable-curves}
For a fixed smooth curve $\Sigma_g$, stabilization gives a map $\pi: \overline{C}_n(\Sigma_g) \to \overline{\mathcal{M}}_{g,n}$ sending $(p_1, \ldots, p_n) \mapsto [\Sigma_g; p_1, \ldots, p_n]^{\mathrm{st}}$. Its image lies in the closure of the locus of curves isomorphic to $\Sigma_g$, and over boundary strata the bubble tree structure maps to the nodal curve structure. This underlies the modular properties of chiral correlators and the identification of conformal anomalies with failure of sections to extend over the boundary.
\end{remark}

\subsubsection{Local coordinates and blow-up charts}

\begin{theorem}[Local holomorphic coordinates near a collision divisor; \ClaimStatusProvedHere]\label{thm:local-coords-boundary}
Let $D_S \subset \partial\overline{C}_n(\Sigma_g)$ be the
Fulton--MacPherson divisor corresponding to collision of
$S=\{i_1,\ldots,i_k\}$, $k\geq 2$, at a general point of the divisor.
After choosing a holomorphic coordinate $\zeta$ on $\Sigma_g$ at the
collision point, there are analytic coordinates
\[
(p,x_S,u_1,\ldots,u_{k-1},w_\alpha)_{\alpha\notin S}
\]
near that point such that:
\begin{enumerate}
\item $p$ is the collision point and $w_\alpha$ are the remaining
 marked points;
\item $D_S$ is the coordinate divisor $\{x_S=0\}$;
\item for $i,j\in S$ one has
 \[
 \zeta(z_i)-\zeta(z_j)=x_S\,u_{ij},
 \]
 where $u_{ij}$ is regular and is nonzero on the open chart of the
 collision screen in which the two labelled directions remain distinct;
\item hence
 \[
 \eta_{ij}=d\log(\zeta(z_i)-\zeta(z_j))
 =d\log x_S+d\log u_{ij}
 \]
 on this chart, so $\operatorname{Res}_{D_S}\eta_{ij}=1$ precisely
 for the pairs separated by the screen coordinate.
\end{enumerate}
At a normal-crossing corner indexed by a nested or disjoint collection
of collision subsets, the corresponding scale parameters
$x_{S_1},\ldots,x_{S_m}$ are simultaneous normal coordinates.
\end{theorem}

\begin{proof}
This is the standard local chart of the iterated blow-up of the small
diagonal. In a holomorphic coordinate, write the relative coordinates
of the colliding points as
\[
\xi_j=\zeta(z_{i_j})-\frac1k\sum_{\ell=1}^k\zeta(z_{i_\ell}),
\qquad \sum_j\xi_j=0.
\]
On an affine chart of the exceptional projective space choose one
nonzero relative coordinate as the scale $x_S$ and write all other
relative coordinates as $x_S u_j$. Then $D_S=\{x_S=0\}$ and pairwise
differences inside the cluster have the asserted form. The
Fulton--MacPherson building set permits exactly nested or disjoint
collections of collision subsets to occur simultaneously; the
corresponding blow-up scales give the stated normal coordinates.
\end{proof}

\begin{remark}[Ordering convention for collisions]\label{rem:collision-ordering}
When a collision subset is written as $\{i_1,\ldots,i_k\}$, the
indices so that $i_1<\cdots<i_k$. This is a notation convention for
labels and signs; it is not an additional geometric ordering of the
Fulton--MacPherson compactification.
\end{remark}

\begin{example}[Explicit coordinates for three points]\label{ex:three-points-coords}
For $n=3$ on $\Sigma_g$, consider the divisor $D_{12}$ where $z_1 \to z_2$.

\emph{Coordinates.}
\begin{itemize}
\item $p \in \Sigma_g$: collision point
\item $x_{12}$: holomorphic normal coordinate $z_1-z_2$ in a chosen
local chart
\item $w = z_3$: third point
\end{itemize}

\emph{Reconstruction.}
\[
z_1=p+\frac{x_{12}}2,\qquad
z_2=p-\frac{x_{12}}2,\qquad z_3=w.
\]

\emph{Divisor.}
\[
D_{12} = \{x_{12}=0\},
\]
parametrized by the collision point $p$ on the first factor and the
spectator $w = z_3$ on the second. For a curve, the normal bundle of
the diagonal $\Delta \hookrightarrow \Sigma_g \times \Sigma_g$ has
complex rank~$1$, so the projectivized fibre $\mathbb{P}^0$ is a point
and no angular coordinate contributes additional moduli. A real
oriented blow-up may write $x_{12}=re^{i\vartheta}$, but the complex
FM chart and the logarithmic residue use $x_{12}$.
\end{example}

\subsubsection{Normal crossing property and residues}

\begin{theorem}[Normal crossings; \ClaimStatusProvedHere]\label{thm:normal-crossings}
\index{normal crossings|textbf}
\index{boundary divisor|textbf}
The boundary divisor $D = \partial\overline{C}_n(\Sigma_g)$ is a \emph{strict normal crossing divisor}.

More precisely, if $D = \bigcup_{\alpha} D_\alpha$ is the decomposition into irreducible components, then:
\begin{enumerate}
\item Each $D_\alpha$ is smooth
\item At any point $x \in D_{\alpha_1} \cap \cdots \cap D_{\alpha_k}$ (intersection of $k$ components), there exist local analytic coordinates $(u_1, \ldots, u_N)$ near $x$ such that:
\[D_{\alpha_j} = \{u_j = 0\} \text{ for } j = 1,\ldots,k\]
\item The components intersect transversely: $\bigcap_{j=1}^k T_x D_{\alpha_j}$ has codimension $k$ in $T_x \overline{C}_n(\Sigma_g)$, equivalently the conormals $du_1, \ldots, du_k$ are linearly independent at $x$
\end{enumerate}
\end{theorem}

\begin{proof}
Normal crossings follow from the blow-up construction.

\medskip
\noindent\emph{Single Divisor ($k=1$).}

Each divisor $D_\alpha=D_S$ is smooth in the FM wonderful
compactification. For $|S|\geq3$ it is an exceptional divisor of a
blow-up of a smooth center; for $|S|=2$ on a curve it is the proper
transform of a smooth Cartier diagonal. In both cases it is smooth.

\medskip
\noindent\emph{Multiple Intersections ($k \geq 2$).}

Suppose $x \in D_{S_1} \cap \cdots \cap D_{S_k}$ where $S_1, \ldots, S_k$ are distinct subsets.

For the divisors to intersect at $x$, the sets must form a nest: any
two of them are either disjoint or one contains the other. Overlapping
non-nested subsets, such as $\{1,2\}$ and $\{2,3\}$, are separated by
the blow-up of the larger collision locus.

\emph{Local coordinates for nested sets.}

Assume $S_1 \subsetneq S_2 \subsetneq \cdots \subsetneq S_k$. Near
$x$, there are coordinates:
\[
(p,\;x_{S_1},\ldots,x_{S_k},\;q_{S_1},\ldots,q_{S_k},\ldots)
\]
where:
\begin{itemize}
\item $x_{S_j}$ is the holomorphic normal coordinate for $D_{S_j}$
\item $q_{S_j}$ is the compactified screen coordinate at level $j$
\item $p \in \Sigma_g$ is the ultimate collision point
\end{itemize}

The divisors are:
\[D_{S_j} = \{x_{S_j} = 0\}\]
These are coordinate hyperplanes, hence normal crossing.
If several collision subsets are disjoint, the local chart is the
product of the corresponding chain charts, with one independent
normal coordinate for each divisor.

\medskip
\noindent\emph{Transversality.}

The tangent spaces satisfy:
\[T_x D_{S_j} = \{\frac{\partial}{\partial x_{S_j}} = 0\} \subset T_x \overline{C}_n(\Sigma_g)\]
Since the $x_{S_j}$ are independent coordinates:
\[\dim(T_x D_{S_1} \cap \cdots \cap T_x D_{S_k}) = \dim(T_x \overline{C}_n(\Sigma_g)) - k\]
which is the expected codimension for the intersection, confirming transversality.
\end{proof}

\subsection{Stratification}

\subsubsection{Incidence relations and poset structure}

\begin{definition}[Stratification poset]\label{def:stratification-poset}
The FM boundary strata are indexed by non-empty nested sets
\[
\mathcal N=\{S_1,\ldots,S_r\},\qquad |S_a|\geq2,
\]
where any two subsets are disjoint or comparable by inclusion. The
open stratum attached to $\mathcal N$ is
\[
D_{\mathcal N}^{\circ}
=
\left(\bigcap_{S\in\mathcal N}D_S\right)
\setminus
\bigcup_{T\notin\mathcal N}D_T .
\]
The partial order is inclusion of nested sets:
\[
\mathcal N\leq \mathcal N'
\quad\Longleftrightarrow\quad
\mathcal N\subseteq\mathcal N'.
\]
Set partitions give only a coarse collision pattern; they do not
remember the compactified screen strata and therefore are not the FM
stratification poset.
\end{definition}

\begin{example}[Poset for \texorpdfstring{$n=3$}{n=3}]\label{ex:poset-n3}
For $n=3$, the boundary divisors are
\[
D_{12},\quad D_{13},\quad D_{23},\quad D_{123}.
\]
The non-empty nested sets are the four singletons
$\{12\}$, $\{13\}$, $\{23\}$, $\{123\}$ and the three nested pairs
\[
\{123,12\},\qquad \{123,13\},\qquad \{123,23\}.
\]
Thus $D_{123}$ is itself a divisor. The closure of $D_{12}^{\circ}$
contains the corner $D_{123}\cap D_{12}$, not the whole open divisor
$D_{123}^{\circ}$.
\end{example}

\begin{theorem}[Closure relations; \ClaimStatusProvedHere]\label{thm:closure-relations}
For a nested set $\mathcal N$, the closure of its open stratum is
\[
\overline{D_{\mathcal N}^{\circ}}
=
\bigcup_{\mathcal N'\supseteq \mathcal N}
D_{\mathcal N'}^{\circ},
\]
where $\mathcal N'$ ranges over nested sets containing $\mathcal N$.
In particular:
\begin{enumerate}
\item $\operatorname{codim}D_{\mathcal N}^{\circ}=|\mathcal N|$;
\item $\overline{D_{\mathcal N}^{\circ}}\cap
\overline{D_{\mathcal M}^{\circ}}$ is nonempty only when
$\mathcal N\cup\mathcal M$ is nested; in that case it contains the
stratum indexed by $\mathcal N\cup\mathcal M$;
\item overlapping non-nested subsets have no common boundary stratum.
\end{enumerate}
\end{theorem}

\begin{proof}
In a chart attached to $\mathcal N$, the divisor $D_S$ for
$S\in\mathcal N$ has equation $x_S=0$. Adding another compatible
divisor $D_T$ means imposing one more independent equation $x_T=0$,
which gives the stratum for $\mathcal N\cup\{T\}$. Iterating gives the
displayed closure formula and the codimension count. If two subsets
overlap without containment, the FM blow-up separates the corresponding
proper transforms; equivalently, there is no chart in which both
boundary equations occur.
\end{proof}

\begin{corollary}[Boundary divisors in the FM compactification \cite{FM94}; \ClaimStatusProvedElsewhere]\label{cor:dimension-strata}
For each subset $S \subset \{1,\ldots,n\}$ with $|S| \geq 2$, the boundary divisor
$D_S \subset \overline{C}_n(\Sigma_g)$ has \emph{codimension~1}. The divisor $D_S$
parametrizes configurations where the points indexed by $S$ collide, together with
the ``screen'' recording how they approach each other.

In particular:
\begin{itemize}
\item Each $D_S$ is a divisor ($\operatorname{codim} = 1$, $\dim = n-1$)
\item The open stratum of $D_S$ fibers over $C_{n-|S|+1}(\Sigma_g)$ (the collision
center replaces the cluster) with fiber
$\overline{C}_{|S|}(\mathbb{C})\big/\mathrm{Aff}(\mathbb{A}^1)$,
the FM screen modulo the affine group of translations and dilations
on the tangent space. Equivalently, the screen is the projectivized
shape space of $|S|$ points on the tangent line $T_q\Sigma_g$ at the
collision point $q$, with the origin and overall scale removed.
\item Deeper strata (codimension $\geq 2$) arise from intersections
$D_{S_1}\cap D_{S_2}$ when $S_1$ and $S_2$ are nested or disjoint.
\end{itemize}
\end{corollary}

\begin{theorem}[Boundary stratification \cite{FM94}; \ClaimStatusProvedElsewhere]
\label{thm:boundary-stratification}
The boundary has a natural stratification:
\[
\partial\overline{C}_n(X)
=
\bigcup_{\varnothing\neq\mathcal N} D_{\mathcal N}^{\circ},
\]
where $\mathcal N$ runs over nested sets of subsets of
$\{1,\ldots,n\}$ of cardinality at least two.
\end{theorem}

\subsection{Logarithmic differential forms}

\begin{definition}[Logarithmic forms]\label{def:log-forms}
\index{logarithmic forms|textbf}
The forms $\eta_{ij} = d\log(z_i - z_j)$ that served as propagators in the Heisenberg bar complex (\S\ref{sec:frame-bar-deg1}) are the simplest instances of the following class.

A differential $k$-form $\omega$ on $\overline{C}_n(\Sigma_g)$ has \emph{logarithmic poles along $D$} if:
\begin{enumerate}
\item $\omega$ is smooth on the interior $C_n(\Sigma_g)$
\item Near each divisor $D_\alpha$ defined locally by $\{f_\alpha = 0\}$, we have:
\[\omega = \frac{df_\alpha}{f_\alpha} \wedge \alpha + \beta\]
where $\alpha$ is a $(k-1)$-form and $\beta$ is a $k$-form, both smooth up to $D_\alpha$
\end{enumerate}

The sheaf of logarithmic $k$-forms is denoted:
\[\Omega^k_{\overline{C}_n(\Sigma_g)}(\log D)\]
\end{definition}

\begin{remark}[Multi-logarithmic extension]\label{rem:multi-log-codim2}
On codimension-$k$ strata where $k$ boundary divisors
$D_{S_1}, \ldots, D_{S_k}$ intersect, logarithmic forms
restrict to \emph{multi-logarithmic} forms with logarithmic
poles along each component simultaneously.
Since $D$ is a normal crossing divisor
(Theorem~\ref{thm:FM}(3)), the sheaf
$\Omega^k(\log D)$ is locally free, so multi-logarithmic
extension across higher-codimension strata follows from the
normal-crossing coordinate calculation.
\end{remark}

\begin{remark}[Logarithmic condition]\label{rem:why-log}
The logarithmic condition is precisely what is needed for well-defined residues.

A general form with poles along $D$ might have:
\[\omega = \frac{\alpha}{f^k}\]
for $k \geq 2$ (higher-order pole). Such forms do not have well-defined residues.

Logarithmic forms have:
\[\omega = \frac{df}{f} \wedge \alpha + \beta\]
which has a \emph{simple pole} with residue $\alpha|_{f=0}$.

For chiral algebras: the OPE has the form
\[\phi_i(z)\phi_j(w) \sim \frac{C_{ij}^k}{(z-w)^\Delta} \phi_k(w)\]
Combined with $\eta_{ij} = \frac{dz-dw}{z-w}$, we get:
\[\frac{1}{(z-w)^\Delta} \cdot \frac{dz-dw}{z-w} = \frac{d(z-w)}{(z-w)^{\Delta+1}}\]
For $\Delta = 0$ (no pole in OPE): this is $\frac{d(z-w)}{z-w}$ = logarithmic.

Logarithmic forms therefore furnish the correct coefficient system for chiral algebras.
\end{remark}

\begin{example}[Logarithmic form for two points]\label{ex:log-form-two-points}
The basic logarithmic 1-form for configuration of two points:
\[\eta_{12} = d\log(z_1 - z_2) = \frac{dz_1 - dz_2}{z_1 - z_2}\]

\emph{Analysis.}
\begin{itemize}
\item On $C_2(\Sigma_g)$ (where $z_1 \neq z_2$): $\eta_{12}$ is smooth
\item Near $D_{12}$ (where $z_1 \to z_2$): Using $\epsilon = z_1 - z_2$, we have:
\[\eta_{12} = \frac{d\epsilon}{\epsilon} + \text{(smooth terms)}\]
This is a logarithmic pole.
\item The residue:
\[\text{Res}_{D_{12}}(\eta_{12}) = 1 \in \Omega^0_{D_{12}} = \mathcal{O}_{D_{12}}\]
\end{itemize}
\end{example}

\begin{theorem}[Logarithmic complex; \ClaimStatusProvedHere]\label{thm:log-complex}
\index{de Rham functor}
The sheaf of logarithmic differential forms
$\Omega^\bullet_{\overline{C}_n(\Sigma_g)}(\log D)$
forms a complex under the de Rham differential:
\[d: \Omega^k(\log D) \to \Omega^{k+1}(\log D)\]

Three properties hold:
\begin{enumerate}
\item $d$ preserves logarithmic poles: if $\omega$ has log poles along $D$, then $d\omega$ also has log poles
\item $d^2 = 0$ (as always for de Rham differential)
\item By the logarithmic Poincar\'e lemma, $H^*(\Omega^\bullet(\log D))$ computes the cohomology of the \emph{open} configuration space $C_n(\Sigma_g) = \overline{C}_n(\Sigma_g) \setminus D$ with coefficients in $\mathbb{C}$
\end{enumerate}
\end{theorem}

\begin{proof}
\emph{Part 1: Preservation of log poles.}

Locally, if $\omega = \frac{df}{f} \wedge \alpha + \beta$ with $\alpha, \beta$ smooth, then:
\[d\omega = d\left(\frac{df}{f}\right) \wedge \alpha + \frac{df}{f} \wedge d\alpha + d\beta\]

Compute:
\[d\left(\frac{df}{f}\right) = -\frac{df \wedge df}{f^2} = 0\]
(since $df \wedge df = 0$)

Therefore:
\[d\omega = \frac{df}{f} \wedge d\alpha + d\beta\]

Since $d\alpha$ and $d\beta$ are smooth, this is again a logarithmic form.

\emph{Part 2: $d^2 = 0$.}
This is the fundamental property of the de Rham differential, independent of logarithmic conditions.

\emph{Part 3: Cohomology.}
By the logarithmic Poincar\'{e} lemma (Deligne~\cite{Deligne71}), the logarithmic de~Rham complex $\Omega^\bullet(\log D)$ on $\overline{C}_n(\Sigma_g)$ computes the cohomology of the \emph{open} complement $\overline{C}_n(\Sigma_g) \setminus D = C_n(\Sigma_g)$:
\[H^*(\overline{C}_n(\Sigma_g),\, \Omega^\bullet(\log D)) \cong H^*(C_n(\Sigma_g);\, \mathbb{C})\]
\end{proof}

\begin{theorem}[Arnold relations and KZ flatness; \ClaimStatusProvedHere]\label{thm:arnold-relations}
\index{Arnold relations|textbf}
\index{propagator!logarithmic|textbf}
\index{Arnold connection!flatness|textbf}
In the Heisenberg bar complex (\S\ref{sec:frame-bar-deg2}), the identity $d^2 = 0$ at degree~$2$ reduces to the three-term relation $\eta_{12} \wedge \eta_{23} + \eta_{23} \wedge \eta_{31} + \eta_{31} \wedge \eta_{12} = 0$. This is the Arnold relation. Geometrically the relation is the scalar differential-form half of the flatness of the universal Arnold connection
$\nabla_{\mathrm{Arnold}}$ on $\operatorname{Conf}_n(\mathbb{C})$
whose connection one-form is $\eta = \sum_{i<j} \eta_{ij}\,T_{ij}$
(the universal KZ form, with $T_{ij}$ infinitesimal-braid generators);
the Lie-algebra half is the infinitesimal braid relation
$[T_{ij},T_{kl}]=0$ for disjoint pairs and
$[T_{ij},T_{ik}+T_{jk}]=0$ for triples;
the chiral specialisation of the affine simple-pole component,
\[
d_{\mathrm{bar}}^{(0),\mathrm{simp}}
=
\mathrm{KZ}^{*}(\nabla_{\mathrm{Arnold}})_{\mathrm{simp}},
\]
is the genus-$0$ Arnold-local corollary proved in
Chapter~\ref{chap:climax-platonic}. The full bar differential also
contains the higher OPE modes and global curve corrections.

The logarithmic 1-forms $\eta_{ij} = d\log(z_i - z_j)$ satisfy fundamental relations:

\begin{enumerate}
\item \emph{Symmetry:} $\eta_{ji} = \eta_{ij}$ (as differential forms, since $d\log(z_j - z_i) = d\log(-(z_i-z_j)) = d\log(z_i-z_j)$, using $d\log(-1)=0$). Antisymmetry $e_{ji} = -e_{ij}$ is imposed at the level of the Orlik--Solomon algebra generators, not the forms themselves.

\item \emph{Arnold relation:} For distinct $i,j,k$:
\[\eta_{ij} \wedge \eta_{jk} + \eta_{jk} \wedge \eta_{ki} + \eta_{ki} \wedge \eta_{ij} = 0\]

\item \emph{Completeness (genus split):} For $\Sigma_0 = \mathbb{C}$, the
classes $\eta_{ij}$ generate $H^*(C_n(\mathbb{C}); \mathbb{C})$ and the
Arnold relations generate all relations. For $X = \mathbb{P}^1$, the
degree-$1$ projective classes are represented after a
$\mathrm{PSL}_2$ gauge fixing and satisfy the projective linear
relations $\sum_{j \neq i}\bar\eta_{ij} = 0$; the full cohomology also
contains the degree-$3$ class of $\mathrm{PSL}_2(\mathbb{C})$. For
$g \geq 1$, the cohomology is computed by the
Totaro--K\v{r}\'{\i}\v{z}--Bezrukavnikov model with additional
generators pulled back from $H^*(\Sigma_g)$ and mixed relations
$\eta_{ij}(\pi_i^*\xi-\pi_j^*\xi)=0$.
\end{enumerate}
\end{theorem}

\begin{proof}
\emph{Part 1: Symmetry of forms.}
\[\eta_{ji} = d\log(z_j - z_i) = \frac{dz_j - dz_i}{z_j - z_i} = \frac{-(dz_i - dz_j)}{-(z_i - z_j)} = \frac{dz_i - dz_j}{z_i - z_j} = \eta_{ij}\]
(The two minus signs cancel. Equivalently, $d\log(z_j-z_i) = d\log(-(z_i-z_j)) = d\log(-1) + d\log(z_i-z_j) = \eta_{ij}$ since $d\log(-1) = 0$. Note: the antisymmetry $e_{ji} = -e_{ij}$ is an Orlik--Solomon \emph{algebra} convention, not a property of the forms.)

\emph{Part 2: Arnold relation.}
Fix any three distinct indices $i,j,k \in \{1,\ldots,n\}$. Write $z_{ab} = z_a - z_b$ and $dz_{ab} = dz_a - dz_b$ for the coordinate differences and their differentials. The fundamental identity is:
\begin{equation}\label{eq:cyclic-sum-zero}
z_{ij} + z_{jk} + z_{ki} = (z_i - z_j) + (z_j - z_k) + (z_k - z_i) = 0
\end{equation}
and identically $dz_{ij} + dz_{jk} + dz_{ki} = 0$ for the differentials.

Each term in the Arnold relation is:
\begin{align*}
\eta_{ij} \wedge \eta_{jk} &= \frac{dz_{ij} \wedge dz_{jk}}{z_{ij} \cdot z_{jk}}, \qquad
\eta_{jk} \wedge \eta_{ki} = \frac{dz_{jk} \wedge dz_{ki}}{z_{jk} \cdot z_{ki}}, \qquad
\eta_{ki} \wedge \eta_{ij} = \frac{dz_{ki} \wedge dz_{ij}}{z_{ki} \cdot z_{ij}}.
\end{align*}

\emph{Step 1: the three numerator $2$-forms coincide.} Using $dz_{ki} = -(dz_{ij} + dz_{jk})$ (from the differential of~\eqref{eq:cyclic-sum-zero}):
\begin{align*}
dz_{jk} \wedge dz_{ki} &= dz_{jk} \wedge (-(dz_{ij} + dz_{jk})) = -dz_{jk} \wedge dz_{ij} = dz_{ij} \wedge dz_{jk}, \\
dz_{ki} \wedge dz_{ij} &= (-(dz_{ij} + dz_{jk})) \wedge dz_{ij} = -dz_{jk} \wedge dz_{ij} = dz_{ij} \wedge dz_{jk},
\end{align*}
using $d\alpha \wedge d\alpha = 0$ and $d\alpha \wedge d\beta = -d\beta \wedge d\alpha$ throughout. Thus all three numerators equal $dz_{ij} \wedge dz_{jk}$.

\emph{Step 2: the coefficient sum vanishes by partial fractions.} Bringing the three terms to the common denominator $z_{ij} \cdot z_{jk} \cdot z_{ki}$:
\begin{align*}
\eta_{ij} \wedge \eta_{jk} + \eta_{jk} \wedge \eta_{ki} + \eta_{ki} \wedge \eta_{ij}
&= \frac{dz_{ij} \wedge dz_{jk}}{z_{ij} \cdot z_{jk} \cdot z_{ki}} \bigl(z_{ki} + z_{ij} + z_{jk}\bigr) \\
&= \frac{dz_{ij} \wedge dz_{jk}}{z_{ij} \cdot z_{jk} \cdot z_{ki}} \cdot 0 = 0
\end{align*}
by~\eqref{eq:cyclic-sum-zero}. This holds for any triple $i,j,k$ and any $n$.

\emph{Alternative direct proof via coordinate reduction.} Set $u = z_i - z_j$ and $v = z_j - z_k$, so that $z_k - z_i = -(u+v)$. Since $d\log$ is insensitive to overall sign ($d\log(-f) = d\log(f)$ for any nonvanishing~$f$), the three propagators become
\[
\eta_{ij} = \frac{du}{u}, \qquad \eta_{jk} = \frac{dv}{v}, \qquad \eta_{ki} = \frac{du + dv}{u + v}.
\]
(For the third: $\eta_{ki} = d\log(z_k - z_i) = d\log(-(u+v)) = d\log(u+v) = (du+dv)/(u+v)$.) Each wedge product is now immediate:
\begin{align*}
\eta_{ij} \wedge \eta_{jk} &= \frac{du \wedge dv}{u \cdot v}, \\[4pt]
\eta_{jk} \wedge \eta_{ki} &= \frac{dv \wedge (du + dv)}{v(u+v)} = \frac{dv \wedge du}{v(u+v)} = -\frac{du \wedge dv}{v(u+v)}, \\[4pt]
\eta_{ki} \wedge \eta_{ij} &= \frac{(du+dv) \wedge du}{(u+v) \cdot u} = \frac{dv \wedge du}{(u+v) \cdot u} = -\frac{du \wedge dv}{(u+v) \cdot u},
\end{align*}
using $d\alpha \wedge d\alpha = 0$ throughout. Summing:
\[
\eta_{ij} \wedge \eta_{jk} + \eta_{jk} \wedge \eta_{ki} + \eta_{ki} \wedge \eta_{ij}
= du \wedge dv \left[\frac{1}{uv} - \frac{1}{v(u+v)} - \frac{1}{u(u+v)}\right].
\]
The bracketed expression reduces by partial fractions over the common denominator $uv(u+v)$:
\[
\frac{u+v}{uv(u+v)} - \frac{u}{uv(u+v)} - \frac{v}{uv(u+v)} = \frac{(u+v) - u - v}{uv(u+v)} = 0.
\]
This completes the proof for any triple of distinct indices and any~$n \geq 3$.

\emph{KZ flatness.}
For
$\omega_{\mathrm{Arnold}}=\sum_{i<j}T_{ij}\eta_{ij}$ one has
$d\eta_{ij}=0$, hence
\[
\nabla_{\mathrm{Arnold}}^2
=\frac12[\omega_{\mathrm{Arnold}},\omega_{\mathrm{Arnold}}].
\]
The summands with disjoint pairs vanish by
$[T_{ij},T_{kl}]=0$. For each triple $\{i,j,k\}$, the remaining
curvature contribution is
\[
[T_{ij},T_{ik}]\,\eta_{ij}\wedge\eta_{ik}
+[T_{ij},T_{jk}]\,\eta_{ij}\wedge\eta_{jk}
+[T_{ik},T_{jk}]\,\eta_{ik}\wedge\eta_{jk}.
\]
The infinitesimal braid relations identify these three Lie brackets up
to the signs forced by the displayed wedge order; the coefficient
therefore factors as one Lie bracket times the Arnold cyclic sum. The
Arnold identity kills that scalar sum. Thus
$\nabla_{\mathrm{Arnold}}^2=0$ on the affine screen.

\emph{Part 3: Completeness.}
The genus split is proved in
Proposition~\ref{prop:arnold-higher-genus}: the affine case follows by
the Fadell--Neuwirth fibration and a Hilbert-series comparison with the
Arnol'd--Orlik--Solomon algebra, the projective case adds the
$\mathrm{PSL}_2$ gauge relations, and the higher-genus case is computed
by the Totaro model \cite[Theorems~1 and~3]{Totaro96} together with its
rational-model refinements due to K\v{r}\'{\i}\v{z} and
Bezrukavnikov.
\end{proof}

\begin{proposition}[Higher-genus correction to the Arnold-only presentation; \ClaimStatusProvedHere]
\label{prop:arnold-higher-genus}
\index{configuration space!genus split}
\index{Totaro relations|textbf}
Let $X$ be a smooth complex curve.
\begin{enumerate}
\item \textup{\textbf{Affine line / formal disk.}}
For $X = \mathbb{A}^1$ \textup{(}equivalently, on a genus-$0$ tangent
screen or formal disk\textup{)}, the algebra
$H^*(C_n(\mathbb{A}^1);\mathbb{Q})$ is generated by the degree-$1$
classes $\eta_{ij} = d\log(z_i-z_j)$ subject only to
$\eta_{ji}=\eta_{ij}$, $\eta_{ij}^2=0$, and the Arnold relations
\[
\eta_{ij}\wedge\eta_{jk}
\;+\;
\eta_{jk}\wedge\eta_{ki}
\;+\;
\eta_{ki}\wedge\eta_{ij}
\;=\;0.
\]
Its Poincar\'e polynomial is
\[
P_{C_n(\mathbb{A}^1)}(t)
\;=\;
\prod_{k=1}^{n-1}(1+kt).
\]

\item \textup{\textbf{Projective line.}}
For $X = \mathbb{P}^1$ and $n \geq 3$, the unique M\"obius
transformation sending $(z_1,z_2,z_3)$ to $(0,1,\infty)$ yields
\[
C_n(\mathbb{P}^1)
\;\cong\;
\mathrm{PSL}_2(\mathbb{C})
\times
C_{n-3}\!\bigl(\mathbb{C}\setminus\{0,1\}\bigr).
\]
Hence
\[
H^*(C_n(\mathbb{P}^1);\mathbb{Q})
\;\cong\;
\Lambda(u)
\otimes
H^*\!\bigl(C_{n-3}(\mathbb{C}\setminus\{0,1\});\mathbb{Q}\bigr),
\qquad |u|=3.
\]
On the gauge-fixed factor, with moving points
$z_4,\ldots,z_n$, the degree-$1$ generators
\[
a_i := d\log z_i,\qquad
b_i := d\log(z_i-1),\qquad
c_{ij} := d\log(z_i-z_j)
\quad (4\leq i<j\leq n)
\]
satisfy the complete relations
\begin{align}
a_i\wedge b_i &= 0,\label{eq:p1-projective-rel-1}\\
a_i\wedge c_{ij} + c_{ij}\wedge a_j + a_j\wedge a_i &= 0,\label{eq:p1-projective-rel-2}\\
b_i\wedge c_{ij} + c_{ij}\wedge b_j + b_j\wedge b_i &= 0,\label{eq:p1-projective-rel-3}\\
c_{ij}\wedge c_{jk} + c_{jk}\wedge c_{ki} + c_{ki}\wedge c_{ij} &= 0.\label{eq:p1-projective-rel-4}
\end{align}
Equivalently, if $\bar\eta_{ij}$ denotes the degree-$1$ projective
class represented on the gauge-fixed factor, then its pairwise
subalgebra has the Arnold relations together with the projective linear
relations
\[
\sum_{j\neq i}\bar\eta_{ij}=0
\qquad (1\leq i\leq n).
\]
This presentation describes the degree-$1$ Orlik--Solomon part; the
tensor factor $\Lambda(u)$, $|u|=3$, is independent of these pairwise
classes.

\item \textup{\textbf{Positive genus.}}
Let $X=\Sigma_g$ with $g\geq 1$, let
$p_i\colon \Sigma_g^n\to\Sigma_g$ be the coordinate projections, choose
a symplectic basis
$\alpha_1,\beta_1,\ldots,\alpha_g,\beta_g\in H^1(\Sigma_g;\mathbb{Q})$,
and write $\omega\in H^2(\Sigma_g;\mathbb{Q})$ for the orientation
class. Then $H^*(C_n(\Sigma_g);\mathbb{Q})$ is computed by the
commutative differential graded algebra
\[
A_{g,n}
\;:=\;
\frac{H^*(\Sigma_g^n;\mathbb{Q})[\eta_{ij}]_{1\leq i<j\leq n}}
{\bigl(
\eta_{ij}^2,\;
\eta_{ij}\eta_{jk}+\eta_{jk}\eta_{ki}+\eta_{ki}\eta_{ij},\;
\eta_{ij}(p_i^*\xi-p_j^*\xi)\ \forall\,\xi\in H^*(\Sigma_g;\mathbb{Q})
\bigr)}
\]
with $|\eta_{ij}|=1$, $d|_{H^*(\Sigma_g^n)}=0$, and
\begin{equation}
\label{eq:totaro-curve-differential}
d\eta_{ij}
\;=\;
\Delta_{ij}
\;:=\;
\omega_i+\omega_j
\;+\;
\sum_{r=1}^g
\bigl(\alpha_{r,i}\beta_{r,j}-\beta_{r,i}\alpha_{r,j}\bigr).
\end{equation}
In particular,
\begin{align}
\eta_{ij}(\alpha_{r,i}-\alpha_{r,j}) &= 0,\label{eq:totaro-curve-rel-1}\\
\eta_{ij}(\beta_{r,i}-\beta_{r,j}) &= 0,\label{eq:totaro-curve-rel-2}\\
\eta_{ij}(\omega_i-\omega_j) &= 0.\label{eq:totaro-curve-rel-3}
\end{align}
Thus the Arnold relation is only the genus-$0$ part of the relation
package: for $g\geq 1$ one must also keep the surface generators and the
mixed period relations, and $\eta_{ij}$ are no longer closed classes in
cohomology.
\end{enumerate}
\end{proposition}

\begin{proof}
\emph{Part 1: $X=\mathbb{A}^1$.}
Let
$p_n\colon C_n(\mathbb{A}^1)\to C_{n-1}(\mathbb{A}^1)$
forget the last point. By Fadell--Neuwirth~\cite{FadellNeuwirth62},
this is a locally trivial fibration with fiber
\[
F_n
\;=\;
\mathbb{A}^1\setminus\{z_1,\ldots,z_{n-1}\}.
\]
The fiber cohomology is concentrated in degrees $0$ and $1$:
$H^0(F_n)=\mathbb{Q}$ and
$H^1(F_n)=\bigoplus_{i=1}^{n-1}\mathbb{Q}\,[\gamma_i]$, where
$\gamma_i$ is a small positively oriented loop around the puncture
$z_i$. The pure braid monodromy fixes each puncture individually, hence
acts trivially on $H^1(F_n)$. Therefore the Serre spectral sequence has
\[
E_2^{p,q}
\;=\;
H^p\!\bigl(C_{n-1}(\mathbb{A}^1);\mathbb{Q}\bigr)
\otimes
H^q(F_n;\mathbb{Q}),
\qquad q\in\{0,1\}.
\]
The class $\eta_{in}=d\log(z_i-z_n)$ restricts on the fiber to
$d\log(z_i-w)$, whose period on $\gamma_i$ is $2\pi i$ and whose period
on $\gamma_j$ for $j\neq i$ is $0$; thus
$[\eta_{1n}],\ldots,[\eta_{n-1,n}]$ form the fiber basis and kill the
only possible transgression. Hence
\[
P_n(t)
\;=\;
(1+(n-1)t)\,P_{n-1}(t),
\qquad
P_1(t)=1,
\]
so
$P_n(t)=\prod_{k=1}^{n-1}(1+kt)$.

Let $\mathrm{OS}_n$ be the graded algebra generated by symbols
$e_{ij}=e_{ji}$, modulo $e_{ij}^2=0$ and the Arnold relations.
Sending $e_{ij}\mapsto[\eta_{ij}]$ defines a surjective map
$\mathrm{OS}_n\twoheadrightarrow H^*(C_n(\mathbb{A}^1);\mathbb{Q})$.
Using the Arnold relation with the index $n$,
\[
e_{in}e_{jn}
\;=\;
e_{ij}(e_{jn}-e_{in}),
\]
every monomial in $\mathrm{OS}_n$ rewrites as a sum of monomials with at
most one factor involving the index $n$. Hence
$\mathrm{OS}_n$ is spanned by
$\mathrm{OS}_{n-1}$ and
$\mathrm{OS}_{n-1}e_{in}$ for $1\leq i\leq n-1$, so
\[
\operatorname{Hilb}(\mathrm{OS}_n;t)
\;\leq\;
(1+(n-1)t)\operatorname{Hilb}(\mathrm{OS}_{n-1};t).
\]
Starting from $\operatorname{Hilb}(\mathrm{OS}_1;t)=1$, this gives
\[
\operatorname{Hilb}(\mathrm{OS}_n;t)
\;\leq\;
\prod_{k=1}^{n-1}(1+kt)
\;=\;
P_n(t).
\]
Because $\mathrm{OS}_n\twoheadrightarrow H^*(C_n(\mathbb{A}^1);\mathbb{Q})$
is surjective and the target has the same Poincar\'e polynomial, the
map is an isomorphism. This is precisely the
Arnol'd--Brieskorn--Orlik--Solomon presentation.

\emph{Part 2: $X=\mathbb{P}^1$.}
For every ordered triple of distinct points on $\mathbb{P}^1$ there is a
unique M\"obius transformation carrying it to $(0,1,\infty)$, so
\[
(g,w_4,\ldots,w_n)
\longmapsto
\bigl(g(0),g(1),g(\infty),g(w_4),\ldots,g(w_n)\bigr)
\]
is an isomorphism
$\mathrm{PSL}_2(\mathbb{C})\times
C_{n-3}(\mathbb{C}\setminus\{0,1\})
\xrightarrow{\sim}
C_n(\mathbb{P}^1)$.
The group $\mathrm{PSL}_2(\mathbb{C})$ deformation retracts onto
$\mathrm{SO}(3)$, so
$H^*(\mathrm{PSL}_2(\mathbb{C});\mathbb{Q})\cong\Lambda(u)$ with
$|u|=3$, giving the stated tensor-product decomposition.

On the gauge-fixed factor, after relabelling the moving points by
$1,\ldots,m=n-3$, the forms $a_i$, $b_i$, and $c_{ij}$ are the pairwise
logarithmic differentials of the points
$0,1,w_1,\ldots,w_m$. The four displayed relations are exactly the
Arnold relations for the triples $(0,1,i)$, $(0,i,j)$, $(1,i,j)$, and
$(i,j,k)$. Forgetting the last moving point gives a fibration with fiber
$\mathbb{C}$ minus the marked set $\{0,1,z_4,\ldots,z_{n-1}\}$, so the
same argument as in Part~1 yields the recurrence
\[
P_m^{01}(t)
\;=\;
(1+(m+1)t)P_{m-1}^{01}(t)
\qquad
(m=n-3)
\]
for the Poincar\'e polynomial of
$C_m(\mathbb{C}\setminus\{0,1\})$.
Let $A_m^{01}$ be the algebra generated by
$a_1,\ldots,a_m$, $b_1,\ldots,b_m$, and $c_{ij}$ modulo the displayed
relations. The same reduction argument as in Part~1, now using
\eqref{eq:p1-projective-rel-1}--\eqref{eq:p1-projective-rel-4}, rewrites
every monomial as a sum of monomials with at most one factor involving
the last point: indeed
\[
a_m b_m = 0,\qquad
a_m c_{im} = a_i(a_m-c_{im}),\qquad
b_m c_{im} = b_i(b_m-c_{im}),\qquad
c_{im}c_{jm} = c_{ij}(c_{jm}-c_{im}),
\]
so every product with two ``$m$-factors'' reduces to a sum with only one
such factor. Hence
$A_m^{01}$ is spanned by
$A_{m-1}^{01}$ together with
$A_{m-1}^{01}a_m$, $A_{m-1}^{01}b_m$, and
$A_{m-1}^{01}c_{im}$ $(1\leq i<m)$, so
\[
\operatorname{Hilb}(A_m^{01};t)
\;\leq\;
(1+(m+1)t)\operatorname{Hilb}(A_{m-1}^{01};t).
\]
Because the cohomology of the gauge-fixed factor satisfies the same
recurrence and is generated by those degree-$1$ classes, the natural map
$A_m^{01}\twoheadrightarrow
H^*(C_m(\mathbb{C}\setminus\{0,1\});\mathbb{Q})$
is an isomorphism. Thus the projective degree-$1$ formulation
$\sum_{j\neq i}\bar\eta_{ij}=0$ is obtained by undoing the gauge fixing
on the Orlik--Solomon factor. The exterior generator
$u\in H^3(\mathrm{PSL}_2(\mathbb{C});\mathbb{Q})$ is not represented by
a pairwise logarithmic form. So $\mathbb{P}^1$ already carries more
than the affine Arnold relation:
projective rigidity contributes a linear relation at each marked point.

\emph{Part 3: $X=\Sigma_g$ with $g\geq 1$.}
Again forget the last point:
\[
p_n\colon C_n(\Sigma_g)\to C_{n-1}(\Sigma_g),
\qquad
F_n=\Sigma_g\setminus\{z_1,\ldots,z_{n-1}\}.
\]
Now
$H^1(F_n;\mathbb{Q})$
has the ambient classes from $H^1(\Sigma_g;\mathbb{Q})$ together with
the puncture classes. Unlike the affine case, the point-pushing action
of $\pi_1(C_{n-1}(\Sigma_g))$ on this $H^1(F_n)$ is nontrivial: moving
the $i$-th marked point around an $A_r$- or $B_r$-cycle changes the
puncture class by the corresponding handle class. So the Serre spectral
sequence is no longer a constant-coefficient copy of the affine proof;
the mixed geometry of punctures and periods must be recorded.

Totaro computes exactly this correction for smooth projective
varieties. Specializing \cite[Theorem~1]{Totaro96} to a complex curve
gives the differential graded algebra generated by
$H^*(\Sigma_g^n;\mathbb{Q})$ and degree-$1$ classes $\eta_{ij}$, with
Arnold, $\eta_{ij}^2=0$, and the mixed relations
$\eta_{ij}(p_i^*\xi-p_j^*\xi)=0$, together with the differential
$d\eta_{ij}=\Delta_{ij}$ of~\eqref{eq:totaro-curve-differential}. For a
curve, the diagonal class is
\[
\Delta_{ij}
\;=\;
\omega_i+\omega_j
\;+\;
\sum_{r=1}^g
\bigl(\alpha_{r,i}\beta_{r,j}-\beta_{r,i}\alpha_{r,j}\bigr),
\]
which is exactly the Abel--Jacobi/period correction tying the puncture
generator to the ambient cohomology. Because $\Sigma_g$ is smooth and
projective, Totaro's degeneration theorem
\cite[Theorem~3]{Totaro96} identifies the rational cohomology of
$C_n(\Sigma_g)$ with the cohomology of this diagonal CDGA; no
additional Arnold-only relation among closed pairwise classes computes
the positive-genus contribution.
The resulting model is equivalent to the rational models of
K\v{r}\'{\i}\v{z}~\cite{Kriz94} and
Bezrukavnikov~\cite{Bezrukavnikov94}. In particular,
equations~\eqref{eq:totaro-curve-rel-1}--\eqref{eq:totaro-curve-rel-3}
show that the higher-genus obstruction is not additional Arnold
relations among closed $\eta_{ij}$ alone, but rather the interaction of
the puncture generators with the period classes of the curve.
\end{proof}

\begin{proposition}[Geometric realization of the universal twisting
morphism; \ClaimStatusProvedHere]
\label{prop:twisting-morphism-propagator}
\index{twisting morphism!geometric realization|textbf}
\index{propagator!as twisting morphism|textbf}
The cogenerator projection underlying the universal twisting morphism
$\pi\colon \barB^{\mathrm{ch}}(\cA) \to \cA$
\textup{(}Proposition~\textup{\ref{prop:universal-twisting-adjunction}(i))}
has the following affine simple-pole geometric component on
$\overline C_2(X)$:
\begin{equation}\label{eq:twisting-propagator-config}
\pi^{(0),\mathrm{simp}}([s^{-1}\bar{a}])(z)
\;=\;
\operatorname{Res}_{w=z}\, \bar{a}(w)\,\eta(z,w),
\end{equation}
where $\eta(z,w) = d\log(z - w)$ is the logarithmic propagator on an
affine genus-$0$ screen \textup{(}the Szeg\H{o}/prime-form propagator
with period correction replaces it on a compact positive-genus
curve\textup{)}. The Maurer--Cartan equation
$\partial(\pi) + \pi \star \pi = 0$
\textup{(}Definition~\textup{\ref{def:twisting-morphism})} has, after
projection to the same affine simple-pole component at bar degree~$2$,
the Arnold relation \textup{(}Theorem~\textup{\ref{thm:arnold-relations})}:
\[
\underbrace{\operatorname{Res}_{w=z}\,
(\bar{a}_{1\,(0)}\bar{a}_2)(w) \cdot \eta(z,w)}_{\partial(\pi)}
\;+\;
\underbrace{\operatorname{Res}_{w_1=z}\,
\operatorname{Res}_{w_2=z}\,
\bar{a}_1(w_1)\, \bar{a}_2(w_2)
\cdot \eta(z,w_1) \wedge \eta(w_1,w_2)}_{\pi \star \pi}
\;=\; 0.
\]
Neither term vanishes individually; their sum is zero by
$\eta_{12} \wedge \eta_{23} + \eta_{23} \wedge \eta_{31}
+ \eta_{31} \wedge \eta_{12} = 0$.
The affine Arnold relation is the local MC equation for the propagator;
the full twisting-morphism equation also contains the higher OPE modes,
the Koszul signs in $T^c(s^{-1}\bar A)$, and the projective or period
terms away from the affine screen. The propagator records only the
logarithmic form factor; the chiral operation supplies the algebraic
residue coefficient.
\end{proposition}

\begin{proof}
The projection $\pi\colon \barB^{\mathrm{ch}}(\cA) \to \cA$ sends
a bar element $[s^{-1}\bar{a}]$ of bar degree~$1$ to~$\bar{a}$
(the cogenerator projection) and kills all higher bar degrees.
Geometrically, bar degree~$1$ elements correspond to sections
of~$\cA$ on $\overline{C}_1(X) = X$, and the projection extracts
the coefficient of the identity operator. The simple-pole residue
formula~\eqref{eq:twisting-propagator-config} is the explicit
realization of this geometric component: the propagator $\eta(z,w)$
extracts the logarithmic residue at $w = z$, while the OPE coefficient
comes from the chiral product.

On an affine genus-$0$ collision screen, the MC equation
$\partial(\pi) + \pi \star \pi = 0$ at bar degree~$2$ reads:
$\operatorname{Res}_{w=z} d_{\barB}([s^{-1}\bar{a}_1 \mid
s^{-1}\bar{a}_2]) \cdot \eta(z,w)
+ \operatorname{Res}_{w_1=z} \operatorname{Res}_{w_2=z}
\bar{a}_1(w_1)\,\bar{a}_2(w_2) \cdot
\eta(z,w_1) \wedge \eta(w_1,w_2) = 0$.
The first term is $\partial(\pi)$ (the bar differential applied to
$\pi$); the second is $\pi \star \pi$ (the convolution square,
a double residue at a compatible codimension-$2$ collision corner).
Their affine simple-pole sum vanishes by the Arnold relation
(Theorem~\ref{thm:arnold-relations}):
$\eta_{12} \wedge \eta_{23} + \eta_{23} \wedge \eta_{31}
+ \eta_{31} \wedge \eta_{12} = 0$ applied to the triple
$(z, w_1, w_2)$.
The projection $\pi$ itself kills cogenerators of bar degree~$\geq2$,
but the convolution MC equation has nontrivial higher components after
precomposition with the bar differential. Those components are the
associativity and Borcherds identities for the higher OPE modes, not
additional statements about the scalar Arnold form.
\end{proof}

\begin{lemma}[Basic logarithmic form; \ClaimStatusProvedHere]\label{lem:basic-log-form-residue}
On a Fulton--MacPherson normal-crossing chart with collision-scale
coordinates $x_S$ for a nested collection of clusters, the form
$\eta_{ij}=d\log(z_i-z_j)$ has the local expression
\[
z_i-z_j
=
u_{ij}\prod_{\{i,j\}\subset S}x_S,
\qquad
u_{ij}\in\mathcal O^\times,
\]
and therefore
\[
\eta_{ij}
=
\sum_{\{i,j\}\subset S}d\log x_S+d\log u_{ij}.
\]
Consequently $\eta_{ij}$ has a simple logarithmic pole along every
collision divisor whose cluster contains both labels $i,j$, with
residue~$1$ on the screen chart where those two labelled directions are
separated, and no pole along collision divisors whose cluster does not
contain both labels.
\end{lemma}

\begin{proof}
Work in the holomorphic coordinates of
Theorem~\ref{thm:local-coords-boundary} for a nested collection of
clusters. Each cluster $S$ containing both $i$ and $j$ contributes one
scale coordinate $x_S$ to the difference $z_i-z_j$; clusters not
containing both labels contribute only to coordinates transverse to that
difference. The remaining factor $u_{ij}$ is invertible on the chart
where the two directions on the last common screen are distinct. Taking
$d\log$ gives the displayed formula. The Poincar\'e residue along
$D_S=\{x_S=0\}$ extracts the coefficient of $dx_S/x_S$, which is $1$
when $\{i,j\}\subset S$ and is $0$ otherwise. A holomorphic change of
coordinate multiplies $z_i-z_j$ by a unit, changing only the regular
term $d\log u_{ij}$, so the residues are intrinsic.
\end{proof}

\begin{theorem}[Residue operations; \ClaimStatusProvedHere]\label{thm:residue-operations}
\index{residue!along divisor|textbf}
For a normal crossing divisor $D = \bigcup_\alpha D_\alpha$ in $\overline{C}_n(\Sigma_g)$, there are well-defined residue maps:
\[\text{Res}_{D_\alpha}: \Omega^{\bullet}_{\overline{C}_n(\Sigma_g)}(\log D) \to \Omega^{\bullet-1}_{D_\alpha}(\log D_\alpha^\partial)\]
where $D_\alpha^\partial := D_\alpha \cap \bigcup_{\beta\neq\alpha} D_\beta$ is the induced boundary on $D_\alpha$.

These satisfy:
\begin{enumerate}
\item \emph{Leibniz rule:} $\text{Res}_{D_\alpha}(\omega \wedge \eta) = \text{Res}_{D_\alpha}(\omega) \wedge \eta|_{D_\alpha} + (-1)^{|\omega|}\omega|_{D_\alpha} \wedge \text{Res}_{D_\alpha}(\eta)$
\item \emph{Commutativity:} If $D_\alpha \cap D_\beta = \emptyset$, then $\text{Res}_{D_\alpha} \circ \text{Res}_{D_\beta} = \text{Res}_{D_\beta} \circ \text{Res}_{D_\alpha}$
\item \emph{Corner sign:} If $D_\alpha$ and $D_\beta$ meet
 transversely, the degree-$-1$ residue operators anticommute:
 \[
 \Res_{D_\alpha}\Res_{D_\beta}
 =
 -\Res_{D_\beta}\Res_{D_\alpha}.
 \]
\item \emph{Residue theorem:} If $d\omega=0$, then the residue family $\{\text{Res}_{D_\alpha}(\omega)\}_\alpha$ satisfies the cocycle condition in the residue sequence of Theorem~\ref{thm:residue-sequence} (equivalently, its total boundary residue vanishes).
\end{enumerate}
\end{theorem}

\begin{proof}
By Theorem~\ref{thm:normal-crossings}, every point of $\overline{C}_n(\Sigma_g)$
admits local coordinates $(z_1,\dots,z_m)$ in which
$D=\{z_1\cdots z_r=0\}$ is a simple normal crossing divisor.
For $i\le r$, write $D_i=\{z_i=0\}$.

In such a chart, every logarithmic form has a unique decomposition
\[
\omega=\omega_0+\sum_{i=1}^r \frac{dz_i}{z_i}\wedge \eta_i,
\]
with $\omega_0$ regular and each $\eta_i$ logarithmic along
$D_i^\partial=\{z_1\cdots \widehat{z_i}\cdots z_r=0\}\subset D_i$, with no
$dz_i/z_i$ component. Define
\[
\Res_{D_i}(\omega):=\eta_i\!\mid_{z_i=0}\in
\Omega^{\bullet-1}_{D_i}(\log D_i^\partial).
\]
Independence of choices follows from uniqueness of the decomposition, so the maps
glue sheaf-theoretically.

For the Leibniz rule, write
$\omega=\omega_0+\frac{dz_i}{z_i}\wedge\omega_1$ and
$\eta=\eta_0+\frac{dz_i}{z_i}\wedge\eta_1$.
Extracting the coefficient of $dz_i/z_i$ in $\omega\wedge\eta$ gives exactly
\[
\Res_{D_i}(\omega\wedge\eta)=
\Res_{D_i}(\omega)\wedge \eta|_{D_i}
 +(-1)^{|\omega|}\omega|_{D_i}\wedge\Res_{D_i}(\eta).
\]

If $D_i\cap D_j=\emptyset$, then locally at any point at most one of $z_i,z_j$
vanishes, so each composition
$\Res_{D_i}\circ\Res_{D_j}$ and $\Res_{D_j}\circ\Res_{D_i}$ is zero; hence they
agree globally.

If $D_i$ and $D_j$ meet transversely, the two-log part of a form is
\[
\omega=\frac{dz_i}{z_i}\wedge\frac{dz_j}{z_j}\wedge\gamma+\cdots .
\]
Changing the order of the two logarithmic factors changes sign. The
two iterated residues therefore extract the same coefficient with
opposite signs, which is the displayed anticommutation identity.

Finally, by Theorem~\ref{thm:residue-sequence}, the residue map is the middle map
of a short exact sequence of complexes. Therefore for $d\omega=0$, the residue
family lies in the kernel of the connecting morphism, i.e. it satisfies the
residue cocycle condition (equivalently, total boundary residue vanishes).
\end{proof}

\begin{proposition}[Residue computation in local coordinates; \ClaimStatusProvedHere]\label{prop:residue-local}
In the holomorphic local coordinates $(p,x_S,u,w)$ near
$D_S=\{x_S=0\}$ from Theorem~\ref{thm:local-coords-boundary}, the
residue operation is:
\[\text{Res}_{D_S}: \Omega^k(\log D_S) \to \Omega^{k-1}_{D_S}\]
given explicitly by:
\[\text{Res}_{D_S}\left(\frac{dx_S}{x_S} \wedge \alpha + \beta\right) = \alpha|_{x_S=0}\]
where $\alpha \in \Omega^{k-1}$ and $\beta \in \Omega^k$ are smooth.
On the real oriented blow-up one may write
$x_S=\epsilon e^{i\theta}$, so
$d\log x_S=d\epsilon/\epsilon+i\,d\theta$; this is the real/log
boundary expression of the same holomorphic residue coordinate.
\end{proposition}

\begin{proof}
In the chart of Theorem~\ref{thm:local-coords-boundary}, $D_S$ is the smooth
divisor $\{x_S=0\}$, and the logarithmic cotangent sheaf is generated by
$dx_S/x_S$ and the regular differentials in the remaining
coordinates. Hence every logarithmic $k$-form can be written uniquely as
\[
\omega=\frac{dx_S}{x_S}\wedge\alpha+\beta,
\]
where neither $\alpha$ nor $\beta$ contains a $dx_S/x_S$ factor.
By definition of the Poincar\'e residue for a smooth divisor, the residue is the
coefficient of $dx_S/x_S$ restricted to $x_S=0$, giving
\[
\Res_{D_S}(\omega)=\alpha|_{x_S=0}.
\]
This is exactly the displayed formula.
\end{proof}

\begin{remark}[Residues and OPE: pole absorption by the propagator]
\label{rem:residues-and-ope}
\label{rem:bar-pole-absorption}
\index{pole absorption!by propagator}
The Poincar\'e residue acts on the logarithmic form factor; the chiral
product supplies the OPE coefficient map. These are coupled in the bar
differential, but they are not the same operation. Recall the OPE:
\[
a(z)b(w)
\;=\;
\sum_{m\geq 0}\frac{(a_{(m)}b)(w)}{(z-w)^{m+1}}
+ \text{regular}.
\]
In the bar complex,
\[
\bar{B}^2(\mathcal{A}) \;=\; \mathcal{A}^{\otimes 2}
\otimes \Omega^1_{\overline{C}_2(\Sigma_g)}(\log D_{12}),
\]
the element $\alpha = \phi_i(z_1)\otimes\phi_j(z_2)\otimes \eta_{12}$
carries the propagator $\eta_{12} = d\log(z_1-z_2)$ as its form
component. The bar differential extracts the residue along $D_{12}$:
\[
d\alpha
\;=\;
\operatorname{Res}_{D_{12}}\!\left[\phi_i(z_1)\phi_j(z_2)\,\eta_{12}\right].
\]

\emph{Pole-order counting.} Put $u=z_1-z_2$. The logarithmic factor is
$du/u$, and the Poincar\'e residue extracts the coefficient of this
factor after the chiral product has expanded the OPE. Translating the
same collision packet into the spectral kernel $r(u)$ lowers the visible
pole order by one:
\[
\frac{a_{(m)}b}{u^{m+1}}\cdot \frac{du}{u}
\quad\leadsto\quad
\frac{a_{(m)}b}{u^m}
\quad\text{in the } r\text{-matrix normalization}.
\]
Thus the residue map extracts all finite OPE modes $a_{(m)}b$; the
Arnold/Jacobi comparison is the simple-pole projection $m=0$.
For the Heisenberg/affine OPE
\[
J^a(z)J^b(w) \;\sim\; \frac{k\delta^{ab}}{(z-w)^2} + \frac{f^{ab}{}_c J^c}{z-w},
\]
the double pole gives the genuine spectral term $k\,\Omega/u$; the
simple-pole bracket is the ordinary chiral product term, not an
additional singular spectral kernel. For the Virasoro OPE the
$u^{-4}$ and $u^{-2}$ terms become
\[
r^{\mathrm{Vir}}(u)=\frac{c/2}{u^3}+\frac{2T}{u},
\]
while the $u^{-1}$ derivative term is a contact/exact contribution and
does not create an independent $u^{-2}$ spectral term.
\end{remark}

\begin{theorem}[Residue sequence; \ClaimStatusProvedHere]\label{thm:residue-sequence}
There is an exact sequence of sheaves:
\[0 \to \Omega^k_{\overline{C}_n(\Sigma_g)} \to \Omega^k_{\overline{C}_n(\Sigma_g)}(\log D) \xrightarrow{\text{Res}} \bigoplus_{\alpha} \Omega^{k-1}_{D_\alpha}(\log D_\alpha^\partial) \to 0\]
where the residue map extracts the logarithmic part along each divisor component $D_\alpha$.

Exactness means that forms with log poles having zero residue along all $D_\alpha$ are smooth, and every $(k-1)$-form on the boundary $D = \bigcup D_\alpha$ arises as the residue of some logarithmic form.
\end{theorem}

\begin{proof}
Exactness is local on $\overline{C}_n(\Sigma_g)$. By
Theorem~\ref{thm:normal-crossings}, near every point there are coordinates
$(z_1,\dots,z_m)$ with $D=\{z_1\cdots z_r=0\}$.

Write $D_i=\{z_i=0\}$ for $1\le i\le r$. Every logarithmic form can be uniquely
written as
\[
\omega=\omega_0+\sum_{i=1}^r \frac{dz_i}{z_i}\wedge\eta_i,
\]
with $\omega_0$ regular and $\eta_i$ logarithmic along
$D_i^\partial=\{z_1\cdots \widehat{z_i}\cdots z_r=0\}\subset D_i$, without a
$dz_i/z_i$ term. The residue map is
\[
\Res(\omega)=\left(\eta_i|_{z_i=0}\right)_{i=1}^r.
\]

The left injection is immediate. If $\Res(\omega)=0$, then each
$\eta_i|_{z_i=0}=0$, so $\eta_i=z_i\theta_i$ for some logarithmic $\theta_i$.
Hence
\[
\frac{dz_i}{z_i}\wedge \eta_i = dz_i\wedge\theta_i
\]
is regular, and therefore $\omega$ is regular. So
$\ker(\Res)=\Omega^k_{\overline{C}_n(\Sigma_g)}$.

For surjectivity, take a local section
$(\rho_i)_{i=1}^r\in\bigoplus_i\Omega^{k-1}_{D_i}(\log D_i^\partial)$.
Choose local lifts $\widetilde{\rho}_i$ to ambient logarithmic forms with no
$dz_i/z_i$ factor, and set
\[
\omega:=\sum_{i=1}^r \frac{dz_i}{z_i}\wedge \widetilde{\rho}_i.
\]
Then $\Res_{D_i}(\omega)=\rho_i$ for each $i$, proving local surjectivity. Since
exactness of sheaf sequences is local, the displayed sequence is exact globally.
\end{proof}

\subsubsection{Functoriality and universal properties}

\begin{theorem}[Functoriality of FM compactification; \ClaimStatusProvedElsewhere]\label{thm:FM-functorial}
The Fulton--MacPherson compactification is functorial in the following
sense \cite{FM94}:

\begin{enumerate}
\item \emph{For embeddings:} If $U \subseteq \Sigma_g$ is an open subset, there is a natural embedding:
\[\overline{C}_n(U) \hookrightarrow \overline{C}_n(\Sigma_g)\]
compatible with boundary stratification.

\item \emph{For open embeddings:} If $f: \Sigma_g \hookrightarrow \Sigma_{g'}$ is an open embedding, there is an induced map:
\[\overline{f^{(n)}}: \overline{C}_n(\Sigma_g) \to \overline{C}_n(\Sigma_{g'})\]
sending $D_S \to D_S$ (same collision pattern). (General smooth maps need not preserve the configuration condition; injectivity on pairs is required.)

\item \emph{For automorphisms:} The group $\text{Aut}(\Sigma_g)$ acts on $\overline{C}_n(\Sigma_g)$ preserving stratification.

\item \emph{For products:} There is a natural product structure:
\[\overline{C}_m(\Sigma_g) \times \overline{C}_n(\Sigma_g) \hookrightarrow \overline{C}_{m+n}(\Sigma_g)\]
(away from mixed collision loci)
\end{enumerate}
\end{theorem}

\begin{remark}[Provenance and citation]
This theorem is imported and treated as \ClaimStatusProvedElsewhere. The
functorial construction of the Fulton--MacPherson compactification and its
boundary-strata compatibility are established in \cite{FM94}; the statements
here are the corresponding specialization to smooth curves/surfaces.
\end{remark}

\begin{theorem}[Universal property: FM right-module structure {\cite{FM94,LV12}}; \ClaimStatusProvedElsewhere]\label{thm:FM-operad}
For a fixed smooth curve $\Sigma_g$, the collection
$\{\overline{C}_n(\Sigma_g)\}_{n\geq 0}$ is a right module over the
genus-$0$ Fulton--MacPherson operad of tangent screens. Equivalently,
each collision divisor carries a gluing map
\[
\overline{C}_{n-|S|+1}(\Sigma_g)
\times
\bigl(\overline{C}_{|S|}(\mathbb{A}^1)/
\mathrm{Aff}(\mathbb{A}^1)\bigr)
\longrightarrow
D_S\subset \overline{C}_n(\Sigma_g),
\]
and these maps are associative along nested collision strata. When
$\Sigma_g=\mathbb{A}^1$ or a formal disk with a chosen coordinate, this
right-module structure is the usual FM operad structure. For a general
fixed curve there is no intrinsic operation inserting an entire global
configuration on $\Sigma_g$ into a point of another such configuration;
the intrinsic inserted object is the affine tangent screen.
\end{theorem}

\begin{remark}[Provenance and citation]
This theorem is imported and treated as \ClaimStatusProvedElsewhere. The
compactified configuration-space gluing formalism is part of the
Fulton--MacPherson framework \cite{FM94}, and the operadic/right-module
packaging used here is standard in operad references such as \cite{LV12}.
\end{remark}

\begin{remark}[Chiral operad structure]\label{rem:chiral-operad}
The FM right-module structure of $\{\overline{C}_n(\Sigma_g)\}$ over
the affine tangent-screen operad is the geometric foundation for the
\emph{chiral operad} structure in Beilinson--Drinfeld \cite{BD04}.

Specifically, the spaces of logarithmic forms
\[\mathcal{P}^{\text{ch}}_n(\Sigma_g) = H^0(\overline{C}_n(\Sigma_g), \Omega^n_{\overline{C}_n(\Sigma_g)}(\log D))\]
assemble into the factorization/right-module operations obtained by
restriction to collision screens; chiral algebras are algebras over the
resulting chiral structure in the appropriate $\infty$-categorical
sense.
\end{remark}

\subsubsection{Connection to factorization homology}

\begin{theorem}[Factorization homology via configuration spaces {\cite{AF15,CG17,BD04}}; \ClaimStatusProvedElsewhere]\label{thm:fact-homology}
\label{thm:bar-factorization-homology}
For a chiral algebra $\mathcal{A}$ on $\Sigma_g$, the factorization homology is computed via:
\[\int_{\Sigma_g} \mathcal{A} = \text{colim}_n \left[ \mathcal{A}^{\boxtimes n} \otimes_{\mathcal{D}_{\overline{C}_n(\Sigma_g)}} \mathcal{O}_{\overline{C}_n(\Sigma_g)} \right]\]

where:
\begin{itemize}
\item $\mathcal{A}^{\boxtimes n} = \mathcal{A} \boxtimes \cdots \boxtimes \mathcal{A}$ is the external tensor product on $\Sigma_g^n$
\item $\mathcal{D}_{\overline{C}_n(\Sigma_g)}$ is the sheaf of differential operators on $\overline{C}_n(\Sigma_g)$
\item The colimit is over the Ran/factorization structure maps induced
by disjoint union and collision-screen gluing
\end{itemize}
\end{theorem}

\begin{remark}[Provenance and citation]
This theorem is imported and treated as \ClaimStatusProvedElsewhere. The
factorization-homology construction and its colimit/locality formalism are
standard in \cite{AF15,CG17}; the chiral-algebra specialization used here is in
the Beilinson--Drinfeld framework \cite[Chapter~3]{BD04}.
\end{remark}

\begin{remark}[Ran space]\label{rem:ran-space}
An alternative perspective uses the \emph{Ran space} $\text{Ran}(\Sigma_g)$:
\[\text{Ran}(\Sigma_g) = \{S \subset \Sigma_g : 0 < |S| < \infty\}\]
equipped with the colimit topology from the natural maps $C_n(\Sigma_g)/S_n \hookrightarrow \text{Ran}(\Sigma_g)$
(in particular, points may collide: the topology allows sets of different cardinalities to be adjacent)

The Ran space parametrizes \emph{finite unordered subsets} of $\Sigma_g$. A chiral algebra structure on $\mathcal{A}$ is equivalent to a factorization algebra $\mathcal{A}_{\text{Ran}}$ on $\text{Ran}(\Sigma_g)$ satisfying chiral locality conditions (encoded by the OPE).

The Fulton--MacPherson compactification provides a ``partial compactification'' of Ran space, adding boundary strata for collision patterns.
\end{remark}

\begin{example}[Factorization for Heisenberg]\label{ex:fact-heisenberg}
For the Heisenberg chiral algebra $\mathcal{H}$ at level $k$:
\[\int_{\Sigma_g} \mathcal{H} \cong \text{Fock space at level } k\]

More precisely:
\begin{itemize}
\item At genus 0: $\int_{\mathbb{P}^1} \mathcal{H} \cong \mathbb{C}[x]$ (polynomial algebra)
\item At genus 1: $\int_{\Sigma_1} \mathcal{H} \cong \text{Hilbert space of } k \text{ particles on } \Sigma_1$
\item At genus $g$: Includes contributions from all homology cycles
\end{itemize}

The computation uses:
\[\int_{\Sigma_g} \mathcal{H} = \text{colim}_n \left[\mathcal{H}^{\boxtimes n} \text{ with Heisenberg OPE along collisions}\right]\]
The OPE $J(z)J(w) \sim \frac{k}{(z-w)^2}$ determines how factors merge at boundaries of $\overline{C}_n(\Sigma_g)$.
\end{example}

The Fulton--MacPherson compactification $\overline{C}_n(\Sigma_g)$ (for a \emph{fixed} curve $\Sigma_g$) stratifies by nestings:
\[\overline{C}_n(\Sigma_g) = \coprod_{S \in \mathrm{Nest}(\{1,\ldots,n\})} C_S\]
where $\mathrm{Nest}(\{1,\ldots,n\})$ denotes the poset of nested collections of subsets of $\{1,\ldots,n\}$, recording which groups of points collide simultaneously. (By contrast, the strata of $\overline{\mathcal{M}}_{g,n}$ are indexed by stable graphs $\mathcal{G}_{g,n}$, which involve \emph{curve} degenerations.)

\section{Logarithmic FM compactification on punctured curves}
\label{sec:log-fm-punctured}
\index{Fulton--MacPherson compactification!logarithmic}
\index{punctured curve!configuration source}

The logarithmic FM compactification developed below is the
ambient geometry on which Theorem~A in its
$(\infty,2)$-properad-level formulation
(Theorem~\ref{thm:A-infinity-2} in
\texttt{theorem\_A\_infinity\_2.tex}) operates. Conceptually,
the Arnold connection is enlarged: in addition to the
interior-collision strata $D_S^{\mathrm{FM}}$ ($|S|\geq2$, where
the Arnold relation lives), the logarithmic strata over the fixed
divisor~$D$ contribute residue faces at which the propagator extracts
the chiral \emph{module} action $\rho_\alpha$ rather than the chiral
algebra OPE. Mok's construction records these faces by grid
subdivisions and planted forests refining the pairwise collision labels.
The resulting residue calculus is the geometric input for
conformal-block recovery
(Corollary~\ref{cor:conformal-blocks-punctured-bar}).

When the base curve $X$ carries marked points $p_1, \ldots, p_r$
(``punctures''), the configuration space $C_m(X \setminus \{p_1,
\ldots, p_r\})$ of $m$ points on the complement is no longer an
open subset of $X^m$ (it is open in $(X \setminus D)^m$ where
$D = \{p_1, \ldots, p_r\}$). The FM compactification must be
modified to account for two geometric mechanisms:
\begin{enumerate}[label=(\roman*)]
\item \emph{FM collisions}: clusters $S\subset\{1,\ldots,m\}$,
 $|S|\geq2$, of moving points collide, encoded by the usual
 Fulton--MacPherson blow-ups;
\item \emph{logarithmic approach to the boundary}: moving points
 approach strata of the fixed divisor~$D$, encoded by toroidal
 monomial modifications of the Artin fan and their pullbacks to the
 configuration space, rather than by blow-ups of~$D$ itself.
\end{enumerate}

\begin{definition}[Logarithmic FM compactification \cite{Mok25}]
\label{def:log-fm-compactification}
\index{logarithmic FM compactification|textbf}
Let $X$ be a smooth projective curve, $D = p_1 + \cdots + p_r$ an
effective divisor. The \emph{logarithmic Fulton--MacPherson
compactification of the fixed pointed curve $(X,D)$},
$\mathrm{FM}_m(X \mathbin{|} D)$, is a compactification of the
labelled configuration space $C_m(X \setminus D)$ with the following
properties:
\begin{enumerate}
\item There is an open embedding
 $j\colon C_m(X \setminus D) \hookrightarrow
 \mathrm{FM}_m(X \mathbin{|} D)$
 with dense image.
\item The boundary
 $\partial\mathrm{FM}_m(X \mathbin{|} D)$ is a logarithmic
 normal-crossings boundary in the toroidal sense. Its local strata
 are indexed by Mok planted-forest combinatorial types
 \[
 (\mathcal F_{\mathcal P}\to\mathcal F_{\Sigma_X},\,m;\,
 T_{v_1},\ldots,T_{v_k}),
 \]
 where $(\mathcal P,m)$ is an $m$-marked grid subdivision of the
 tropicalisation~$\Sigma_X$ of the pair $(X,D)$, and each
 $T_{v_i}$ is a rooted tree whose legs are the markings supported
 at the vertex~$v_i$.
\item The classical FM collision divisors are the special case in
 which the grid subdivision is trivial and a rooted tree records a
 collision cluster. The strata over~$D$ are the grid-subdivision
 part of the same Mok type. Thus pairwise labels $D_{ij}$ and
 single-point puncture labels $D_{i,\alpha}$ are only local
 first-order shadows of the full planted-forest stratification.
\end{enumerate}
\end{definition}

\begin{remark}[Fixed versus relative log FM spaces]
\label{rem:log-fm-fixed-vs-relative}
Definition~\ref{def:log-fm-compactification} keeps the pointed
curve $(X,D)$ fixed. Its boundary records only configuration
collisions and logarithmic approach to the fixed divisor~$D$,
organised by Mok planted forests and grid subdivisions. It does
not contain additional components from degeneration of the
underlying curve. Those curve-degeneration strata appear only in
the relative logarithmic Fulton--MacPherson compactification of a
semistable family, and in this manuscript only conditionally after
descent to the universal pointed stable family over
$\overline{\mathcal{M}}_{g,n}$.
That relative space, not the fixed-curve space of
Definition~\ref{def:log-fm-compactification}, is the one used in
the ambient square-zero argument of
Theorem~\ref{thm:ambient-d-squared-zero}.
\end{remark}

\begin{example}[Log FM recovers moduli of rational curves;
\ClaimStatusProvedElsewhere]
\label{ex:log-fm-moduli-rational}
\index{logarithmic FM compactification!moduli of rational curves}
For $X = \mathbb{P}^1$ with $D = \{0, 1, \infty\}$ (three marked points),
Mok~\cite[\S3.4.1]{Mok25} proves the identification
\begin{equation}\label{eq:log-fm-m0n}
\operatorname{FM}_{n-3}\bigl(\mathbb{P}^1 \mathbin{|}
 0{+}1{+}\infty\bigr)
\;\cong\;
\overline{\cM}_{0,n}.
\end{equation}
This gives the log~FM space a direct modular interpretation:
the $n-3$ configuration points together with the $3$~punctures
produce $n$ marked points on~$\mathbb{P}^1$, and the log~FM
compactification is precisely the Deligne--Mumford--Knudsen
compactification. The FM-type boundary strata record
bubble trees (collisions among moving points), while the
puncture-collision strata correspond, under
\eqref{eq:log-fm-m0n}, to the boundary components where a
moving point bubbles off with a fixed puncture. Together they
reproduce the full boundary stratification of
$\overline{\cM}_{0,n}$ by dual trees.
This identification means that the planted-forest coefficient
algebra $\mathbb{G}_{\mathrm{pf}}$
(Definition~\ref{def:planted-forest-coefficient-algebra}),
when specialised to $\mathbb{P}^1$ with three punctures, encodes
exactly the boundary combinatorics of
$\overline{\cM}_{0,n}$, the same data that controls the
genus-$0$ convolution bracket on
$\mathfrak{g}^{\mathrm{mod}}_\cA$.
\end{example}

\begin{theorem}[Bar complex on punctured curves; \ClaimStatusConditional]
\label{thm:bar-punctured-curve}
\index{bar complex!punctured curve}
Let $\cA$ be a chiral algebra on $X$, and let
$\mathcal{M}_1, \ldots, \mathcal{M}_r$ be $\cA$-modules
supported at $p_1, \ldots, p_r$. The bar complex with module
insertions is:
\begin{equation}\label{eq:bar-punctured}
\bar{B}_n(\cA;\, \mathcal{M}_1, \ldots, \mathcal{M}_r;\, X)
= \Gamma\bigl(\mathrm{FM}_{n}(X \mathbin{|} D),\;
 j_* j^*(\cA^{\boxtimes n})
 \otimes \mathcal{M}_1 \otimes \cdots \otimes \mathcal{M}_r
 \otimes \Omega^n_{\log}\bigr)
\end{equation}
where $\Omega^n_{\log} = \Omega^n_{\mathrm{FM}_{n}(X|D)}
(\log \partial)$ is the sheaf of logarithmic $n$-forms
on $\mathrm{FM}_{n}(X \mathbin{|} D)$. The $r$ modules are fixed
boundary insertions at the divisor $D=\{p_1,\ldots,p_r\}$; they are not
additional moving configuration points.
The differential
$d = d_{\mathrm{FM}} + d_{\mathrm{punc}}$
has two components:
\begin{enumerate}[label=\textup{(\alph*)}]
\item $d_{\mathrm{FM}}$: residues along FM collision faces indexed by
 collision clusters $S$, extracting the iterated $\cA$-algebra OPE
 \textup{(}with the pairwise faces giving the usual Arnold
 computation of Theorem~\textup{\ref{thm:bar-nilpotency-complete})};
\item $d_{\mathrm{punc}}$: logarithmic residues along faces where a
 moving insertion approaches the fixed boundary component
 $p_\alpha$, extracting the
 $\cA$-module action $\rho_\alpha\colon \cA \boxtimes
 \mathcal{M}_\alpha \to \Delta_* \mathcal{M}_\alpha$.
\end{enumerate}
\end{theorem}

\begin{proof}
The key identity is $d^2 = 0$, which follows from
the fact that $\mathrm{FM}_{n}(X \mathbin{|} D)$ is attached to the
fixed pointed curve $(X,D)$. Its boundary has collision faces and log
approach-to-$D$ faces, but no curve-degeneration faces. The
curve-degeneration boundary appears only in the relative
universal-family space used in
Theorem~\ref{thm:ambient-d-squared-zero}. Thus the double-residue
cancellations split into the following cases:
\begin{enumerate}
\item \emph{FM-FM:} on a pairwise collision screen the
 $\operatorname{Res}_{D_{ij}}\operatorname{Res}_{D_{jk}}$ terms cancel
 by the Arnold
 relation $\eta_{ij} \wedge \eta_{jk} + \eta_{jk} \wedge
 \eta_{ki} + \eta_{ki} \wedge \eta_{ij} = 0$, exactly as
 in the unpunctured case
 (Theorem~\ref{thm:bar-nilpotency-complete}). Nested collision faces
 are handled by the same boundary-of-boundary cancellation in the
 Fulton--MacPherson normal-crossings stratification.
\item \emph{FM-puncture:} the mixed residue where a collision cluster
 approaches $p_\alpha$ involves the OPE
 $\cA(z_i) \cdot \cA(z_j)$ followed by the module action
 $\rho_\alpha$. The cancellation uses the module axiom:
 the $\cA$-module action is associative, so
 $\rho_\alpha \circ (m \boxtimes \mathrm{id}_{\mathcal{M}})
 = \rho_\alpha \circ (\mathrm{id}_\cA \boxtimes \rho_\alpha)$
 on the triple-collision locus $z_i = z_j = p_\alpha$.
\item \emph{Puncture-puncture:} For $\alpha \neq \beta$,
 $D_{i,\alpha}^{\mathrm{punc}} \cap D_{j,\beta}^{\mathrm{punc}}$
 is empty when $i = j$ (a single point cannot approach two
 distinct punctures), so these terms vanish. For $i \neq j$,
 the normal-crossings property gives commutation of the underlying
 restrictions and anticommutation of the odd residue operators in the
 boundary differential:
 $\operatorname{Res}_{D_{i,\alpha}}
 \operatorname{Res}_{D_{j,\beta}}
 =-\operatorname{Res}_{D_{j,\beta}}
 \operatorname{Res}_{D_{i,\alpha}}$
 after the orientation/Koszul sign is included.
\end{enumerate}
Each case contributes zero to $d^2$, establishing
$d^2 = 0$.
\end{proof}

\begin{corollary}[Conformal blocks from punctured bar complex;
\ClaimStatusConditional]
\label{cor:conformal-blocks-punctured-bar}
The zeroth cohomology of the punctured bar complex recovers
the space of conformal blocks:
\[
H^0(\bar{B}_\bullet(\cA;\, \vec{\mathcal{M}};\, X))
\;\cong\;
\mathbb{V}(\cA;\, \mathcal{M}_1, \ldots, \mathcal{M}_r;\, X).
\]
\end{corollary}

\begin{proof}
The zeroth bar cohomology computes the space of coinvariants:
elements of $\mathcal{M}_1 \otimes \cdots \otimes \mathcal{M}_r$
modulo the $\cA(X \setminus \{p_1, \ldots, p_r\})$-action.
By definition, this is $\mathbb{V}(\cA;\, \vec{\mathcal{M}};\, X)$
(Definition~\ref{def:conformal-blocks}).
The identification follows from the quasi-isomorphism
of the bar complex with the Čech resolution of the
coinvariant functor, applied to the Weiss cover
of $X \setminus D$ subordinate to the FM stratification.
\end{proof}

\begin{example}[Kac--Moody modules on \texorpdfstring{$\mathbb{P}^1$}{P1} with
\texorpdfstring{$3$}{3} punctures]\label{ex:log-fm-km}
For $\cA = \widehat{\mathfrak{g}}_k$ on $\mathbb{P}^1$ with
simple modules $L(\Lambda_i)$ at $0, 1, \infty$:
\[
\bar{B}_1(\widehat{\mathfrak{g}}_k;\,
L(\Lambda_1), L(\Lambda_2), L(\Lambda_3);\, \mathbb{P}^1)
= \Gamma(\mathrm{FM}_1(\mathbb{P}^1 \mathbin{|} \{0,1,\infty\}),\;
j_*j^*(\widehat{\mathfrak{g}}_k)
\otimes L(\Lambda_1)\otimes L(\Lambda_2)\otimes L(\Lambda_3)
\otimes \Omega^1_{\log})
\]
The bar differential extracts:
\begin{itemize}
\item $d_{\mathrm{FM}}$: the algebra OPE
 $J^a(z_1) J^b(z_2) \sim k\delta^{ab}/(z_1 - z_2)^2
 + f^{ab}{}_c J^c/(z_1 - z_2)$;
\item $d_{\mathrm{punc}}$: the module action
 $J^a(z) \cdot v_\alpha
 = \operatorname{Res}_{z=p_\alpha}
 [J^a(z) \cdot v_\alpha \, dz/(z - p_\alpha)]$
 for $v_\alpha \in L(\Lambda_\alpha)$.
\end{itemize}
The cohomology $H^0$ computes the fusion coefficient
$N_{\Lambda_1 \Lambda_2}^{\Lambda_3^*}$
\textup{(}Theorem~\textup{\ref{thm:verlinde-bar})}.
\end{example}

\section{Bordered Fulton--MacPherson compactification}
\label{sec:bordered-fm}
\index{Fulton--MacPherson compactification!bordered|textbf}
\index{bordered FM compactification|textbf}
\index{open-closed configuration space|textbf}

The logarithmic FM compactification of \S\ref{sec:log-fm-punctured}
resolves collisions of interior points with fixed punctures. The
bordered variant developed in this section adds a second geometric
layer: the punctures $p_1, \ldots, p_r$ are promoted to
\emph{boundary circles} carrying ordered boundary insertions, and
the compactification records point--point, boundary--boundary, and
interior--boundary interactions.
The resulting \emph{bordered FM compactification}
$\overline{C}{}^{\,\mathrm{oc}}_{k,\vec{m}}(X, D, \tau)$ is the
geometric carrier of the open-closed bar complex, and its
codimension-one boundary decomposes into exactly four types: interior
collision, consecutive boundary collision, mixed bubbling, and nodal
clutching. The Stokes theorem on this boundary is the geometric
origin of the modular Maurer--Cartan equation with clutching
(\S\ref{sec:modular-mc-clutching}).

\subsection{Tangential log curves}
\label{subsec:tangential-log-curves}
\index{tangential log curve|textbf}
\index{real oriented blowup|textbf}

\begin{definition}[Tangential log curve]
\label{def:tangential-log-curve}
A \emph{tangential log curve} is a triple $(X, D, \tau)$ where:
\begin{enumerate}[label=\textup{(\roman*)}]
\item $X$ is a smooth projective curve over~$\mathbb{C}$
 of genus~$g$;
\item $D = p_1 + \cdots + p_r$ is an effective reduced divisor
 (the \emph{puncture locus});
\item $\tau = (\tau_1, \ldots, \tau_r)$ is a collection of
 \emph{tangential basepoints}: for each $j = 1, \ldots, r$,
 a nonzero tangent direction
 $\tau_j \in (T_{p_j} X) \setminus \{0\}$, well-defined up to
 positive real scaling:
 $\tau_j \in (T_{p_j} X \setminus \{0\}) / \mathbb{R}_{>0}
 \cong S^1_{p_j}$.
\end{enumerate}
The tangential basepoints determine preferred points on the
boundary circles of the real oriented blowup
(Construction~\ref{constr:real-oriented-blowup}). Two
tangential log curves $(X, D, \tau)$ and $(X, D, \tau')$ are
\emph{equivalent} if they differ by a positive-real rotation
at each puncture.
\end{definition}

\begin{construction}[Real oriented blowup]
\label{constr:real-oriented-blowup}
\index{real oriented blowup!construction|textbf}
Let $(X, D, \tau)$ be a tangential log curve. The
\emph{real oriented blowup} of~$X$ along~$D$ is the smooth
manifold with boundary
\begin{equation}\label{eq:real-oriented-blowup}
\widetilde{X}_D
\;:=\;
\bigl[X \setminus D\bigr]
\;\cup\;
\bigsqcup_{j=1}^{r} S^1_{p_j},
\end{equation}
constructed as follows. Choose a local holomorphic coordinate
$z_j$ centered at $p_j$ (so $z_j(p_j) = 0$); write
$z_j = \rho_j e^{i\theta_j}$ in polar form with
$\rho_j \geq 0$ and $\theta_j \in \mathbb{R} / 2\pi\mathbb{Z}$.
\begin{enumerate}[label=\textup{(\arabic*)}]
\item The underlying set of $\widetilde{X}_D$ is obtained by
 replacing each puncture $p_j$ with the circle of directions:
 \[
 \widetilde{X}_D
 \;=\;
 (X \setminus D)
 \;\sqcup\;
 \bigsqcup_{j=1}^{r}
 \bigl\{(\rho_j, \theta_j) : \rho_j = 0,\;
 \theta_j \in S^1\bigr\}.
 \]
\item The smooth structure near
 $S^1_{p_j}$ is given by the coordinates
 $(\rho_j, \theta_j)$ with $\rho_j \geq 0$: the chart
 $U_j = \{|z_j| < \epsilon\}$ in~$X$ lifts to
 $\widetilde{U}_j = [0, \epsilon) \times S^1 \hookrightarrow
 \widetilde{X}_D$, making $\rho_j$ a boundary-defining function.
\item The boundary of $\widetilde{X}_D$ is the disjoint union of
 circles:
 \begin{equation}\label{eq:boundary-circles}
 \partial \widetilde{X}_D
 \;=\;
 \bigsqcup_{j=1}^{r} S^1_{p_j}
 \;\cong\;
 \bigsqcup_{j=1}^{r}
 \bigl\{\rho_j = 0\bigr\}.
 \end{equation}
\item The tangential basepoint $\tau_j$ determines a preferred
 point $\star_j \in S^1_{p_j}$: the direction $\tau_j$ is the
 unit tangent vector at $\theta_j = 0$ in the local coordinate
 $z_j$ chosen so that $\partial/\partial z_j|_{p_j}$ is
 positively aligned with~$\tau_j$. The identification is:
 \begin{equation}\label{eq:basepoint-identification}
 \star_j
 \;=\;
 \lim_{\rho_j \to 0^+}
 \bigl(\rho_j,\, \arg(z_j / \tau_j)\bigr)
 \;=\;
 (0,\, 0)
 \;\in\;
 S^1_{p_j}.
 \end{equation}
\end{enumerate}
The result is independent of the choice of local coordinate $z_j$
up to diffeomorphism: a coordinate change $z_j \mapsto w_j(z_j)$
with $w_j(0) = 0$ and $w_j'(0) \neq 0$ induces a diffeomorphism
of $[0, \epsilon) \times S^1$ that rotates $S^1_{p_j}$ by
$\arg(w_j'(0))$ and rescales $\rho_j$ to first order. The
tangential basepoint absorbs this rotation.
\end{construction}

\begin{definition}[Boundary intervals]
\label{def:boundary-intervals}
\index{boundary interval|textbf}
For a tangential log curve $(X, D, \tau)$, the
\emph{boundary interval} at $p_j$ is the complement of the
tangential basepoint in the boundary circle:
\begin{equation}\label{eq:boundary-interval}
I_{p_j}
\;:=\;
S^1_{p_j} \setminus \{\star_j\}
\;\cong\;
(0, 2\pi)
\;\cong\;
\mathbb{R}.
\end{equation}
The interval $I_{p_j}$ carries the cyclic orientation induced from
the complex orientation of~$X$: the boundary circle $S^1_{p_j}$
inherits the orientation where $\theta_j$ increases
counterclockwise, and removing the basepoint
$\star_j = \{0\}$ produces a \emph{totally ordered} interval
with the standard orientation of~$\mathbb{R}$.
\end{definition}

\begin{remark}[Geometric meaning of the interval]
\label{rem:geometric-interval-meaning}
The boundary interval $I_{p_j}$ parametrizes the angular
directions from which interior points can approach the
puncture~$p_j$, with the basepoint direction removed to break
the cyclic symmetry and obtain a linear ordering. In the
open-string interpretation, $I_{p_j}$ is the worldsheet boundary
associated to the $j$-th D-brane, and the basepoint $\star_j$ is
the distinguished insertion point for the identity of the
boundary algebra.
\end{remark}

\subsection{Mixed configuration spaces}
\label{subsec:mixed-config-spaces}
\index{mixed configuration space|textbf}
\index{configuration space!open-closed}

\begin{definition}[Mixed open-closed configuration space]
\label{def:mixed-config-space}
Let $(X, D, \tau)$ be a tangential log curve with $D = p_1 +
\cdots + p_r$, and let $\vec{m} = (m_1, \ldots, m_r) \in
\mathbb{Z}_{\geq 0}^r$ be a vector of boundary multiplicities.
The \emph{mixed open-closed configuration space} is
\begin{equation}\label{eq:mixed-config-space}
\operatorname{Conf}^{\mathrm{oc}}_{k, \vec{m}}(X, D, \tau)
\;:=\;
C_k(X \setminus D)
\;\times\;
\prod_{j=1}^{r}
\operatorname{Conf}^{\mathrm{ord}}_{m_j}(I_{p_j}),
\end{equation}
where:
\begin{itemize}
\item $C_k(X \setminus D)$ is the configuration space of
 $k$~distinct ordered \emph{interior} (``closed-string'') points
 in the complement $X \setminus D$;
\item $\operatorname{Conf}^{\mathrm{ord}}_{m_j}(I_{p_j})$ is the
 configuration space of $m_j$~distinct \emph{ordered}
 \emph{boundary} (``open-string'') points in $I_{p_j}$:
 \begin{equation}\label{eq:ordered-boundary-config}
 \operatorname{Conf}^{\mathrm{ord}}_{m_j}(I_{p_j})
 \;=\;
 \bigl\{(t_1^{(j)}, \ldots, t_{m_j}^{(j)}) \in
 I_{p_j}^{m_j}
 \,:\,
 0 < t_1^{(j)} < t_2^{(j)} < \cdots < t_{m_j}^{(j)} < 2\pi
 \bigr\}.
 \end{equation}
 The ordering is with respect to the linear coordinate
 $\theta_j \in (0, 2\pi) \cong I_{p_j}$ from
 Definition~\ref{def:boundary-intervals}.
\end{itemize}
The total dimension is
$\dim_{\mathbb{R}}\,
\operatorname{Conf}^{\mathrm{oc}}_{k,\vec{m}}
= 2k + |\vec{m}|$,
where $|\vec{m}| = m_1 + \cdots + m_r$.
\end{definition}

\begin{remark}[Boundary ordering and the associahedron]
\label{rem:boundary-ordering-associahedron}
\index{associahedron!boundary points}
The ordered boundary configuration space
$\operatorname{Conf}^{\mathrm{ord}}_{m}(I_p)$ is contractible:
it is the interior of the standard simplex
$\{0 < t_1 < \cdots < t_m < 2\pi\} \cong
\operatorname{int}(\Delta^{m-1})$. The compactification of this
simplex is the $A_\infty$ associahedron $K_m$: boundary strata
correspond to consecutive sub-blocks of boundary points colliding
along the interval. The global geometry therefore separates into a
\emph{contractible boundary factor} (the simplex) and an
\emph{interesting interior factor} (the FM compactification of the
interior points).
\end{remark}

\subsection{The bordered FM compactification}
\label{subsec:bordered-fm-construction}
\index{bordered FM compactification!construction|textbf}

\begin{construction}[Bordered FM compactification]
\label{constr:bordered-fm}
Let $(X, D, \tau)$ be a tangential log curve with
$D = p_1 + \cdots + p_r$, let $k \geq 0$ be the number of
interior points, and let $\vec{m} = (m_1, \ldots, m_r)$ be the
vector of boundary multiplicities. Set $M = |\vec{m}| = \sum_j m_j$.
The \emph{bordered FM compactification}
$\overline{C}{}^{\,\mathrm{oc}}_{k,\vec{m}}(X, D, \tau)$ is
constructed by the following iterated blowup procedure,
extending the classical FM construction
(Theorem~\ref{thm:FM}) and the logarithmic variant
(Definition~\ref{def:log-fm-compactification}).

\medskip
\noindent\textsc{Step~1: The ambient partial compactification.}
Begin with the product
\begin{equation}\label{eq:bordered-ambient}
\mathring{Y}
\;:=\;
(X \setminus D)^k
\;\times\;
\prod_{j=1}^{r}
I_{p_j}^{m_j}.
\end{equation}
This is an open subset of
$\widetilde{X}_D^k \times \prod_j [0, 2\pi]^{m_j}$. To
compactify, first close the boundary factors: replace each open
interval $I_{p_j} \cong (0, 2\pi)$ by the closed interval
$\overline{I}_{p_j} = [0, 2\pi]$, and replace $(X \setminus D)^k$
by the logarithmic compactification
$\operatorname{FM}_k(X \mathbin{|} D)$ of
Definition~\ref{def:log-fm-compactification}. This gives the
\emph{rough compactification}:
\begin{equation}\label{eq:bordered-rough}
Y
\;:=\;
\operatorname{FM}_k(X \mathbin{|} D)
\;\times\;
\prod_{j=1}^{r}
\overline{I}_{p_j}^{\,m_j}.
\end{equation}
The space $Y$ is compact but has non-transverse singularities: when
two interior points collide while simultaneously a boundary point
approaches the basepoint, the resulting degeneration is not
normal-crossing.

\medskip
\noindent\textsc{Step~2: Interior blowups (FM type).}
For each subset $S \subseteq \{1, \ldots, k\}$ with $|S| \geq 2$,
the diagonal $\Delta_S^{\mathrm{int}}
= \{z_i = z_j \text{ for all } i, j \in S\}$ inside
$\operatorname{FM}_k(X \mathbin{|} D)$ has already been resolved
by the FM construction. No further blowups are needed for purely
interior collisions.

\medskip
\noindent\textsc{Step~3: Boundary blowups (associahedral type).}
For each boundary component~$j$ and each \emph{consecutive
block} $B = \{a, a{+}1, \ldots, b\} \subseteq
\{1, \ldots, m_j\}$ with $|B| \geq 2$,
the diagonal
\begin{equation}\label{eq:consecutive-diagonal}
\Delta_B^{\mathrm{bdy}}
\;=\;
\bigl\{t_a^{(j)} = t_{a+1}^{(j)} = \cdots = t_b^{(j)}\bigr\}
\;\subset\;
\overline{I}_{p_j}^{\,m_j}
\end{equation}
is a smooth submanifold of codimension $|B| - 1$. Blow up these
diagonals in decreasing order of $|B|$
(equivalently, decreasing order of codimension).
The blowups only involve \emph{consecutive} blocks because the
boundary points are ordered: a non-consecutive collision (e.g.,
$t_1 = t_3$ with $t_2$ distinct) does not occur for ordered
configurations. The result is the
\emph{Stasheff associahedron}~$K_{m_j}$: the compactified moduli
space of $m_j$ ordered points on~$\mathbb{R}$.

\medskip
\noindent\textsc{Step~4: Mixed blowups (open-closed interaction).}
For each interior point $i \in \{1, \ldots, k\}$ and each
boundary component~$j \in \{1, \ldots, r\}$ with associated
boundary points $B_j = \{t_1^{(j)}, \ldots, t_{m_j}^{(j)}\}$,
and for each subset $S \subseteq \{1, \ldots, k\}$ with
$i \in S$ and each \emph{nonempty} consecutive block
$B \subseteq \{1, \ldots, m_j\}$ (so that $|S| \geq 1$
and $|B| \geq 1$; the purely interior and purely boundary
cases are already resolved in Steps~2 and~3), define the
\emph{mixed collision locus}:
\begin{equation}\label{eq:mixed-collision}
\Delta_{S, B, j}^{\mathrm{mix}}
\;:=\;
\bigl\{
z_i = p_j \text{ for all } i \in S,\;
t_a^{(j)} = \star_j \text{ for all } a \in B
\bigr\}
\;\subset\; Y.
\end{equation}
This records a degeneration where interior points in~$S$
approach the puncture~$p_j$ while boundary points in the
block~$B$ simultaneously approach the basepoint~$\star_j$.
Blow up these mixed loci in the order dictated by decreasing
codimension $|S| + |B|$, after the purely interior and
purely boundary blowups.

\medskip
\noindent\textsc{Step~5: Nodal blowups (clutching type).}
When the base curve $(X, D)$ varies in the family
$\overline{\mathcal{M}}_{g,r}$, additional boundary strata
arise from curve degeneration. For each stable graph
$\Gamma$ with $\sum_v g(v) + h^1(\Gamma) = g$ and
$r$~marked legs, the clutching map
$\xi_\Gamma \colon \prod_v \overline{\mathcal{M}}_{g(v), n(v)} \to
\overline{\mathcal{M}}_{g,r}$ produces a boundary stratum
$D_\Gamma^{\mathrm{nod}}$. Over this stratum, the bordered FM
space acquires a fiber that is a product of bordered FM spaces on
the component curves:
\begin{equation}\label{eq:nodal-fiber}
\overline{C}{}^{\,\mathrm{oc}}_{k,\vec{m}}(X, D, \tau)
\big|_{D_\Gamma^{\mathrm{nod}}}
\;\cong\;
\prod_{v \in V(\Gamma)}
\overline{C}{}^{\,\mathrm{oc}}_{k_v, \vec{m}_v}
(X_v, D_v, \tau_v),
\end{equation}
where $(k_v, \vec{m}_v)$ is the restriction of interior and
boundary multiplicities to the component~$v$, and~$(X_v, D_v,
\tau_v)$ is the tangential log curve on the $v$-th irreducible
component with nodal points added to~$D_v$ and tangential
directions induced from the gluing parameter.

\medskip
The bordered FM compactification is:
\begin{equation}\label{eq:bordered-fm-result}
\overline{C}{}^{\,\mathrm{oc}}_{k,\vec{m}}(X, D, \tau)
\;:=\;
\text{iterated blowup of } Y
\text{ along }
\Delta_B^{\mathrm{bdy}},\;
\Delta_{S,B,j}^{\mathrm{mix}},\;
D_\Gamma^{\mathrm{nod}}
\end{equation}
in the order prescribed above.
\end{construction}

\begin{theorem}[Properties of the bordered FM compactification;
\ClaimStatusProvedHere]
\label{thm:bordered-fm-properties}
\index{bordered FM compactification!properties}
The bordered FM compactification
$\overline{C}{}^{\,\mathrm{oc}}_{k,\vec{m}}(X, D, \tau)$ of
Construction~\textup{\ref{constr:bordered-fm}} has the following
properties:
\begin{enumerate}[label=\textup{(\roman*)}]
\item \emph{Smooth manifold with corners.}
 $\overline{C}{}^{\,\mathrm{oc}}_{k,\vec{m}}$ is a compact smooth
 manifold with corners of dimension $2k + |\vec{m}|$.
\item \emph{Dense open embedding.}
 There is an open embedding
 $j \colon \operatorname{Conf}^{\mathrm{oc}}_{k,\vec{m}}(X, D, \tau)
 \hookrightarrow
 \overline{C}{}^{\,\mathrm{oc}}_{k,\vec{m}}(X, D, \tau)$
 with dense image.
\item \emph{Normal crossings.}
 The boundary $\partial\,\overline{C}{}^{\,\mathrm{oc}}_{k,\vec{m}}$
 is a union of smooth codimension-one faces intersecting in
 normal crossings.
\item \emph{Functoriality.}
 A morphism of tangential log curves $(X, D, \tau) \to
 (X', D', \tau')$ (i.e., a morphism $f \colon X \to X'$ with
 $f_* D \leq D'$ and compatible tangential data) induces a
 smooth map
 $\overline{C}{}^{\,\mathrm{oc}}_{k,\vec{m}}(X, D, \tau) \to
 \overline{C}{}^{\,\mathrm{oc}}_{k,\vec{m}}(X', D', \tau')$
 compatible with boundary stratification.
\end{enumerate}
\end{theorem}

\begin{proof}
\emph{(i)} Each blowup in Steps~2--5 has smooth center
(the diagonals $\Delta_B^{\mathrm{bdy}}$ are smooth because the
ordering prevents non-transverse intersections; the mixed loci
$\Delta_{S,B,j}^{\mathrm{mix}}$ are smooth because they are
transverse intersections of interior FM divisors and boundary
diagonals). The ordering by decreasing codimension ensures that
at each stage the center is a smooth submanifold of the current
blowup. The corners arise from the boundary factor
$\prod_j \overline{I}_{p_j}^{\,m_j}$.

\emph{(ii)} The interior
$\operatorname{Conf}^{\mathrm{oc}}_{k,\vec{m}}$ is the complement
of all collision loci, which embeds as a dense open subset of the
blowup by the universal property of blowups.

\emph{(iii)} Normal crossings. The interior FM factor has
normal-crossings boundary by Theorem~\ref{thm:normal-crossings}.
The boundary factor is a product of associahedra, whose
face structure is that of coordinate hyperplanes (Stasheff).
It remains to verify that the mixed blowups (Step~4) preserve
normal crossings. At each step, the center of the blowup is a
smooth submanifold $Z$ that is the intersection of the
interior-collision locus $\epsilon_{i,j} = 0$ with the
boundary-proximity locus $\rho_{i,j} = 0$. Since $\epsilon_{i,j}$
and $\rho_{i,j}$ are independent local coordinates (the former is
an interior distance, the latter an interior-to-boundary distance),
their joint vanishing locus $Z$ is a smooth subvariety of
codimension~$2$ in $\overline{C}_k(X) \times K_{\vec{m}}$.
The blowup $\operatorname{Bl}_Z$ replaces $Z$ by its projectivized
normal bundle $\bP(N_{Z/M})$; the exceptional divisor $E$ is
transverse to the pre-existing boundary divisors because $Z$
meets them transversally (the defining functions
$\{\epsilon_S,\, \delta_B,\, t_e\}$ of the existing boundary are
functionally independent from $\{\epsilon_{i,j},\, \rho_{i,j}\}$
at every point of~$Z$). Iterating over all mixed pairs $(i,j)$
in any order produces the same result: the centers are pairwise
disjoint (or nested) at each step, so the blowups commute
(by the clean-intersection criterion: disjoint or nested centers
commute under blowup).

\emph{(iv)} Functoriality follows from the naturality of the
blowup construction: if $f$ preserves diagonals and puncture
loci, it lifts through each iterated blowup.
\end{proof}

\begin{lemma}[Nested blowup commutativity; \ClaimStatusProvedElsewhere]
\label{lem:nested-blowup-commutativity}
\index{blowup!commutativity of nested}
Let $M$ be a smooth variety with smooth subvarieties
$Z_1, Z_2 \subset M$. If $Z_1$ and $Z_2$ are either
disjoint or nested \textup{(}$Z_1 \subseteq Z_2$ or
$Z_2 \subseteq Z_1$\textup{)}, then
$\operatorname{Bl}_{Z_2}(\operatorname{Bl}_{Z_1}(M))
\cong
\operatorname{Bl}_{Z_1}(\operatorname{Bl}_{Z_2}(M))$.
More generally, if $Z_1, \ldots, Z_r$ are pairwise nested
or disjoint, the iterated blowup
$\operatorname{Bl}_{Z_r} \cdots \operatorname{Bl}_{Z_1}(M)$
is independent of the ordering~\cite{FM94}.
\end{lemma}

\subsection{Codimension-one boundary decomposition}
\label{subsec:boundary-decomposition}
\index{bordered FM compactification!boundary strata|textbf}

\begin{proposition}[Four-type boundary decomposition;
\ClaimStatusProvedHere]
\label{prop:four-type-boundary}
\index{boundary strata!four types}
The codimension-one boundary of the bordered FM
compactification
$\overline{C}{}^{\,\mathrm{oc}}_{k,\vec{m}}(X, D, \tau)$
decomposes into four types of smooth faces:
\begin{equation}\label{eq:boundary-four-types}
\partial\,\overline{C}{}^{\,\mathrm{oc}}_{k,\vec{m}}
\;=\;
\partial_{\mathrm{I}}\,\overline{C}{}^{\,\mathrm{oc}}_{k,\vec{m}}
\;\cup\;
\partial_{\mathrm{II}}\,\overline{C}{}^{\,\mathrm{oc}}_{k,\vec{m}}
\;\cup\;
\partial_{\mathrm{III}}\,\overline{C}{}^{\,\mathrm{oc}}_{k,\vec{m}}
\;\cup\;
\partial_{\mathrm{IV}}\,\overline{C}{}^{\,\mathrm{oc}}_{k,\vec{m}}.
\end{equation}

\medskip
\noindent\textbf{Type~I: Interior collision strata.}
\index{boundary strata!Type I (interior collision)}
For each subset $S \subseteq \{1, \ldots, k\}$ with
$|S| \geq 2$, the \emph{interior collision face} is
\begin{equation}\label{eq:typeI-stratum}
\partial_S^{\mathrm{I}}\,
\overline{C}{}^{\,\mathrm{oc}}_{k,\vec{m}}
\;\cong\;
\overline{C}{}^{\,\mathrm{oc}}_{k-|S|+1,\, \vec{m}}
(X, D, \tau)
\;\times\;
\overline{C}_{|S|}(\mathbb{A}^1).
\end{equation}
This records the collision of interior points indexed by~$S$
at a point $q \in X \setminus D$. The first factor records the
collision center~$q$ together with the remaining $k - |S|$
interior points and all $|\vec{m}|$ boundary points; the second
factor is the genus-$0$ FM screen recording relative positions on
the tangent space $T_q X \cong \mathbb{A}^1$.

\emph{Geometric fiber.}
The fiber over a point of the first factor is
$\overline{C}_{|S|}(\mathbb{A}^1)$, the FM compactification of
$|S|$ points in the affine line. This is a smooth manifold of
dimension $2(|S| - 1) - 1 = 2|S| - 3$ (after quotienting by the
translation--dilation group).

\emph{Normal bundle.}
The exceptional divisor $D_S^{\mathrm{I}}$ is the projectivization
of the normal bundle to $\Delta_S^{\mathrm{int}}$ inside the FM
space. The normal bundle is
$N_{\Delta_S/\operatorname{FM}_k} \cong
T_q X \otimes (\mathbb{C}^{|S|}/\mathbb{C})
\cong T_q X^{\oplus(|S|-1)}$,
and the exceptional divisor $D_S^{\mathrm{I}}$ is its
projectivization: a $\mathbb{P}^{|S|-2}$-bundle over
the collision center locus.
As a boundary-defining function, the collision scale
$\epsilon_S = \max_{i \in S} |z_i - q|$ (in local coordinates
near~$q$) satisfies $D_S^{\mathrm{I}} = \{\epsilon_S = 0\}$.

\emph{Residue map.}
The residue along $D_S^{\mathrm{I}}$ is:
\begin{equation}\label{eq:typeI-residue}
\operatorname{Res}_{D_S^{\mathrm{I}}}
\;\colon\;
\Omega^p_{\log}\bigl(
\overline{C}{}^{\,\mathrm{oc}}_{k,\vec{m}}\bigr)
\;\longrightarrow\;
\Omega^{p-1}_{\log}\bigl(
\partial_S^{\mathrm{I}}\,
\overline{C}{}^{\,\mathrm{oc}}_{k,\vec{m}}\bigr).
\end{equation}
Concretely, for a logarithmic form
$\omega = \frac{d\epsilon_S}{\epsilon_S} \wedge \alpha + \beta$:
$\operatorname{Res}_{D_S^{\mathrm{I}}}(\omega) = \alpha|_{\epsilon_S=0}$.
This extracts the OPE coefficient of the $\cA$-algebra product
$\mu_S \colon \cA^{\boxtimes |S|} \to \cA$ from the leading
collision singularity.

\medskip
\noindent\textbf{Type~II: Consecutive boundary collision strata.}
\index{boundary strata!Type II (boundary collision)}
For each boundary component~$j$ and each consecutive block
$B = \{a, a{+}1, \ldots, b\} \subseteq \{1, \ldots, m_j\}$
with $|B| \geq 2$, the
\emph{consecutive boundary collision face} is
\begin{equation}\label{eq:typeII-stratum}
\partial_{B,j}^{\mathrm{II}}\,
\overline{C}{}^{\,\mathrm{oc}}_{k,\vec{m}}
\;\cong\;
\overline{C}{}^{\,\mathrm{oc}}_{k,\, \vec{m}'}(X, D, \tau)
\;\times\;
K_{|B|},
\end{equation}
where $\vec{m}' = (m_1, \ldots, m_j - |B| + 1, \ldots, m_r)$
(the colliding block is replaced by a single point) and $K_{|B|}$
is the $|B|$-th Stasheff associahedron parametrizing the relative
positions of $|B|$ ordered boundary points approaching a common
limit. Consecutiveness is essential: the ordered configuration
$(t_a < t_{a+1} < \cdots < t_b)$ can only degenerate by
shrinking the gaps $t_{l+1} - t_l \to 0$ for $l \in B$, which
preserves the ordering.

\emph{Geometric fiber.}
The fiber over a point of the first factor is the Stasheff
associahedron $K_{|B|}$, a convex polytope of dimension
$|B| - 2$.

\emph{Normal bundle.}
The boundary-defining function is the total gap
$\delta_B = t_b^{(j)} - t_a^{(j)} \geq 0$, with
$\partial_{B,j}^{\mathrm{II}} = \{\delta_B = 0\}$.
The normal bundle is trivial of rank~$1$ (the real line
generated by $\partial/\partial\delta_B$).

\emph{Residue map.}
The residue along $\partial_{B,j}^{\mathrm{II}}$ extracts the
boundary $A_\infty$-composition:
\begin{equation}\label{eq:typeII-residue}
\operatorname{Res}_{\partial_{B,j}^{\mathrm{II}}}
\;\colon\;
\omega \;\longmapsto\;
\omega|_{\delta_B = 0}
\;\leadsto\;
\mu_B^{\mathrm{bdy}}
\colon \mathcal{M}_j^{\otimes |B|} \to \mathcal{M}_j,
\end{equation}
the higher multiplication on the boundary module
$\mathcal{M}_j$ supported at~$p_j$.

\medskip
\noindent\textbf{Type~III: Mixed bubbling strata.}
\index{boundary strata!Type III (mixed bubbling)}
For each interior point~$i$ approaching boundary component~$j$,
and each consecutive boundary block $B \subseteq
\{1, \ldots, m_j\}$ (possibly empty) colliding simultaneously,
the \emph{mixed bubbling face} is
\begin{equation}\label{eq:typeIII-stratum}
\partial_{i,B,j}^{\mathrm{III}}\,
\overline{C}{}^{\,\mathrm{oc}}_{k,\vec{m}}
\;\cong\;
\overline{C}{}^{\,\mathrm{oc}}_{k-1,\, \vec{m}'}(X, D, \tau)
\;\times\;
\overline{C}{}^{\,\mathrm{oc}}_{1,\,
 (0,\ldots,|B|,\ldots,0)}
(\mathbb{H}, \{0\}, \tau_j),
\end{equation}
where $\mathbb{H} = \{z \in \mathbb{C} : \operatorname{Im}(z) \geq 0\}$
is the closed upper half-plane (the local model for the real
oriented blowup near~$p_j$),
$\vec{m}' = (m_1, \ldots, m_j - |B|, \ldots, m_r)$,
and the second factor is the bordered FM space of $1$~interior
point and $|B|$~boundary points on the half-plane.

Geometrically, as $z_i \to p_j$ along a direction $\theta_i$,
the interior point bubbles off onto the tangent half-plane
$T_{p_j}^+ X := \{v \in T_{p_j} X :
\operatorname{arg}(v/\tau_j) \in (0, 2\pi)\}$; boundary points
in~$B$ that simultaneously approach $\star_j$ join this bubble.
The mixed bubble carries both a bulk insertion (from the
interior point) and boundary insertions (from the collapsing
boundary points), recording the $\cA$-module action
$\rho_j \colon \cA \otimes \mathcal{M}_j^{\otimes |B|} \to
\mathcal{M}_j$.

\emph{Geometric fiber.}
The fiber is the bordered FM compactification of one interior
point and $|B|$ boundary points on the upper half-plane: the
\emph{Voronov Swiss-cheese} space
$\mathrm{SC}_{1, |B|} :=
\overline{C}{}^{\,\mathrm{oc}}_{1, (|B|)}(\mathbb{H}, \{0\}, 1)$,
which is a manifold with corners of dimension $|B| + 1$.

\emph{Normal bundle.}
The boundary-defining function is the distance
$\rho_{i,j} = |z_i - p_j|$ in the real oriented blowup, with
$\partial_{i,B,j}^{\mathrm{III}} = \{\rho_{i,j} = 0\}$. The
conormal is spanned by $d\rho_{i,j}$; the normal bundle is
trivial of rank~$1$.

\emph{Residue map.}
The residue extracts the bulk-to-boundary OPE:
\begin{equation}\label{eq:typeIII-residue}
\operatorname{Res}_{\partial_{i,B,j}^{\mathrm{III}}}
\;\colon\;
\omega \;\longmapsto\;
\omega|_{\rho_{i,j} = 0}
\;\leadsto\;
\rho_j^B
\colon \cA \otimes \mathcal{M}_j^{\otimes |B|} \to
\mathcal{M}_j,
\end{equation}
extracting the open-closed intertwining map. When $B =
\emptyset$, this is the pure bulk-to-boundary operation
$\rho_j \colon \cA \to \operatorname{End}(\mathcal{M}_j)$.

\medskip
\noindent\textbf{Type~IV: Nodal clutching strata.}
\index{boundary strata!Type IV (nodal clutching)}
These arise from stable-curve degeneration, present only when
the base curve $(X, D)$ varies in moduli. For each stable
graph~$\Gamma$ with one edge~$e$ (a codimension-one
degeneration), the clutching face is:
\begin{equation}\label{eq:typeIV-stratum}
\partial_\Gamma^{\mathrm{IV}}\,
\overline{C}{}^{\,\mathrm{oc}}_{k,\vec{m}}
\;\cong\;
\prod_{v \in V(\Gamma)}
\overline{C}{}^{\,\mathrm{oc}}_{k_v,\vec{m}_v}
(X_v, D_v, \tau_v),
\end{equation}
where the product is over the vertices of~$\Gamma$, with the
distribution of interior and boundary multiplicities determined
by the edge data of~$\Gamma$. The edge~$e$ corresponds to a
node: a point $q$ on the degenerate curve where two irreducible
components are identified. The gluing parameter
$t_e$ with $|t_e| < \epsilon$ is the transverse coordinate,
with $t_e = 0$ defining the nodal locus.

The clutching strata separate into two sub-types, following
the modular operad structure
(Definition~\ref{def:stable-graph},
Theorem~\ref{thm:boundary-higher-genus}):
\begin{enumerate}[label=\textup{(\alph*)}]
\item \emph{Separating clutching}
 ($\Gamma$ disconnects upon removing~$e$):
 \begin{equation}\label{eq:separating-clutching}
 \partial_\Gamma^{\mathrm{IV,sep}}
 \;\cong\;
 \overline{C}{}^{\,\mathrm{oc}}_{k_1, \vec{m}_1}
 (X_1, D_1, \tau_1)
 \;\times\;
 \overline{C}{}^{\,\mathrm{oc}}_{k_2, \vec{m}_2}
 (X_2, D_2, \tau_2),
 \end{equation}
 with $g_1 + g_2 = g$, $k_1 + k_2 = k$,
 $\vec{m}_1 + \vec{m}_2 = \vec{m}$ (distributed between the
 two components), and $D_j$ augmented by the nodal point
 on component~$j$.
\item \emph{Non-separating clutching}
 ($\Gamma$ remains connected after removing~$e$):
 \begin{equation}\label{eq:nonseparating-clutching}
 \partial_\Gamma^{\mathrm{IV,ns}}
 \;\cong\;
 \overline{C}{}^{\,\mathrm{oc}}_{k, \vec{m}'}
 (X', D', \tau'),
 \end{equation}
 where $(X', D', \tau')$ is the partial normalization at the
 node, with genus $g - 1$ and $D' = D + q^+ + q^-$ (two new
 punctures at the preimages of the node), and $\vec{m}'$ is
 $\vec{m}$ augmented by $m_{q^+} = m_{q^-} = 0$.
\end{enumerate}

\emph{Normal bundle.}
The boundary-defining function for the clutching face is the
modulus $|t_e|$ of the gluing parameter, with
$\partial_\Gamma^{\mathrm{IV}} = \{t_e = 0\}$.

\emph{Residue map.}
The residue along $\partial_\Gamma^{\mathrm{IV}}$ produces the
clutching operation:
\begin{equation}\label{eq:typeIV-residue}
\operatorname{Res}_{\partial_\Gamma^{\mathrm{IV}}}
\;\colon\;
\omega \;\longmapsto\;
\xi_\Gamma^*(\omega|_{t_e = 0})
\end{equation}
where $\xi_\Gamma$ is the clutching map of the modular operad.
For separating clutching, this contracts a pair of legs
using the invariant pairing $\langle -, - \rangle_\cA$:
\[
\xi_\Gamma^*\colon
\operatorname{End}_\cA(n_1) \otimes \operatorname{End}_\cA(n_2)
\;\xrightarrow{\;\operatorname{id} \otimes \langle -,-\rangle
\otimes \operatorname{id}\;}
\operatorname{End}_\cA(n_1 + n_2 - 2).
\]
For non-separating clutching, this traces a pair of inputs
against the pairing:
$\Delta_{\mathrm{ns}} \colon
\operatorname{End}_\cA(n+2) \to \operatorname{End}_\cA(n)$,
$\Delta_{\mathrm{ns}}(\phi)(a_1, \ldots, a_n)
= \sum_i \phi(a_1, \ldots, a_n, e_i, e^i)$.
\end{proposition}

\begin{proof}
It suffices to verify that every codimension-one face of
$\overline{C}{}^{\,\mathrm{oc}}_{k,\vec{m}}$ is of one of the
four types. A codimension-one face is defined by one
boundary-defining function reaching zero. The possible
degenerations are:
\begin{enumerate}
\item $\epsilon_S = 0$ for some $S \subseteq \{1, \ldots, k\}$
 with $|S| \geq 2$: interior collision (Type~I).
\item $\delta_B = 0$ for some consecutive block~$B$ on
 boundary~$j$ with $|B| \geq 2$: consecutive boundary collision
 (Type~II).
\item $\rho_{i,j} = 0$ for some interior point~$i$ and
 boundary component~$j$: mixed bubbling (Type~III). The
 consecutive block~$B$ of simultaneously collapsing boundary
 points is determined by the limiting configuration.
\item $t_e = 0$ for some edge~$e$ of a stable graph: nodal
 clutching (Type~IV).
\end{enumerate}
Any other degeneration has codimension $\geq 2$ (e.g.,
simultaneous collisions of two independent clusters, or a
non-consecutive boundary collision). The four types are
mutually exclusive because their boundary-defining functions
($\epsilon_S$, $\delta_B$, $\rho_{i,j}$, $t_e$) are
independent local coordinates. Exhaustiveness follows from
the compactness of
$\overline{C}{}^{\,\mathrm{oc}}_{k,\vec{m}}$: every sequence
in the interior has a convergent subsequence, and the limit
is classified by the minimal degeneration pattern, which is
one of the four types.
\end{proof}

\begin{remark}[Comparison with classical FM]
\label{rem:comparison-classical-fm}
When $D = \emptyset$ (no punctures), the bordered FM
compactification reduces to the classical FM compactification:
$\overline{C}{}^{\,\mathrm{oc}}_{k,\vec{0}}(X, \emptyset, -)
= \overline{C}_k(X)$.
Only Type~I and Type~IV faces are present. The Type~II and
Type~III faces are \emph{new} in the bordered setting and encode
the open-string sector. When the curve is fixed
(no moduli variation), Type~IV faces are absent and the boundary
has three types, matching the Swiss-cheese operad structure
of Voronov~\cite{Voronov99}.
\end{remark}


\section{Modular Maurer--Cartan equation with clutching}
\label{sec:modular-mc-clutching}
\index{modular Maurer--Cartan equation!with clutching|textbf}
\index{clutching!in modular MC equation|textbf}

The bordered FM compactification of
\S\ref{sec:bordered-fm} carries, via the Stokes theorem on
its boundary, the \emph{modular Maurer--Cartan equation with
clutching}: the master equation
$d\Theta + \tfrac{1}{2}[\Theta, \Theta]
+ \hbar\Delta_{\mathrm{clutch}}(\Theta) = 0$
governing the genus expansion of chiral correlators on
punctured stable curves. The modular open-closed convolution
algebra carries the equation; the MC element~$\Theta_\cA$ and the
equation itself follow from Stokes' theorem on the four-type
boundary of~$\overline{C}{}^{\,\mathrm{oc}}_{k,\vec{m}}$.

\subsection{The modular open-closed convolution algebra}
\label{subsec:modular-oc-convolution}
\index{modular convolution algebra!open-closed|textbf}

\begin{definition}[Open-closed modular convolution algebra]
\label{def:oc-modular-convolution}
Let $\cA$ be a cyclic chiral algebra on a smooth projective
curve~$X$, with $\cA$-modules
$\mathcal{M}_1, \ldots, \mathcal{M}_r$ supported at punctures
$D = p_1 + \cdots + p_r$, and let
$(X, D, \tau)$ be a tangential log curve. The \emph{open-closed
modular convolution algebra} is
\begin{equation}\label{eq:oc-convolution}
\mathfrak{g}^{\mathrm{oc}}_{\cA, \vec{\mathcal{M}}}
\;:=\;
\prod_{\substack{g, k, \vec{m} \\ 2g - 2 + k + |\vec{m}| > 0}}
\operatorname{Hom}\!\bigl(
C_*\bigl(
\overline{C}{}^{\,\mathrm{oc}}_{k,\vec{m}}(X, D, \tau)
\bigr),\;
\operatorname{End}^{\mathrm{oc}}_{\cA, \vec{\mathcal{M}}}
(k, \vec{m})
\bigr),
\end{equation}
where the endomorphism module is
\begin{equation}\label{eq:oc-endomorphism}
\operatorname{End}^{\mathrm{oc}}_{\cA, \vec{\mathcal{M}}}
(k, \vec{m})
\;:=\;
\operatorname{Hom}\!\bigl(
\cA^{\otimes k}
\otimes
\bigotimes_{j=1}^{r} \mathcal{M}_j^{\otimes m_j},\;
\cA \otimes \bigotimes_{j=1}^r \mathcal{M}_j
\bigr),
\end{equation}
equipped with the internal differential from~$\cA$ and
the~$\mathcal{M}_j$. The completed product is taken with respect
to the \emph{genus--degree weight}
$w(g, k, \vec{m}) = 2g - 2 + k + |\vec{m}|$
and the genus grading $\hbar^g$.
\end{definition}

\begin{theorem}[dg~Lie structure on the open-closed convolution
algebra; \ClaimStatusConditional]
\label{thm:oc-convolution-dg-lie}
\index{modular convolution algebra!open-closed!dg Lie structure}
The open-closed convolution algebra
$\mathfrak{g}^{\mathrm{oc}}_{\cA, \vec{\mathcal{M}}}$ carries
a dg~Lie structure with:

\noindent\textbf{Differential.}
The differential $D^{\mathrm{oc}}$ has \emph{seven} components,
extending the five-component differential of
Theorem~\textup{\ref{thm:convolution-dg-lie-structure}}:
\begin{equation}\label{eq:Doc-seven}
D^{\mathrm{oc}}
\;=\;
\underbrace{d_{\mathrm{int}}
+ [\tau,-]
+ d_{\mathrm{sew}}
+ d_{\mathrm{pf}}
+ \hbar\Delta}_{\text{closed sector}}
\;+\;
\underbrace{d_{\mathrm{bdy}}
+ d_{\mathrm{mix}}}_{\text{open sector}}.
\end{equation}
The five closed-sector components are exactly those of
equation~\eqref{eq:D-five-components}; the two new open-sector
components are:
\begin{enumerate}[label=\textup{(\alph*)}]
\item $d_{\mathrm{bdy}}$: the \emph{boundary differential},
 transporting the boundary $\partial$ on the associahedral factor
 $\prod_j K_{m_j}$ through the Hom functor. Geometrically, this
 sums over Type~II faces, extracting the $A_\infty$-composition
 on each boundary module~$\mathcal{M}_j$:
 \begin{equation}\label{eq:d-bdy}
 d_{\mathrm{bdy}}(\alpha)_{g, k, \vec{m}}
 \;=\;
 \sum_{j=1}^r \;
 \sum_{\substack{B \subseteq \{1, \ldots, m_j\} \\
 B \text{ consecutive,}\; |B| \geq 2}}
 \alpha_{g, k, \vec{m}'} \circ_{B,j}
 \mu_{|B|}^{\mathcal{M}_j},
 \end{equation}
 where $\mu_{|B|}^{\mathcal{M}_j}$ is the $|B|$-ary
 multiplication on~$\mathcal{M}_j$ (the $A_\infty$-structure
 map), composed at the location of block~$B$ on boundary~$j$.
\item $d_{\mathrm{mix}}$: the \emph{mixed differential},
 transporting the boundary $\partial$ on the mixed bubbling faces
 through the Hom functor. Geometrically, this sums over
 Type~III faces:
 \begin{equation}\label{eq:d-mix}
 d_{\mathrm{mix}}(\alpha)_{g, k, \vec{m}}
 \;=\;
 \sum_{i=1}^{k}\;
 \sum_{j=1}^{r}\;
 \sum_{\substack{B \subseteq \{1,\ldots,m_j\} \\
 B \text{ consecutive}}}
 \alpha_{g, k-1, \vec{m}'} \circ_{i,B,j}
 \rho_j^B,
 \end{equation}
 where $\rho_j^B \colon \cA \otimes \mathcal{M}_j^{\otimes |B|}
 \to \mathcal{M}_j$ is the open-closed intertwining map, composed
 at interior slot~$i$ and boundary block~$B$ on boundary~$j$.
\end{enumerate}

\noindent\textbf{Lie bracket.}
The bracket $[-,-]^{\mathrm{oc}}$ extends the convolution bracket of
Construction~\textup{\ref{const:explicit-convolution-bracket}} by
allowing both interior and boundary legs to be composed via the
invariant pairing. Explicitly, for
$\alpha \in \mathfrak{g}^{\mathrm{oc}}(g_1, k_1, \vec{m}_1)$ and
$\beta \in \mathfrak{g}^{\mathrm{oc}}(g_2, k_2, \vec{m}_2)$:
\begin{equation}\label{eq:oc-bracket}
[\alpha, \beta]^{\mathrm{oc}}
\;=\;
\sum_{\substack{i \leq k_1,\, j \leq k_2}}
(\alpha \circ_{i,j}^{\mathrm{int}} \beta
- (-1)^{|\alpha||\beta|}\,
\beta \circ_{j,i}^{\mathrm{int}} \alpha),
\end{equation}
where $\circ_{i,j}^{\mathrm{int}}$ is the single-edge composition
that joins the $i$-th interior leg of~$\alpha$ to the $j$-th
interior leg of~$\beta$ by contracting with the bulk pairing
$\langle -, - \rangle_\cA$ and pushing forward along the
separating clutching map. The boundary legs are
\emph{spectators} in the bracket: they are carried through
unchanged.

\noindent\textbf{Square-zero.}
$(D^{\mathrm{oc}})^2 = 0$.
\end{theorem}

\begin{proof}
The seven-component differential $D^{\mathrm{oc}}$ is the
transport of the total boundary operator~$\partial$ on
$C_*(\overline{C}{}^{\,\mathrm{oc}}_{k,\vec{m}})$ through
the Hom functor, together with the internal differentials.
Since $\partial^2 = 0$ on the chain complex of a manifold with
corners (by the same boundary-of-boundary cancellation as in
Theorem~\ref{thm:convolution-d-squared-zero}), and the
internal differentials satisfy $d_{\mathrm{int}}^2 = 0$,
the total differential satisfies
$(D^{\mathrm{oc}})^2 = 0$.

The component matching between $D^{\mathrm{oc}}$ and the
four boundary types is:
\begin{center}
\renewcommand{\arraystretch}{1.2}
\begin{tabular}{lll}
\emph{Boundary type} & \emph{Differential component} &
 \emph{Algebraic operation} \\
\hline
Type~I (interior collision)
 & $d_{\mathrm{pf}} + [\tau, -]$
 & Bulk OPE / tree-level composition \\
Type~II (boundary collision)
 & $d_{\mathrm{bdy}}$
 & $A_\infty$-multiplication on $\mathcal{M}_j$ \\
Type~III (mixed bubbling)
 & $d_{\mathrm{mix}}$
 & Bulk-to-boundary OPE \\
Type~IV(a) (separating node)
 & $d_{\mathrm{sew}}$
 & Separating clutching \\
Type~IV(b) (non-separating node)
 & $\hbar\Delta$
 & BV operator / trace \\
Internal
 & $d_{\mathrm{int}}$
 & Algebra \& module differentials
\end{tabular}
\end{center}

The identity $(D^{\mathrm{oc}})^2 = 0$ expanded into
cross-terms gives, in addition to the closed-sector identities of
Theorem~\ref{thm:ambient-d-squared-zero}, the new open-sector
relations:
\begin{align}
d_{\mathrm{bdy}}^2 &= 0
 &&\text{($A_\infty$ relations on $\mathcal{M}_j$)},
 \label{eq:dbdy-squared} \\
d_{\mathrm{mix}}^2 &= 0
 &&\text{(iterated bulk-to-boundary is consistent)},
 \label{eq:dmix-squared} \\
[d_{\mathrm{bdy}}, d_{\mathrm{mix}}] &= 0
 &&\text{(boundary composition commutes with mixed bubbling)},
 \label{eq:bdy-mix-commute} \\
[d_{\mathrm{bdy}}, d_{\mathrm{sew}}]
+ [d_{\mathrm{mix}}, d_{\mathrm{pf}}] &= 0
 &&\text{(open-closed codimension-2 face cancellation)}.
 \label{eq:oc-codim2}
\end{align}
Each identity follows from the boundary-of-boundary vanishing
$\partial^2 = 0$ on the corresponding codimension-$2$ stratum
of $\overline{C}{}^{\,\mathrm{oc}}_{k,\vec{m}}$:
\eqref{eq:dbdy-squared} from the codimension-$2$ associahedral
faces (nested consecutive blocks);
\eqref{eq:dmix-squared} from two interior points independently
approaching the same boundary;
\eqref{eq:bdy-mix-commute} from a boundary collision simultaneous
with a mixed bubble;
\eqref{eq:oc-codim2} from a boundary collision simultaneous with
a nodal degeneration.

For~\eqref{eq:bdy-mix-commute}, consider a
codimension-$2$ stratum~$S$ where an interior point~$z_i$
approaches puncture~$p_j$ (mixed) while simultaneously two
boundary points $t_a, t_{a+1}$ on interval~$\mathcal{I}_j$
collide (boundary). The stratum~$S$ is an intersection
$S = D_{\mathrm{mix}} \cap D_{\mathrm{bdy}}$, where
$D_{\mathrm{mix}} = \{\rho_{i,j} = 0\}$ and
$D_{\mathrm{bdy}} = \{\delta_{\{a,a+1\}} = 0\}$.
By the normal-crossings property
(Theorem~\ref{thm:bordered-fm-properties}(iii)),
these are transverse:
$\rho_{i,j}$ is a distance in the interior-to-boundary direction
and $\delta_{\{a,a+1\}}$ is a distance along the boundary interval,
so $d\rho_{i,j} \wedge d\delta_{\{a,a+1\}} \neq 0$.
The stratum~$S$ appears as a face of both
$D_{\mathrm{mix}}$ (where $\partial D_{\mathrm{mix}}$ includes the
sub-stratum with a simultaneous boundary collision) and
$D_{\mathrm{bdy}}$ (where $\partial D_{\mathrm{bdy}}$ includes the
sub-stratum with a simultaneous mixed bubble). The
$\partial^2 = 0$ identity applied to the cell dual to~$S$ gives
$d_{\mathrm{bdy}} \circ d_{\mathrm{mix}}
+ d_{\mathrm{mix}} \circ d_{\mathrm{bdy}} = 0$,
which is~\eqref{eq:bdy-mix-commute}.

The bracket $[-,-]^{\mathrm{oc}}$ satisfies the Jacobi identity
by the same pre-Lie argument as
Remark~\ref{rem:convolution-pre-lie-jacobi}: the
single-edge composition $\circ_{i,j}^{\mathrm{int}}$ is pre-Lie,
and antisymmetrization of a pre-Lie product is Lie.
\end{proof}

\subsection{The modular MC element with clutching}
\label{subsec:modular-mc-element-clutching}
\index{modular Maurer--Cartan element!with clutching|textbf}

\begin{definition}[Modular twisting morphism]
\label{def:modular-twisting-morphism-oc}
\index{modular twisting morphism|textbf}
Let $\cA$ be a modular Koszul chiral algebra on $(X, D, \tau)$
with the genus-completed bar differential
$D_\cA = \sum_{g \geq 0} \hbar^g\, d_\cA^{(g)}$ of
Theorem~\ref{thm:bar-modular-operad}. The
\emph{modular twisting morphism} is
\begin{equation}\label{eq:theta-oc}
\Theta_\cA
\;:=\;
D_\cA^{\mathrm{oc}} - \dzero
\;=\;
\sum_{g \geq 1} \hbar^g\, d_\cA^{(g)}
\;+\;
d_{\mathrm{bdy}}^{(0)}
\;+\;
d_{\mathrm{mix}}^{(0)}
\;\in\;
\mathfrak{g}^{\mathrm{oc}}_{\cA, \vec{\mathcal{M}}},
\end{equation}
where $D_\cA^{\mathrm{oc}}$ is the genus-completed open-closed
bar differential (the coderivation encoding the full
$\mathrm{FCom}^{\mathrm{oc}}$-algebra structure on the
open-closed bar complex), $\dzero$ is the genus-$0$ closed-sector
bar differential (the Arnold-relation part), and the three
summands separate the positive-genus closed corrections from the
genus-$0$ boundary and mixed terms.

More precisely, $\Theta_\cA$ is the element of
$\mathfrak{g}^{\mathrm{oc}}_{\cA, \vec{\mathcal{M}}}$ whose
$(g, k, \vec{m})$-component is the composition
\begin{equation}\label{eq:theta-components}
\Theta_\cA^{(g, k, \vec{m})}
\;=\;
\int_{\overline{C}{}^{\,\mathrm{oc}}_{k,\vec{m}}}
\omega^{(g)}_{k, \vec{m}}(\cA, \vec{\mathcal{M}}),
\end{equation}
where $\omega^{(g)}_{k, \vec{m}}$ is the chiral amplitude form:
the logarithmic top-form on
$\overline{C}{}^{\,\mathrm{oc}}_{k,\vec{m}}$ with
coefficients in
$\operatorname{End}^{\mathrm{oc}}_{\cA, \vec{\mathcal{M}}}$,
built from the propagator $\eta^{(g)}_{ij}$ (for interior pairs),
the boundary propagator $\eta^{(g)}_{i, t_a^{(j)}}$ (for
interior-boundary pairs), and the boundary-boundary propagator
$\eta^{\mathrm{bdy}}_{a,b}$ (for consecutive boundary pairs on
the same interval).
\end{definition}

\begin{theorem}[Modular MC equation with clutching;
\ClaimStatusConditional]
\label{thm:modular-mc-clutching}
\index{modular Maurer--Cartan equation!with clutching}
The modular twisting morphism $\Theta_\cA$ of
Definition~\textup{\ref{def:modular-twisting-morphism-oc}}
satisfies the \emph{modular Maurer--Cartan equation with
clutching}:
\begin{equation}\label{eq:modular-mc-clutching}
\boxed{
d\,\Theta_\cA
+ \tfrac{1}{2}[\Theta_\cA, \Theta_\cA]^{\mathrm{oc}}
+ \hbar\,\Delta_{\mathrm{clutch}}(\Theta_\cA)
= 0
}
\end{equation}
where:
\begin{enumerate}[label=\textup{(\roman*)}]
\item $d\,\Theta_\cA := (d_{\mathrm{int}} + d_{\mathrm{bdy}}
 + d_{\mathrm{mix}})\,\Theta_\cA$ is the full open-closed
 chain differential applied to~$\Theta_\cA$, where
 $d_{\mathrm{int}}\,\Theta_\cA = [\dzero, \Theta_\cA]$ is the
 interior (genus-$0$ collision) component,
 $d_{\mathrm{bdy}}$ is the boundary associahedral component, and
 $d_{\mathrm{mix}}$ is the mixed (bulk-to-boundary) component;
\item $[\Theta_\cA, \Theta_\cA]^{\mathrm{oc}}$ is the
 bracket in the open-closed convolution Lie algebra
 \textup{(}Theorem~\textup{\ref{thm:oc-convolution-dg-lie}},
 equation~\eqref{eq:oc-bracket}\textup{)};
\item $\Delta_{\mathrm{clutch}}(\Theta_\cA)$ is the
 non-separating clutching operator applied to $\Theta_\cA$:
 \begin{equation}\label{eq:delta-clutch-def}
 \Delta_{\mathrm{clutch}}(\Theta_\cA)_{g, k, \vec{m}}
 \;:=\;
 \sum_a
 \Theta_\cA^{(g-1, k+2, \vec{m})}
 (\underbrace{x_1, \ldots, x_k}_{\text{original}},\, e_a,\, e^a;\;
 \vec{t}),
 \end{equation}
 where $\{e_a\}$, $\{e^a\}$ are dual bases for the invariant
 pairing $\langle -, - \rangle_\cA$, inserted in the
 $(k{+}1)$-st and $(k{+}2)$-nd interior slots and summed over
 the basis index~$a$.
 This contracts the two additional interior legs of the
 $(g{-}1, k{+}2, \vec{m})$-component by the pairing, producing a
 $(g, k, \vec{m})$-component with genus increased by~$1$.
\end{enumerate}
\end{theorem}

\begin{proof}
The proof is a Stokes-theorem argument on the bordered FM
compactification, following the same logic as
Theorem~\ref{thm:mc2-bar-intrinsic} but with the four-type
boundary decomposition of
Proposition~\ref{prop:four-type-boundary} in place of the
boundary of $\overline{\cM}_{g,n}$.

\medskip
\noindent\textsc{Setup.}
Consider a generic one-parameter family in the bordered FM
compactification: a smooth chain
$\sigma \in C_{2k + |\vec{m}|}
(\overline{C}{}^{\,\mathrm{oc}}_{k,\vec{m}})$. The
chiral amplitude form
$\omega^{(g)}_{k,\vec{m}}(\cA, \vec{\mathcal{M}})$ of
Definition~\ref{def:modular-twisting-morphism-oc} is a
logarithmic top-form on
$\overline{C}{}^{\,\mathrm{oc}}_{k,\vec{m}}$ with values in
$\operatorname{End}^{\mathrm{oc}}_{\cA, \vec{\mathcal{M}}}$.
The amplitude form $\omega^{(g)}_{k,\vec{m}}$ is
\emph{logarithmic}: its singularities along the boundary
divisors $D_I \subset
\overline{C}{}^{\,\mathrm{oc}}_{k,\vec{m}}$ are of the form
$f \cdot d\log(\epsilon_{I_1}) \wedge \cdots \wedge
d\log(\epsilon_{I_p})$ with $f$ smooth, hence $L^1$-integrable
(the $d\log$ singularities are integrable in one real dimension).
The normal-crossings property
(Theorem~\ref{thm:bordered-fm-properties}(iii)) ensures that
the boundary divisors meet transversally, so the multi-logarithmic
forms at corners remain integrable. Stokes' theorem for
manifolds with corners (see~\cite{FM94}, Appendix~B for
the classical FM case) then gives
\begin{equation}\label{eq:stokes-bordered}
\int_\sigma d\omega^{(g)}_{k,\vec{m}}
\;=\;
\int_{\partial\sigma} \omega^{(g)}_{k,\vec{m}}
\;=\;
\sum_{F \in \mathcal{F}(\partial\sigma)}
\int_F \omega^{(g)}_{k,\vec{m}},
\end{equation}
where $\mathcal{F}(\partial\sigma)$ is the set of
codimension-one faces of~$\partial\sigma$. By
Proposition~\ref{prop:four-type-boundary}, these faces
decompose into four types, and each type contributes a
specific term to the MC equation.

\medskip
\noindent\textsc{Type~I contribution:
chain differential $d\,\Theta_\cA$.}

The integral over Type~I faces produces the collision
residues of the genus-$0$ bar differential:
\[
\sum_{S}\,
\int_{\partial_S^{\mathrm{I}}\sigma}
\omega^{(g)}_{k,\vec{m}}
\;=\;
[\dzero,\, \Theta_\cA^{(g)}]_{k, \vec{m}}
\;=\;
(d\,\Theta_\cA)^{(g)}_{k, \vec{m}}.
\]
This uses the factorization~\eqref{eq:typeI-stratum}: the
screen factor $\overline{C}_{|S|}(\mathbb{A}^1)$ computes the
OPE collision $\mu_S$, and the complementary factor
$\overline{C}{}^{\,\mathrm{oc}}_{k-|S|+1, \vec{m}}$ carries
the remaining amplitude, giving the commutator
$[\dzero, \Theta_\cA]$ (the bar differential applied to
$\Theta_\cA$).

\medskip
\noindent\textsc{Type~II + Type~III contribution:
boundary and mixed differentials in $d\,\Theta_\cA$.}

The Type~II and Type~III faces contribute the
\emph{boundary} and \emph{mixed} components of the
open-closed differential~$d$. These are \emph{not} part of
the Lie bracket $[-,-]^{\mathrm{oc}}$: the bracket only
composes interior legs via separating nodal degenerations
(Type~IV(a)), whereas Types~II and~III involve
boundary-leg operations.

Type~II faces (consecutive boundary collisions): the
associahedral fiber $K_{|B|}$ produces the
$A_\infty$-composition on the boundary module. This is the
boundary differential $d_{\mathrm{bdy}}$ acting on
$\Theta_\cA$:
\[
\sum_{B,j}\,
\int_{\partial_{B,j}^{\mathrm{II}}\sigma}
\omega^{(g)}
\;=\;
(d_{\mathrm{bdy}}\,\Theta_\cA)^{(g)}_{k, \vec{m}}.
\]

Type~III faces (mixed bubbling): the Swiss-cheese fiber
$\mathrm{SC}_{1, |B|}$ produces the bulk-to-boundary
intertwining. This is the mixed differential
$d_{\mathrm{mix}}$ acting on $\Theta_\cA$:
\[
\sum_{i,B,j}\,
\int_{\partial_{i,B,j}^{\mathrm{III}}\sigma}
\omega^{(g)}
\;=\;
(d_{\mathrm{mix}}\,\Theta_\cA)^{(g)}_{k, \vec{m}}.
\]

Together with the Type~I contribution
($d_{\mathrm{int}}\,\Theta_\cA$), the three face types
produce the \emph{full} open-closed differential
$d = d_{\mathrm{int}} + d_{\mathrm{bdy}} + d_{\mathrm{mix}}$
acting on~$\Theta_\cA$. The Lie bracket
$[-,-]^{\mathrm{oc}}$ sums only over \emph{separating nodal
degenerations} (Type~IV(a)), not over boundary or mixed
collisions.

\medskip
\noindent\textsc{Type~IV contribution:
bracket $\tfrac{1}{2}[\Theta_\cA, \Theta_\cA]^{\mathrm{oc}}$
and clutching $\hbar\,\Delta_{\mathrm{clutch}}(\Theta_\cA)$.}

The Type~IV faces produce the nodal-degeneration terms and
split into two geometrically distinct classes.
The \emph{separating} nodal faces Type~IV(a) are the
\emph{sole} source of the bracket
$\frac{1}{2}[\Theta_\cA, \Theta_\cA]^{\mathrm{oc}}$:
each separating node joins two distinct components, composing
two operations via the interior pairing.
The \emph{non-separating} clutching faces Type~IV(b)
contribute the genuinely new term
$\hbar\,\Delta_{\mathrm{clutch}}(\Theta_\cA)$.
(Note: the Type~IV faces arise from the genus-graded modular
structure of the bar differential: they record nodal
degenerations of the underlying curve, not collisions of marked
points on a fixed curve. The fixed-curve Stokes argument on
$\overline{C}{}^{\,\mathrm{oc}}_{k,\vec{m}}$ produces
Types~I--III; the Type~IV contribution enters through the
moduli variation encoded in the genus summation
$\sum_g \hbar^g$ and the modular operad algebra structure of the bar
complex (Theorem~\textup{\ref{thm:bar-modular-operad}}).

\emph{Precise accounting.}
The factor of $\frac{1}{2}$ arises because each separating
nodal degeneration (Type~IV(a)) is counted twice (once for
each orientation of the separating edge): the
$(g_1, k_1, \vec{m}_1)
\circ (g_2, k_2, \vec{m}_2)$ composition with
$g_1 + g_2 = g$ appears both as ``left $\circ$ right'' and
``right $\circ$ left,'' and the Lie bracket
$[\Theta_\cA, \Theta_\cA]$ accounts for both with the
antisymmetry producing the factor of~$2$ that cancels the
$\frac{1}{2}$.

\emph{Separation of clutching from bracket.}
The distinction is structural. In
Getzler--Kapranov's formalism~\cite[\S5]{GeK98}, the
\emph{modular operad bracket} $[\![-,-]\!]^{\mathrm{mod}}$ sums
over \emph{all} stable-graph compositions, including both
separating and non-separating gluings. However, the
\emph{Lie bracket} $[-,-]^{\mathrm{oc}}$ of
Theorem~\ref{thm:oc-convolution-dg-lie} only sums over
\emph{separating} compositions (those that join two distinct
elements of the convolution algebra). The non-separating
compositions (those that contract two legs of a
\emph{single} element) form a separate operation, the
\emph{BV operator}~$\Delta$.

Explicitly, the modular operad differential decomposes as:
\begin{equation}\label{eq:gk-decomposition}
D_{\mathrm{GK}}
\;=\;
d_{\mathrm{int}}
+ [\tau, -]
+ \underbrace{\sum_{\text{sep.\ edges}} \circ_e}_{=\,
 \text{Lie bracket part}}
+ \underbrace{\sum_{\text{non-sep.\ edges}}
 \Delta_e}_{=\, \hbar\Delta}.
\end{equation}
A separating edge $e$ in a stable graph $\Gamma$
disconnects $\Gamma$ into two components upon removal;
the corresponding composition $\circ_e$ joins two
operations (one from each component) and is captured by the
Lie bracket $[-,-]^{\mathrm{oc}}$. A non-separating edge
$e$ leaves $\Gamma$ connected upon removal; the
corresponding operation $\Delta_e$ traces two inputs of a
\emph{single} operation against the pairing
$\langle -, - \rangle_\cA$ and is the BV
operator~$\hbar\Delta$.

In the Maurer--Cartan equation, this separation produces:
\begin{itemize}
\item The differential $d\,\Theta_\cA$ from all collision
 faces (Types~I + II + III), where $d = d_{\mathrm{int}}
 + d_{\mathrm{bdy}} + d_{\mathrm{mix}}$ is the full
 open-closed differential;
\item The bracket $\frac{1}{2}[\Theta_\cA, \Theta_\cA]^{\mathrm{oc}}$
 from all separating nodal degenerations (Type~IV(a));
\item The BV term $\hbar\,\Delta_{\mathrm{clutch}}(\Theta_\cA)$
 from all non-separating nodal degenerations (Type~IV(b)).
\end{itemize}
The non-separating clutching increases genus by~$1$
($g \mapsto g + 1$) and decreases degree by~$2$
($n \mapsto n - 2$), which is why it is weighted by $\hbar$
(the genus-tracking parameter).

\medskip
\noindent\textsc{Assembly.}
Summing the contributions from all four boundary types:
\begin{align}
0 &= \int_{\partial\sigma} \omega^{(g)} \notag \\
 &= \underbrace{\text{Type~I} + \text{Type~II}
 + \text{Type~III}}_{\displaystyle d\,\Theta_\cA}
 \;+\;
 \underbrace{\text{Type~IV(a)}}_{\displaystyle
 \frac{1}{2}[\Theta_\cA, \Theta_\cA]^{\mathrm{oc}}}
 \;+\;
 \underbrace{\text{Type~IV(b)}}_{\displaystyle
 \hbar\,\Delta_{\mathrm{clutch}}(\Theta_\cA)}.
 \label{eq:stokes-assembly}
\end{align}
Since this holds for every smooth chain~$\sigma$ (by the
Stokes theorem on the manifold with corners
$\overline{C}{}^{\,\mathrm{oc}}_{k,\vec{m}}$,
applied using the normal-crossings property of
Theorem~\ref{thm:bordered-fm-properties}(iii)), we obtain
the MC equation~\eqref{eq:modular-mc-clutching}:
\[
d\,\Theta_\cA
+ \tfrac{1}{2}[\Theta_\cA, \Theta_\cA]^{\mathrm{oc}}
+ \hbar\,\Delta_{\mathrm{clutch}}(\Theta_\cA)
= 0.
\]

\emph{Consistency check.}
Restricting to the closed sector
($\vec{m} = \vec{0}$, no boundary modules) recovers the
bar-intrinsic MC equation of Theorem~\ref{thm:mc2-bar-intrinsic}:
the Type~II and Type~III faces are absent, the bracket reduces
to the interior convolution bracket, and
$\Delta_{\mathrm{clutch}}$ reduces to the non-separating BV
operator~$\hbar\Delta$ of
equation~\eqref{eq:D-five-components}. The open-sector terms
($d_{\mathrm{bdy}}$ and $d_{\mathrm{mix}}$) only contribute when
boundary modules are present.
\end{proof}

\begin{remark}[The QME versus the MC equation]
\label{rem:qme-vs-mc}
\index{quantum master equation!versus MC equation}
The modular MC equation~\eqref{eq:modular-mc-clutching} is
equivalent to the \emph{quantum master equation} (QME) in the
BV formalism. Writing $S = \Theta_\cA$ as the ``action
functional,'' the QME reads
$\hbar\,\Delta(S) + \tfrac{1}{2}\{S, S\} = 0$, where
$\{-,-\}$ is the BV antibracket. The identification is:
\begin{itemize}
\item $\{S, S\} = 2(d\,\Theta + \tfrac{1}{2}[\Theta, \Theta])$:
 the BV antibracket combines the chain differential and the
 convolution bracket.
\item $\hbar\,\Delta(S) = \hbar\,\Delta_{\mathrm{clutch}}(\Theta)$:
 the BV operator is the non-separating clutching.
\end{itemize}
The factor of $\frac{1}{2}$ in~\eqref{eq:modular-mc-clutching}
is the QME factor (cf.\ Appendix~\ref{app:signs}, equation~\eqref{eq:qme-cg}: the QME
has $\hbar\Delta S + \frac{1}{2}\{S, S\} = 0$,
not $\hbar\Delta S + \{S, S\} = 0$).
\end{remark}

\begin{remark}[Genus expansion of the MC equation]
\label{rem:genus-expansion-mc}
\index{modular Maurer--Cartan equation!genus expansion}
Expanding~\eqref{eq:modular-mc-clutching} in powers
of~$\hbar$, with $\Theta_\cA = \sum_{g \geq 0}
\hbar^g\, \Theta^{(g)}$ (where $\Theta^{(0)}$ includes the
genus-$0$ boundary and mixed terms):
\begin{enumerate}[label=\textup{(\arabic*)}]
\item \emph{Genus~$0$:}
 $d\,\Theta^{(0)} + \tfrac{1}{2}[\Theta^{(0)}, \Theta^{(0)}]
 = 0$.
 This is the classical MC equation for the $A_\infty$-structure
 on the boundary module (no $\hbar$-corrections).
\item \emph{Genus~$1$:}
 $d\,\Theta^{(1)} + [\Theta^{(0)}, \Theta^{(1)}]
 + \Delta_{\mathrm{clutch}}(\Theta^{(0)}) = 0$.
 The curvature $\kappa(\cA) \cdot \omega_1$ appears through
 $\Theta^{(1)}$, and the non-separating clutching of the
 genus-$0$ data gives the one-loop obstruction.
\item \emph{Genus~$g \geq 2$:}
 $d\,\Theta^{(g)} + \frac{1}{2}\sum_{g_1 + g_2 = g}
 [\Theta^{(g_1)}, \Theta^{(g_2)}]
 + \Delta_{\mathrm{clutch}}(\Theta^{(g-1)}) = 0$.
 Each higher genus is determined inductively from lower genera,
 with the clutching term acting as a ``genus-raising''
 obstruction.
\end{enumerate}
\end{remark}

\begin{corollary}[Recovery of the bar-intrinsic MC element;
\ClaimStatusConditional]
\label{cor:recovery-bar-intrinsic}
\index{bar-intrinsic MC element!recovery from bordered FM}
In the closed sector $($no boundary modules:
$\vec{m} = \vec{0}$, $r = 0)$, the modular MC equation with
clutching~\eqref{eq:modular-mc-clutching} reduces to
\begin{equation}\label{eq:mc-closed-recovery}
[\dzero,\, \Theta_\cA]
+ \tfrac{1}{2}[\Theta_\cA, \Theta_\cA]
+ \hbar\,\Delta(\Theta_\cA) = 0,
\end{equation}
which is precisely the MC equation of
Theorem~\textup{\ref{thm:mc2-bar-intrinsic}}, expressed in the
modular convolution algebra
$\gAmod$
\textup{(}Definition~\textup{\ref{def:modular-convolution-dg-lie})}.
The equivalence $(D_\cA - \dzero)^2 = 0 \iff
[\dzero, \Theta] + \frac{1}{2}[\Theta, \Theta] = 0$ of
Theorem~\textup{\ref{thm:mc2-bar-intrinsic}} is the closed-sector
specialization: the $\hbar\Delta$ term is already contained in
$D_\cA$ via the modular operad differential~\eqref{eq:gk-decomposition},
and the expansion
$0 = (D_\cA)^2 = (\dzero + \Theta_\cA + \hbar\Delta)^2$
reproduces~\eqref{eq:mc-closed-recovery} upon collecting terms.
\end{corollary}

\begin{proof}
When $r = 0$, the boundary intervals $I_{p_j}$, the
boundary differentials $d_{\mathrm{bdy}}$ and $d_{\mathrm{mix}}$,
and all Type~II and Type~III faces are absent. The bordered FM
compactification reduces to $\overline{C}_k(X)$ (over a fixed
curve) or to $\overline{\cM}_{g,k}$ (over moduli). The
open-closed bracket $[-,-]^{\mathrm{oc}}$ reduces to the
interior convolution bracket $[-,-]$ of
Construction~\ref{const:explicit-convolution-bracket}, and
$\Delta_{\mathrm{clutch}}$ reduces to the non-separating BV
operator $\hbar\Delta$ of
Theorem~\ref{thm:convolution-dg-lie-structure}(i).
Equation~\eqref{eq:modular-mc-clutching} then becomes
\[
[\dzero, \Theta_\cA]
+ \tfrac{1}{2}[\Theta_\cA, \Theta_\cA]
+ \hbar\Delta(\Theta_\cA) = 0,
\]
which is the expansion of $D_\cA^2 = 0$ with
$D_\cA = \dzero + \Theta_\cA$ as in the proof of
Theorem~\ref{thm:mc2-bar-intrinsic}(i). (The $\hbar\Delta$
term appears explicitly because we have separated the
non-separating clutching from the rest of $D_\cA$; in the
bar-intrinsic formulation, it is absorbed into the
genus-completed differential.)
\end{proof}


\section{Period coordinates at higher genus}

At genus $g$, additional coordinates come from:
\begin{itemize}
\item Period matrix $\Omega \in \mathcal{H}_g$ (Siegel upper half-space)
\item Marking of homology basis $\{a_i, b_i\}_{i=1}^g$
\item Choice of spin structure (quadratic refinement)
\end{itemize}

These appear in correlation functions through:
\[\langle \prod_i \phi_i(z_i) \rangle_g = \sum_{\text{spin}} \int_{\mathcal{F}_g} d\mu(\Omega) \, F(\Omega, z_i, \phi_i)\]
where $\mathcal{F}_g$ is a fundamental domain for $\text{Sp}(2g, \mathbb{Z})$.

\section{The genus-stratified bar construction}

The total bar complex becomes:
\[\barB(\cA) = \bigoplus_{g=0}^{\infty} \bigoplus_{n=0}^{\infty} \barB^{(g),n}(\cA)\]
with differential respecting the genus stratification (Definition~\ref{def:bar-differential-complete}):
\[d\colon \barB^{(g),n} \to \barB^{(g),n-1} \;\oplus\; \barB^{(g-1),n+2} \;\oplus\; \bigoplus_{\substack{g_1+g_2=g \\ n_1+n_2=n}} \barB^{(g_1),n_1+1} \otimes \barB^{(g_2),n_2+1}\]

The three components correspond to:
\begin{itemize}
\item \emph{Collision} (genus-preserving): merging two insertions at a smooth point
\item \emph{Non-separating node}: $\Sigma_g \to \Sigma_{g-1}$ with two new marked points (pinching a non-separating cycle)
\item \emph{Separating node}: $\Sigma_g \to \Sigma_{g_1} \cup \Sigma_{g_2}$, $g_1 + g_2 = g$ (pinching a separating cycle)
\end{itemize}

% (Configuration space definition already given above in enhanced form)

\begin{proposition}[Fundamental group across genera \cite{Arnold69,Brieskorn73}; \ClaimStatusProvedElsewhere]
\label{prop:fundamental-group-genera}
The fundamental group $\pi_1(C_n(\Sigma_g))$ varies with genus:
the low-genus structure is summarized as follows.
\begin{itemize}
\item \emph{Genus 0:} pure braid group $P_n$ on $n$ strands.
\item \emph{Genus 1:} an extension of $P_n$ by an elliptic braid group with modular data.
\item \emph{Genus $g \geq 2$:} an extension by a surface braid group with mapping class group action.
\end{itemize}

\index{braid group}
For genus 0 ($X = \mathbb{C}$), this is the kernel of $B_n \to S_n$ where $B_n$ is the Artin braid group with generators $\sigma_i$ ($i = 1, \ldots, n-1$) and relations:
\begin{align}
\sigma_i\sigma_j &= \sigma_j\sigma_i \quad \text{if } |i-j| > 1 \\
\sigma_i\sigma_{i+1}\sigma_i &= \sigma_{i+1}\sigma_i\sigma_{i+1} \quad \text{(braid relations)}
\end{align}
\end{proposition}

% (Fulton--MacPherson compactification already covered in detail above)

\begin{example}[Configuration spaces across genera]\label{ex:config-genera}
\emph{Genus 0 ($\mathbb{P}^1$):} The compactification $\overline{M}_{0,4}$ is explicit: it is the Deligne--Mumford compactification of the moduli space of $4$ distinct points on~$\mathbb{P}^1$, and it governs the genus-$0$ bar complex at degree~$1$.
\begin{enumerate}
\item Fix three points using $\mathrm{PSL}_2(\mathbb{C})$: $(z_1, z_2, z_3) = (0, 1, \infty)$. The fourth point $z_4 = \lambda \in \mathbb{C} \setminus \{0,1\}$ is the remaining modular parameter.

\item The compactification $\overline{M}_{0,4}$ adds three boundary divisors:
 \begin{itemize}
 \item $D_{14|23}$: $\lambda \to 0$ (collision $z_4 \to z_1$)
 \item $D_{24|13}$: $\lambda \to 1$ (collision $z_4 \to z_2$)
 \item $D_{34|12}$: $\lambda \to \infty$ (collision $z_4 \to z_3$)
 \end{itemize}

\item Result: $\overline{M}_{0,4} \cong \mathbb{P}^1$ with three marked points $\{0, 1, \infty\}$. These three boundary divisors encode the $A_\infty$ relation $m_2(m_2(a,b),c) - m_2(a, m_2(b,c)) = 0$ (associativity).
\end{enumerate}

\begin{remark}
The FM compactification $\overline{C}_n(\mathbb{P}^1)$ (without the $\mathrm{PSL}_2$ quotient) is a higher-dimensional space: $\dim_{\mathbb{C}} C_n(\mathbb{P}^1) = n$, and $\dim_{\mathbb{C}} \overline{M}_{0,n+1} = n - 2$. The bar complex uses the automorphism-free configuration space $C_n(X)$ for curves $X$ of genus~$g \geq 1$, where $\mathrm{Aut}(X)$ is discrete.
\end{remark}

\emph{Genus 1 (Torus).} For $\Sigma_1 = \mathbb{C}/(\mathbb{Z} + \tau\mathbb{Z})$:
\begin{enumerate}
\item for a fixed elliptic curve, $\tau$ is fixed and the FM boundary
contains only point-collision divisors;
\item in the relative family over the modular curve, $\tau$ varies and
the compactified total space also has the nodal degeneration
$q=e^{2\pi i\tau}=0$;
\item the positive-genus correction is carried by period classes and by
the diagonal differential in Proposition~\ref{prop:arnold-higher-genus}.
\end{enumerate}

\emph{Genus $g \geq 2$.} For $\Sigma_g$:
\begin{enumerate}
\item for a fixed curve, the period matrix is fixed and
$\overline C_n(\Sigma_g)$ has only FM collision strata;
\item stable-graph boundary strata appear in the relative
Deligne--Mumford compactification, not in the fixed-curve FM space;
\item spin structures and theta characteristics enter the analytic
choice of prime-form or Szeg\H{o} propagator, not the Arnold
Orlik--Solomon presentation.
\end{enumerate}

The logarithmic forms at each genus:
\begin{itemize}
\item \emph{Genus 0:} Standard forms $\eta_{ij} = d\log(z_i - z_j)$
\item \emph{Genus 1:} Elliptic forms $\eta_{ij}^{(1)} = d\log\vartheta_1(z_i - z_j|\tau)$ with modular parameter
\item \emph{Genus $g \geq 2$:} Prime forms $\eta_{ij}^{(g)} = d_z\log E(z_i, z_j)$ where $E$ is the Fay--Schottky prime form (a $(-\tfrac{1}{2},-\tfrac{1}{2})$-differential with a unique simple zero at $z_i = z_j$)
\end{itemize}

Key relations (Arnold relations extended):
\begin{itemize}
\item \emph{Genus 0 (Arnold relation):}
$\eta_{ij} \wedge \eta_{jk} + \eta_{jk} \wedge \eta_{ki} + \eta_{ki} \wedge \eta_{ij} = 0$
as a 2-form identity on $C_3(\mathbb{C})$. (Note: the \emph{sum} of 1-forms
$\eta_{12} + \eta_{23} + \eta_{13} = d\log[(z_1{-}z_2)(z_2{-}z_3)(z_3{-}z_1)]$
is exact but nonzero; this is a different identity.)
\item \emph{Genus 0, projective screen:} on $C_n(\mathbb{P}^1)$ one must add
the projective linear relations
$\sum_{j\neq i}\bar\eta_{ij}=0$
from Proposition~\ref{prop:arnold-higher-genus}.
\item \emph{Genus $g \geq 1$:} the curve model carries the mixed period
relations
$\eta_{ij}(p_i^*\xi-p_j^*\xi)=0$
and the diagonal differential
$d\eta_{ij}=\Delta_{ij}$
of Proposition~\ref{prop:arnold-higher-genus}; on elliptic and
higher-genus compactifications these appear analytically as the modular
and prime-form corrections of
Proposition~\ref{prop:elliptic-arnold-relations}.
\end{itemize}

At genus~$0$, the Arnold relation between propagator forms is the precise
identity that forces $d^2 = 0$ on the bar differential at degree~$2$:
the composition of two residue operations is a wedge product of two
propagators, and the Arnold sum vanishes. At genus~$g \geq 1$ the
elliptic and prime-form propagators satisfy period-corrected relations:
the local Arnold term survives on the tangent screen, but the global
configuration-space complex also sees the projective or mixed
Totaro-type relations of Proposition~\ref{prop:arnold-higher-genus}.
The remaining analytic corrections
\textup{(}Proposition~\ref{prop:elliptic-arnold-relations} for $g = 1$\textup{)}
are absorbed into the curvature class
$\kappa(\cA)\cdot \omega_g$.

Configuration spaces therefore encode local OPE data through the
collision divisors and (at genus~$g \geq 1$) global monodromy data
through the period integrals on $\Sigma_g$. On the affine genus-$0$
screen only the local data is present; on $\mathbb{P}^1$ the
projective linear relations remain, and at genus~$g \geq 1$ the period
data also contribute.
\end{example}

\begin{theorem}[FM compactification and associahedra {\cite{Sta63}}; \ClaimStatusProvedElsewhere]\label{thm:fm-associahedron}
For $n \geq 2$, the real Fulton--MacPherson moduli space $\overline{M}_{0,n+1}(\mathbb{R})$ is a cell complex whose face poset is isomorphic to the Stasheff associahedron~$K_{n-1}$. The boundary strata of the FM operad $\{\overline{C}_n(\mathbb{R})\}$ (modulo translations and dilations) correspond bijectively to the faces of~$K_{n-1}$, with codimension-$k$ strata indexed by trees with $k$ internal edges.
\end{theorem}

% Logarithmic forms and blow-up coordinates: see Definition~\ref{def:log-forms},
% Theorem~\ref{thm:local-coords-boundary}, and Proposition~\ref{prop:eta}.

\begin{proposition}[Properties of \texorpdfstring{$\eta_{ij}$}{eta-ij}; \ClaimStatusProvedHere]\label{prop:eta}
On an affine genus-$0$ screen, and in a holomorphic FM chart near a
pair-collision divisor, the forms $\eta_{ij}$ satisfy:
\begin{enumerate}
\item Symmetry: $\eta_{ji} = \eta_{ij}$ as differential forms, since
$d\log(-1)=0$.
\item Blow-up expansion: if $x_{ij}$ is a local equation for $D_{ij}$
and $u_{ij}$ is invertible on the screen chart, then
\[
\eta_{ij}=d\log x_{ij}+d\log u_{ij}.
\]
\item Residue: $\Res_{D_{ij}}\eta_{ij}=1$ with the convention that
$\Res_{x=0}(d\log x)=1$.
\item Closure: $d\eta_{ij}=0$ on the affine/formal screen away from
the pole divisor.
\end{enumerate}
On $\mathbb P^1$ and on compact positive-genus curves this proposition
is only the local tangent-screen statement; the global model also has
the projective and period corrections of
Proposition~\ref{prop:arnold-higher-genus}.
\end{proposition}

\begin{proof}
(1) is immediate from the definition. For (2), compute in the
normal-crossing chart:
\[
z_i-z_j=x_{ij}u_{ij},
\qquad u_{ij}\in\mathcal O^\times .
\]
Therefore
\[
d\log(z_i-z_j)=d\log x_{ij}+d\log u_{ij}.
\]
For (3), the residue extracts the coefficient of $d\log x_{ij}$, which
is $1$ by this normalization.

For (4), since $\eta_{ij}$ is locally $d$ of a function away from other collision divisors, we have $d\eta_{ij} = d^2\log(x_i - x_j) = 0$.
\end{proof}

\subsection{The Orlik--Solomon algebra}

The logarithmic forms $\eta_{ij}$ generate a differential graded algebra with the following properties:

The Arnold relations (Theorem~\ref{thm:arnold-relations}) hold in the OS algebra.

\begin{theorem}[Cohomology via Orlik--Solomon {\cite{Arnold69,OS80}}; \ClaimStatusProvedElsewhere]\label{thm:os-cohomology-config}
\index{Orlik--Solomon algebra|textbf}
\index{braid group!arrangement}
\index{Poincare series}
For $X = \mathbb{C}$, the cohomology of the \emph{open} configuration space $C_n(\mathbb{C})$ is:
\[
H^*(C_n(\mathbb{C})) \cong \text{OS}(A_{n-1})
\]
where $\text{OS}(A_{n-1})$ is the Orlik--Solomon algebra of the braid arrangement $A_{n-1}$. The Poincar\'e polynomial is:
\[
\sum_{k=0}^{n-1} \dim H^k(C_n(\mathbb{C})) \cdot t^k = \prod_{i=1}^{n-1}(1 + it)
\]
\end{theorem}

\begin{remark}[Provenance and citation]
This theorem is imported and treated as \ClaimStatusProvedElsewhere. The
cohomology identification with the Orlik--Solomon algebra and the corresponding
Poincar\'e polynomial are classical results in configuration-space/arrangement
theory; see~\cite{Arnold69,OS80}.
\end{remark}

\subsection{No-broken-circuit bases}

Explicit computations use concrete cohomology bases:

\begin{definition}[Broken circuit]
Fix a total order on pairs $(i, j)$ with $i < j$; this chapter uses
lexicographic order. A \emph{broken circuit} is a set obtained by
removing the minimal element from a circuit (minimal dependent set) in
the graphical matroid on $K_n$.
\end{definition}

\begin{definition}[NBC basis]\label{def:nbc-basis}
A \emph{no-broken-circuit (NBC)} set is a collection of pairs that contains no broken circuit. These correspond bijectively to:
\begin{itemize}
\item Acyclic directed graphs on $[n]$ (forests)
\item Independent sets in the graphical matroid
\item Monomials in $\eta_{ij}$ that do not vanish by Arnold relations
\end{itemize}
\end{definition}

\begin{theorem}[NBC basis theorem {\cite{OS80}}; \ClaimStatusProvedElsewhere]\label{thm:NBC}
The NBC sets provide a basis for $H^*(C_n(\mathbb{C}))$. More precisely, if $F$ is an NBC forest with edges $E(F) = \{(i_1, j_1), \ldots, (i_k, j_k)\}$, then:
\[
\omega_F = \eta_{i_1j_1} \wedge \cdots \wedge \eta_{i_kj_k}
\]
forms a basis element of $H^k(C_n(\mathbb{C}))$
(the cohomology of the \emph{open} configuration space, not the FM compactification).
\end{theorem}

\begin{remark}[Provenance and citation]
This theorem is imported and treated as \ClaimStatusProvedElsewhere. The NBC
basis theorem for the Orlik--Solomon algebra is standard in the arrangement
literature; see~\cite{OS80}.
\end{remark}

\begin{example}[NBC basis for \texorpdfstring{$n = 4$}{n = 4}]\label{ex:NBC4}
For $C_4(\mathbb{C})$, using the lexicographic order on pairs, the NBC basis consists of:
\begin{itemize}
\item Degree 0: $1$
\item Degree 1: $\eta_{12}, \eta_{13}, \eta_{14}, \eta_{23}, \eta_{24}, \eta_{34}$ (6 elements)
\item Degree 2: $\eta_{12} \wedge \eta_{34}, \eta_{13} \wedge \eta_{24}, \eta_{14} \wedge \eta_{23}$, plus 8 other terms (11 total)
\item Degree 3: $\eta_{12} \wedge \eta_{23} \wedge \eta_{34}$ and 5 other spanning trees (6 total)
\end{itemize}
Total: $1 + 6 + 11 + 6 = 24 = 4!$ basis elements, confirming $\dim H^*(C_4(\mathbb{C})) = 4!$.
\end{example}

The operadic framework, configuration spaces, and logarithmic forms
developed above form the ingredients for the geometric bar complex
constructed in Chapter~\ref{chap:bar-cobar}.

\section{Configuration spaces, factorization and higher genus}

\subsection{The Ran space and chiral operations}

\begin{definition}[D-module category]
\index{D-module|textbf}
Let $\text{D-mod}_{rh}(X)$ denote the category of regular holonomic
D-modules on $X$.
These are D-modules $\mathcal{M}$ satisfying:
\begin{enumerate}
\item Finite presentation: locally finitely generated over $\mathcal{D}_X$
\item Regular singularities: characteristic variety is Lagrangian
\item Holonomicity: $\text{dim}(\text{Char}(\mathcal{M})) = \text{dim}(X)$
\end{enumerate}
This category has:
\begin{itemize}
\item Six functors: $f^*, f_*, f^!, f_!, \otimes^L, \mathcal{RHom}$
\item Riemann--Hilbert correspondence with perverse sheaves
\item Well-defined maximal extension $j_*j^*$ for $j: U \hookrightarrow X$ open
\end{itemize}
\end{definition}

\begin{definition}[Ran space via categorical colimit]\label{def:ran-precise}
\index{Ran space|textbf}
Let $X$ be a smooth algebraic curve over $\mathbb{C}$. The \emph{Ran space} of $X$ is the ind-scheme
defined as the colimit:
\[
\text{Ran}(X) = \underset{I \in \text{FinSet}^{\text{surj,op}}}{\text{colim}} \, X^I
\]
where:
\begin{itemize}
\item $\text{FinSet}^{\text{surj}}$ is the category of finite sets with surjections as morphisms
\item For a surjection $\phi: I \twoheadrightarrow J$, the induced map $X^J \to X^I$ is the diagonal
embedding on fibers $\phi^{-1}(j)$
\item The colimit is taken in the category of ind-schemes with the Zariski topology
\end{itemize}
Explicitly, a point in $\text{Ran}(X)$ is a non-empty finite subset $S \subset X$ (without multiplicities: the colimit over surjections identifies $\{x, x\}$ with $\{x\}$). This is distinct from the space of effective 0-cycles $\text{Sym}^+(X) = \coprod_n \text{Sym}^n(X)$, which records multiplicities.
\end{definition}

\begin{remark}[Set-theoretic description]
The underlying set of $\text{Ran}(X)$ is the set of all non-empty finite subsets of~$X$
(without multiplicities, by the identification convention in the colimit). This is distinct from
the free commutative monoid on~$X$, which records multiplicities. The ind-scheme structure
on~$\text{Ran}(X)$ encodes the deformation theory of point configurations.
\end{remark}

The Ran space carries a fundamental monoidal structure encoding disjoint union:

\begin{definition}[Factorization structure]\label{def:factorization}
\index{factorization!locality}
The union correspondence on the Ran space is not a proper morphism:
\[
\Delta\colon \Ran(X)\times\Ran(X)\longrightarrow \Ran(X),
\qquad (S,T)\longmapsto S\cup T .
\]
Consequently the expression
\[
\mathcal{M} \otimes^{\text{ch}} \mathcal{N}
= \Delta_! \left( \rho_1^* \mathcal{M} \otimes^! \rho_2^* \mathcal{N} \right)
\]
is not a functor on all $\mathcal D$-modules on $\Ran(X)$. The
factorization product is defined instead on the disjoint locus, with
extension data across the diagonals supplied by the chiral product and
the factorization algebra axioms.
\end{definition}

\begin{definition}[Factorization algebra]\label{def:fact-algebra-correct}
A \emph{factorization algebra} $\mathcal{F}$ on $X$ consists of:
\begin{enumerate}
\item A quasi-coherent $\mathcal{D}$-module $\mathcal{F}_S$ for each finite set $S \subset X$
\item For disjoint $S_1, S_2$, a factorization isomorphism:
 \[\mu_{S_1,S_2}: \mathcal{F}_{S_1} \boxtimes \mathcal{F}_{S_2} \xrightarrow{\sim} \mathcal{F}_{S_1 \sqcup S_2}\]
\item These satisfy:
 \begin{itemize}
 \item \emph{Associativity:} For disjoint $S_1, S_2, S_3$:
 \begin{center}
 \begin{tikzcd}
 \mathcal{F}_{S_1} \boxtimes \mathcal{F}_{S_2} \boxtimes \mathcal{F}_{S_3}
 \arrow[r, "\mu_{S_1,S_2} \boxtimes \text{id}"]
 \arrow[d, "\text{id} \boxtimes \mu_{S_2,S_3}"'] &
 \mathcal{F}_{S_1 \sqcup S_2} \boxtimes \mathcal{F}_{S_3}
 \arrow[d, "\mu_{S_1 \sqcup S_2, S_3}"] \\
 \mathcal{F}_{S_1} \boxtimes \mathcal{F}_{S_2 \sqcup S_3}
 \arrow[r, "\mu_{S_1, S_2 \sqcup S_3}"'] &
 \mathcal{F}_{S_1 \sqcup S_2 \sqcup S_3}
 \end{tikzcd}
 \end{center}
 \item \emph{Commutativity:} $\mu_{S_2,S_1} = \sigma_{S_1,S_2} \circ \mu_{S_1,S_2}$ where $\sigma$ is the swap
 \item \emph{Unit:} $\mathcal{F}_\emptyset = \mathbb{C}$ with canonical isomorphisms $\mathcal{F}_S \cong \mathbb{C} \boxtimes \mathcal{F}_S$
 \end{itemize}
\end{enumerate}
\end{definition}

\begin{remark}[Kontsevich's formulation]
Factorization algebras encode locality: observables
on disjoint regions combine independently, and the factorization isomorphisms express the commutativity of spacelike-separated observables (Kontsevich; Costello--Gwilliam~\cite{CG17}).
\end{remark}

\begin{theorem}[Chiral algebras as factorization algebras \cite{BD04}; \ClaimStatusProvedElsewhere]\label{thm:chiral-as-fact}
\index{chiral algebra!as factorization algebra}
Every chiral algebra $\mathcal{A}$ on $X$ determines a factorization algebra $\mathcal{F}_\mathcal{A}$ where:
\begin{itemize}
\item $\mathcal{F}_\mathcal{A}(S) = \mathcal{A}^{\boxtimes S}$ for finite $S \subset X$
\item The factorization structure comes from the chiral multiplication
\item This defines a fully faithful functor $\text{ChirAlg}(X) \to \text{FactAlg}(X)$
\end{itemize}
\end{theorem}

\begin{proof}[Proof following Beilinson--Drinfeld]
Chiral multiplication provides exactly the factorization isomorphisms needed.
The Jacobi identity for chiral algebras translates to associativity of factorization. The technical
issue with properness is avoided because we work fiberwise over finite sets rather than globally on Ran space.
\end{proof}
% Add precise D-module structure


\begin{theorem}[Factorization monoidal structure {\cite{BD04,CG17}}; \ClaimStatusProvedElsewhere]\label{thm:fact-monoidal-corrected}
The category $\text{FactAlg}(X)$ of factorization algebras forms a symmetric monoidal
category with:
\begin{enumerate}
\item Tensor product: $(\mathcal{F} \otimes_{\text{fact}} \mathcal{G})(S) = \bigoplus_{S_1 \sqcup S_2 = S} \mathcal{F}(S_1) \otimes \mathcal{G}(S_2)$
\item Unit: The vacuum factorization algebra $\mathbb{1}$ with $\mathbb{1}(S) = \begin{cases} \mathbb{C} & S = \emptyset \\ 0 & \text{otherwise} \end{cases}$
\item Associativity isomorphism satisfying the pentagon axiom
\item Braiding isomorphism induced by the symmetric group action
\end{enumerate}

There is a fully faithful embedding:
\[\text{ChirAlg}(X) \hookrightarrow \text{FactAlg}(X)\]
sending a chiral algebra $\mathcal{A}$ to its associated factorization algebra $\mathcal{F}_{\mathcal{A}}$.
\end{theorem}

\begin{proof}
The tensor product is defined by Day convolution over the symmetric monoidal structure on finite sets under disjoint union. The pentagon and hexagon axioms for $\otimes_{\mathrm{fact}}$ reduce to the corresponding axioms for disjoint union of finite sets, which are immediate. The unit axiom $\mathbb{1} \otimes_{\mathrm{fact}} \mathcal{F} \simeq \mathcal{F}$ follows from $\mathbb{1}(\emptyset) = \mathbb{C}$.

For the embedding $\mathrm{ChirAlg}(X) \hookrightarrow \mathrm{FactAlg}(X)$: given a chiral algebra $\mathcal{A}$, define $\mathcal{F}_{\mathcal{A}}(S) = j_{S*}j_S^*\mathcal{A}^{\boxtimes S}$ for finite $S$, with factorization isomorphisms coming from the chiral product (Beilinson--Drinfeld~\cite{BD04}, Theorem~3.4.9). Fully faithfulness follows because chiral algebras are determined by their factorization data on pairs of points (the chiral product $\mu$).
\end{proof}

\emph{Underlying D-modules.} A collection $\{\mathcal{A}_n\}_{n \geq 0}$ where each $\mathcal{A}_n$ is a quasi-coherent $\mathcal{D}_{X^n}$-module, meaning:
\begin{itemize}
\item $\mathcal{A}_n$ is a sheaf of modules over the sheaf of differential operators $\mathcal{D}_{X^n}$
\item The action satisfies the Leibniz rule: $\partial(fs) = (\partial f)s + f(\partial s)$ for local functions $f$ and sections $s$
\item $\mathcal{A}_n$ is quasi-coherent as an $\mathcal{O}_{X^n}$-module
\end{itemize}

\subsection{Elliptic configuration spaces and theta functions}

\subsubsection{The genus 1 realm: elliptic curves as quotients}

For genus 1, we work with elliptic curves $E_\tau = \mathbb{C}/(\mathbb{Z} + \tau\mathbb{Z})$ where $\tau \in \mathfrak{h}$ lies in the upper half-plane:

\begin{definition}[Elliptic configuration space]
For an elliptic curve $E_\tau$, the configuration space of $n$ points is:
\[
C_n(E_\tau) = \{(z_1, \ldots, z_n) \in E_\tau^n \mid z_i \neq z_j \text{ mod } \Lambda_\tau\}
\]
where $\Lambda_\tau = \mathbb{Z} + \tau\mathbb{Z}$ is the period lattice.
\end{definition}

\begin{theorem}[Elliptic compactification {\cite{Fay73}}; \ClaimStatusProvedElsewhere]
\label{thm:elliptic-compactification}
The compactification $\overline{C_n(E_\tau)}$ is constructed via:
\begin{enumerate}
\item \emph{Local blow-ups}: Near collision points, use elliptic blow-up coordinates
\item \emph{Global structure}: The compactified space admits a stratification by \emph{stable elliptic graphs}
\item \emph{Modular invariance}: Under $SL_2(\mathbb{Z})$ action on $\tau$, the construction is equivariant
\end{enumerate}
\end{theorem}

\begin{proof}[Construction]
Near a collision point $z_i \to z_j$ on $E_\tau$, introduce elliptic blow-up coordinates:
\begin{align}
\epsilon_{ij} &= |z_i - z_j|_{E_\tau} \quad \text{(elliptic distance)} \\
\theta_{ij} &= \arg(z_i - z_j) \quad \text{(angular parameter)} \\
u_{ij} &= \frac{z_i + z_j}{2} \quad \text{(center on } E_\tau)
\end{align}

The key difference from genus 0: the elliptic distance involves the Weierstrass $\sigma$-function:
\[
|z_i - z_j|_{E_\tau} = |\sigma(z_i - z_j; \tau)|e^{-\eta_1\,\text{Im}(z_i - z_j)^2/\text{Im}(\tau)}
\]
where $\eta_1 = \zeta(\tfrac{1}{2}; \Lambda_\tau)$ is the quasi-period of the Weierstrass zeta function (not to be confused with the Dedekind eta function $\eta(\tau)$).
\end{proof}

\subsubsection{Theta functions as building blocks}

The logarithmic forms on elliptic curves are replaced by forms built from theta functions:

\begin{definition}[Elliptic logarithmic forms]
On $\overline{C_n(E_\tau)}$, define the elliptic analogs of $\eta_{ij}$:
\[
\eta_{ij}^{(1)} = d\log\theta_1\left(\frac{z_i - z_j}{2\pi i}; \tau\right) + \frac{\pi}{\mathrm{Im}(\tau)}\, d(\overline{z_i - z_j})
\]
where $\theta_1(z|\tau) = -i\sum_{n \in \mathbb{Z}}(-1)^n q^{(n+1/2)^2/2}e^{i\pi(2n+1)z}$ with $q = e^{2\pi i\tau}$.
\end{definition}

\begin{proposition}[Elliptic correction to the Arnold relation \cite{Fay73}; \ClaimStatusProvedElsewhere]
\label{prop:elliptic-arnold-relations}
The elliptic logarithmic forms do not satisfy the affine Arnold
relation as a global identity on $C_3(E_\tau)$. Their cyclic wedge sum
is a period correction:
\[
\eta_{ij}^{(1)} \wedge \eta_{jk}^{(1)}
+ \eta_{jk}^{(1)} \wedge \eta_{ki}^{(1)}
+ \eta_{ki}^{(1)} \wedge \eta_{ij}^{(1)}
= \mathcal P_{ijk}(\tau),
\]
where $\mathcal P_{ijk}(\tau)$ is the translation-invariant
$(1,1)$ period term determined by the normalized elliptic propagator
and the diagonal class in the genus-one
Totaro--K\v{r}\'{\i}\v{z} model. Its scalar coefficient depends on the
theta-function normalization; the only invariant point used here is
that the term is generally nonzero globally and vanishes on the
tangent-screen associated graded. Thus the affine Arnold relation is
recovered on a formal disk. The form $\eta_{ij}^{(1)}$ is \emph{not}
holomorphic: the regularization term
$\frac{\pi}{\mathrm{Im}(\tau)}\,d\bar{z}$ ensures double periodicity
but introduces a $(0,1)$-component, so the cyclic wedge can produce a
$(1,1)$ period term.
\end{proposition}

The non-vanishing right-hand side encodes the central extension that appears at genus 1.

\subsection{Higher genus configuration spaces}

\subsubsection{Hyperbolic surfaces and Teichmüller theory}

For genus $g \geq 2$, the underlying curve $\Sigma_g$ admits a
hyperbolic metric, but the fixed-curve configuration space is still the
open complement of diagonals in $\Sigma_g^n$. The metric and period
matrix enter the propagator; they do not add stable-curve boundary
components to the Fulton--MacPherson compactification of a fixed curve.

\begin{definition}[Higher genus configuration]
For a compact Riemann surface $\Sigma_g$ of genus $g \geq 2$:
\[
C_n(\Sigma_g) = \{(p_1, \ldots, p_n) \in \Sigma_g^n \mid p_i \neq p_j \text{ for } i \neq j\}
\]
The compactification $\overline{C}_n(\Sigma_g)$ is the
Fulton--MacPherson compactification of this fixed smooth curve. Its
boundary consists of collision divisors $D_S$ and their nested
intersections. Stable curves and Deligne--Mumford boundary strata enter
only after replacing the fixed curve by a relative family over
$\overline{\mathcal M}_g$:
\begin{itemize}
\item fixed curve: FM collision strata inside $\overline{C}_n(\Sigma_g)$;
\item relative modular family: collision strata together with the
 pullback of the log boundary of $\overline{\mathcal M}_g$;
\item stable graphs: clutching strata in the relative modular operad,
 not boundary components of $\overline{C}_n(\Sigma_g)$ for a fixed
 smooth~$\Sigma_g$.
\end{itemize}
\end{definition}

\begin{remark}[Fiberwise and total higher-genus differentials]
\label{rem:period-integrals-bar}
On a fixed positive-genus curve there is a fiberwise collision-residue
operator $\dfib$ built from the prime-form propagator. It is not the
strict total differential unless the curvature vanishes:
\[
\dfib^{\,2}=\kappa(\cA)\,\omega_g\cdot\mathrm{id}
\]
with $\omega_g$ the Arakelov $(1,1)$-form, as in
Proposition~\ref{prop:genus-g-curvature-package}. The strict relative
differential is
\[
\Dg{g}=\dfib+d_{\mathrm{per}}^{(g)}
+\nabla_{\mathrm{KS}}^{\mathrm{GM}},
\qquad
\Dg{g}^{\,2}=0,
\]
where the period operator comes from the Abel--Jacobi and period-matrix
terms in the Arakelov-normalized prime-form propagator, and the
Gauss--Manin term appears only when the curve varies. Thus the
decomposition is not a sum of three independent differentials on
$\overline{C}_n(\Sigma_g)$. The fixed-curve object carries collision
residues and period-corrected propagators; the relative modular object
adds base variation and clutching.
\end{remark}

\subsection{Finite-window residue control}

\begin{definition}[Finite-window residue datum]
\label{def:finite-window-residue-datum}
A chiral algebra $\cA=\omega_X\oplus\bar\cA$ has a
\emph{finite-window residue datum} if it is conformally graded with
finite-dimensional weight spaces, its OPEs have finite singular part in
each pair of homogeneous inputs, and the reduced-weight filtration
\[
\rho(sa_1|\cdots|sa_\ell)
=\sum_{r=1}^{\ell}(\wt(a_r)-1)
\]
has finite windows
\[
K_q(\cA)=B_{\rho(\le q)}(\cA)
\]
stable under the bar differential. For an infinite tower of algebras or
generators, the transition maps between the windows are required to be
Mittag--Leffler so that cohomology commutes with the inverse limit.
\end{definition}

\begin{conjecture}[Analytic FM integral convergence]\label{conj:FM-convergence}
Absolute convergence of unrenormalized configuration-space integrals is
an analytic condition, not the definition of the chiral bar complex. In
a fixed finite window $K_q(\cA)$ and fixed genus range, one may ask for
exponents or counterterms for which the ordinary integrals
\[
\int_{\ConfigSpace{n}} |\phi_1(z_1)\cdots\phi_n(z_n)|^2
\prod_{i<j}|z_i-z_j|^{2\alpha_{ij}}
\]
converge on the FM compactification. This conjectural analytic
condition is separate from the algebraic existence of the bar
differential, which is defined by Poincar\'e residues of logarithmic
forms and finite OPE singular parts.
\end{conjecture}

\begin{remark}[Residue definition versus \(L^2\)-integrability]
Near a collision divisor $D_{ij}$, an unrenormalized integrand with
lowest target weight $h_{\min}$ has the model growth
\[
|\phi_i(z_i)\phi_j(z_j)|^2
\sim |z_i-z_j|^{-2(h_i+h_j-h_{\min})},
\qquad
|d\log(z_i-z_j)|^2
=\frac{|dz_i-dz_j|^2}{|z_i-z_j|^2}.
\]
The local polar integral behaves as
\[
\int_\epsilon^1
r^{1-2(h_i+h_j-h_{\min}+1)}\,dr,
\]
so naive $L^2$ convergence would force
$h_i+h_j-h_{\min}<0$. This condition is usually false in positive-energy
theories. The bar differential therefore uses the logarithmic
Poincar\'e residue and the chiral product, not an $L^2$ integral of the
unregularized OPE. Under the finite-window datum
Definition~\ref{def:finite-window-residue-datum}, each residue
calculation is finite, and the completed object is the inverse limit of
the finite windows, as in
Theorem~\ref{thm:reduced-weight-finiteness} and
Definition~\ref{def:finite-window}.
\end{remark}

\begin{remark}[Regularization]
When analytic integrals rather than residue operations are required,
regularization must be specified as part of the analytic model:
\begin{enumerate}
\item analytic continuation in exponents or dimension;
\item point-splitting with counterterms supported on FM boundary strata;
\item Pauli--Villars or BV renormalization for quantum corrections.
\end{enumerate}
\end{remark}

\subsection{Orientation conventions for configuration spaces}

\begin{definition}[Oriented configuration space]
\index{orientation!line bundle}
\index{configuration space!unordered}
The (ordered) configuration space $C_n(X)$ inherits an orientation from $X^n$ via restriction:
\[\text{or}(C_n(X)) = \text{or}(X)^{\otimes n}\big|_{C_n(X)}\]
since $C_n(X) \subset X^n$ is an open submanifold. (For the \emph{unordered} configuration space $\mathrm{UConf}_n(X) = C_n(X)/S_n$, one must also account for the sign representation of $S_n$.)
\end{definition}

\begin{definition}[Orientation of compactification]
The Fulton--MacPherson compactification $\ConfigSpace{n}$ is oriented by:
\begin{enumerate}
\item Choose orientation on $C_n(X)$ as above
\item At each blow-up, orient exceptional divisors via the outward normal convention
\item The boundary $\partial\ConfigSpace{n} = D$ inherits the outward normal orientation
\end{enumerate}
\end{definition}

\begin{lemma}[Orientation compatibility; \ClaimStatusProvedHere]
\label{lem:orientation-compatibility}
For the stratification of $\partial\ConfigSpace{n}$:
\[\partial\ConfigSpace{n} = \bigcup_{I \subset \{1,\ldots,n\}, |I| \geq 2} D_I\]
The orientations satisfy:
\[\text{or}(\partial D_I) = (-1)^{\text{codim}(D_I)} \text{or}(D_I)\]
\end{lemma}

\begin{proof}
Proceed by induction on codimension.

\emph{Codimension 1}: $D_{ij}$ has orientation from the normal bundle:
\[\text{or}(D_{ij}) = \text{or}(N_{D_{ij}}) \wedge \text{or}(\ConfigSpace{n-1})\]
where $N_{D_{ij}}$ is oriented by $d\epsilon_{ij}$ (radial coordinate).

\emph{Codimension 2}: At $D_{ijk}$ (the stratum where all three points $z_i, z_j, z_k$ collide simultaneously):
\[\text{or}(D_{ijk}) = \text{or}(N_{D_{ij}}) \wedge \text{or}(N_{D_{jk}|D_{ij}}) \wedge \text{or}(\ConfigSpace{n-2})\]

The key sign:
\[\text{or}(D_{ijk})|_{D_{ij} \to D_{ijk}} = -\text{or}(D_{ijk})|_{D_{jk} \to D_{ijk}}\]

This ensures Stokes' theorem holds:
\[\int_{\partial D_{ij}} \omega = \sum_{k} \epsilon_k \int_{D_{ijk}} \omega\]
with appropriate signs $\epsilon_k = \pm 1$.
\end{proof}

\begin{theorem}[Stokes on configuration spaces \cite{FM94}; \ClaimStatusProvedElsewhere]
\label{thm:stokes-config-spaces}
\index{Stokes theorem!on configuration spaces}
For $\omega \in \Omega^{n-1}(\ConfigSpace{n})$:
\[\int_{\ConfigSpace{n}} d\omega = \int_{\partial\ConfigSpace{n}} \omega = \sum_{I} \epsilon_I \int_{D_I} \omega\]
where $\epsilon_I$ is determined by the orientation convention.
\end{theorem}

% The explicit D-module formulation of chiral multiplication
% (extending Definition~\ref{def:fact-algebra-correct}) uses the
% maximal extension $j_{ij*}j_{ij}^*$ to encode regularity away
% from coincident points, with multiplication maps
% $\mu_{ij}\colon j_{ij*}j_{ij}^*(\cA_m \boxtimes \cA_n)
% \to \Delta_*\cA_{m+n-1}$ extracting the OPE singularity.
% Associativity of these multiplication maps is encoded by
% the commutativity of the triple-collision diagram
% (see Beilinson--Drinfeld~\cite[Chapter~3]{BD04}).

\subsection{The chiral endomorphism operad}

For any D-module $\mathcal{M}$ on $X$, we construct the operad controlling chiral algebra structures:

\begin{definition}[Chiral endomorphisms]\label{def:chiral-endo}
The \emph{chiral endomorphism operad} of a D-module $\mathcal{M}$ on $X$ is defined by:
\[
\text{End}_{\mathcal{M}}^{\text{ch}}(n) = \text{Hom}_{\mathcal{D}(X^n)}\left(j_*j^*\mathcal{M}^{\boxtimes n}, \Delta_*\mathcal{M}\right)
\]
where:
\begin{itemize}
\item $j: C_n(X) \hookrightarrow X^n$ is the inclusion of the configuration space
\item $\Delta: X \hookrightarrow X^n$ is the small diagonal
\item The morphisms are taken in the derived category of D-modules
\end{itemize}
\end{definition}

\begin{proposition}[Operadic structure; \ClaimStatusProvedHere]
\label{prop:operadic-structure}
$\text{End}_{\mathcal{M}}^{\text{ch}}$ forms an operad in the category of D-modules with:
\begin{enumerate}
\item Composition: For $f \in \text{End}_{\mathcal{M}}^{\text{ch}}(k)$ and $g_i \in \text{End}_{\mathcal{M}}^{\text{ch}}(n_i)$,
\[
f \circ (g_1, \ldots, g_k) = f \circ \left(\Delta_{n_1,\ldots,n_k}^* (g_1 \boxtimes \cdots \boxtimes g_k)\right)
\]
where $\Delta_{n_1,\ldots,n_k}: X^{n_1 + \cdots + n_k} \to X^k \times X^{n_1} \times \cdots \times X^{n_k}$

\item Unit: The identity map $\text{id}_{\mathcal{M}} \in \text{End}_{\mathcal{M}}^{\text{ch}}(1)$

\item The composition satisfies associativity up to coherent isomorphism
\end{enumerate}
\end{proposition}

\begin{proof}
Associativity follows from the functoriality of the diagonal embeddings. Consider the diagram:
\[
X^{n_1 + \cdots + n_k} \xrightarrow{\Delta_{n_1,\ldots,n_k}} X^k \times \prod_i X^{n_i}
\xrightarrow{\text{id} \times \prod_i \Delta_{m_{i1},\ldots}} X^k \times \prod_i \prod_j X^{m_{ij}}
\]
The two ways of composing correspond to different factorizations of the total diagonal, which are
canonically isomorphic. The coherence follows from the coherence theorem for operads.
\end{proof}

\begin{theorem}[Chiral algebras as algebra objects \cite{BD04}; \ClaimStatusProvedElsewhere]
\label{thm:chiral-algebra-objects}
A chiral algebra structure on $\mathcal{M}$ is equivalent to an algebra structure over the operad
$\text{End}_{\mathcal{M}}^{\text{ch}}$ in the symmetric monoidal category of D-modules. This equivalence is functorial and preserves quasi-isomorphisms.
\end{theorem}


 \section{Chain-level constructions and simplicial models}

\subsection{NBC bases and computational optimality}

The NBC basis (Definition~\ref{def:nbc-basis}) gives an exact normal
form for the affine Orlik--Solomon part of the geometric bar complex.

\begin{theorem}[Exact NBC reduction for the affine OS component; \ClaimStatusProvedHere]
\label{thm:nbc-basis-optimality}
For the braid arrangement component of $C_n(\mathbb C)$, the NBC basis
satisfies:
\begin{enumerate}
\item Each basis element is
$\eta_F=\bigwedge_{e\in F}\eta_e$ for an NBC forest $F$, with the wedge
ordered by the fixed edge order.
\item The raw residue along an edge $e$ in the $p$-th wedge position is
edge deletion with sign $(-1)^{p-1}$; any further integer coefficients
come only from reducing non-NBC monomials by the Orlik--Solomon circuit
relations.
\item The identity $d^2=0$ is exactly the Arnold/Orlik--Solomon circuit
cancellation. There is no additional geometric compactification relation
hidden in the affine scalar propagator.
\item The Poincare polynomial is
\[
P_{C_n(\mathbb C)}(t)=\prod_{q=1}^{n-1}(1+qt),
\]
so the total NBC count is $n!$ and the top degree count is
$(n-1)!$.
\end{enumerate}
\end{theorem}

\begin{proof}
The NBC basis theorem (Theorem~\ref{thm:NBC}) identifies the NBC
monomials with a basis of the Orlik--Solomon algebra of the braid
arrangement. If
$\eta_F=\eta_{e_1}\wedge\cdots\wedge\eta_{e_k}$ and $e_p=e$, the
one-edge residue/deletion operator gives
\[
\operatorname{Res}_{e}(\eta_F)
=
(-1)^{p-1}
\eta_{e_1}\wedge\cdots\wedge\widehat{\eta_{e_p}}\wedge\cdots
\wedge\eta_{e_k}.
\]
This is the Koszul sign convention used by the ordered tensor coalgebra.
When deletion or contraction produces a broken circuit, the result is
rewritten by the OS circuit relation; this is the only source of
integer coefficients beyond the raw sign.

The square of the residue differential is a signed sum over deleting two
edges. For every circuit, the two deletion orders occur with opposite
Koszul signs, and the remaining discrepancy is exactly the
Arnold--Orlik--Solomon relation of
Theorem~\ref{thm:arnold-orlik-solomon}. Thus $d^2=0$ is not a statement
that cancellations are absent. It is the statement that all
codimension-$2$ affine scalar-propagator contributions cancel through
the OS circuit ideal.

Finally,
Theorem~\ref{thm:os-cohomology-config} gives
$P_{C_n(\mathbb C)}(t)=\prod_{q=1}^{n-1}(1+qt)$. Evaluation at $t=1$
gives the total dimension $n!$, and the coefficient of $t^{n-1}$ is
$(n-1)!$. Since the NBC monomials are a basis, these are the
corresponding total and top NBC counts.
\end{proof}

\begin{proposition}[NBC sparsity analysis; \ClaimStatusProvedHere]\label{prop:nbc-sparsity}
For the geometric bar complex, the differential matrix (in the NBC basis at degree~$n$) has at most $O(n^3)$ total non-zero entries due to weight constraints.
\end{proposition}

\begin{proof}
Consider NBC forests $F_1, F_2$ on $n$ vertices. A non-zero differential $\langle dF_1, F_2 \rangle$ requires:
\begin{enumerate}
\item $F_2$ obtained from $F_1$ by contracting one edge $(i,j)$
\item The weight condition $h_{\phi_i} + h_{\phi_j} = h_{\phi_k} + 1$ for some resulting field $\phi_k$
\end{enumerate}

For a chiral algebra with $r$ generators of weights $\{h_1, \ldots, h_r\}$:
- Each vertex can be labeled by one of $r$ generators
- Weight-preserving collisions form a sparse $r \times r$ matrix $M_{ij}$
- $M_{ij} \neq 0$ only if $h_i + h_j \in \{h_k + 1 : k = 1, \ldots, r\}$

The sparsity factor is:
$\rho = \frac{|\{(i,j,k) : h_i + h_j = h_k + 1\}|}{r^3} \leq \frac{r^2}{r^3} = \frac{1}{r}$

Total non-zero entries: $\leq n \cdot \binom{n-1}{2} \cdot \rho \cdot |\text{NBC}(n)| = O(n^3)$ after sparsity.
\end{proof}

\begin{theorem}[Presentation independence; \ClaimStatusProvedHere]\label{thm:presentation-independence}
 The geometric bar complex satisfies:
 \begin{enumerate}
 \item \emph{Functoriality:} A morphism $\phi: \mathcal{A}_1 \to \mathcal{A}_2$ induces
 $\bar{B}^{\text{ch}}(\phi): \bar{B}^{\text{ch}}(\mathcal{A}_1) \to \bar{B}^{\text{ch}}(\mathcal{A}_2)$

 \item \emph{Quasi-isomorphism invariance:} If $\phi$ is a quasi-isomorphism, so is $\bar{B}^{\text{ch}}(\phi)$

 \item \emph{Presentation independence within equivalence class:} Two presentations
 $\mathcal{A} = \text{Free}^{\text{ch}}(V_1)/R_1 = \text{Free}^{\text{ch}}(V_2)/R_2$
 yield quasi-isomorphic bar complexes if and only if:
 \begin{itemize}
 \item Conformal weights are preserved modulo integers
 \item Relations differ only by Jacobi identity consequences
 \item Only tautological generators/relations are added/removed
 \end{itemize}

 \item \emph{Criticality obstruction:} Different weight assignments satisfying different criticality
 conditions yield non-quasi-isomorphic complexes
 \end{enumerate}
 \end{theorem}

 \begin{proof}
 \emph{(1) Functoriality.} The bar construction $\bar{B}^{\mathrm{ch}}$ is defined via the cofree conilpotent chiral coalgebra on the augmentation ideal, with differential induced by the chiral operations. A morphism $\phi\colon \mathcal{A}_1 \to \mathcal{A}_2$ preserves augmentation ideals and chiral operations, hence induces a coalgebra map $\bar{B}^{\mathrm{ch}}(\phi)$ by the universal property of the cofree coalgebra.

 \emph{(2) Quasi-isomorphism invariance.} This is a standard consequence of the fact that $\bar{B}^{\mathrm{ch}}$ is a derived functor: it preserves quasi-isomorphisms between augmented chiral algebras (see \cite[Theorem~2.2.4]{LV12} for the operadic statement, which applies verbatim in the chiral setting via the chiral operad formalism of Theorem~\ref{thm:geometric-equals-operadic-bar}).

 \emph{(3) Presentation independence.} Given two presentations $\mathcal{A} = \mathrm{Free}^{\mathrm{ch}}(V_1)/R_1 = \mathrm{Free}^{\mathrm{ch}}(V_2)/R_2$, the identity map $\mathrm{id}\colon \mathcal{A} \to \mathcal{A}$ induces an isomorphism $\bar{B}^{\mathrm{ch}}(\mathrm{id})$ by~(1). Conversely, if the three listed data (conformal weights, OPE structure, relations modulo Jacobi) differ, then the bar differentials, which are completely determined by these data, act differently on the generators, so the resulting complexes are not isomorphic. The ``only if'' direction uses the fact that the bar complex detects conformal weights via the residue pairing (Remark~\ref{rem:criticality-conformal-weight} below).

 \emph{(4) Criticality obstruction.} Follows from~(3): conformal weights enter the bar differential through the condition $h_i + h_j = h_k + 1$ for non-vanishing residue contributions. Different weight assignments change which matrix entries of the differential are non-zero, yielding non-quasi-isomorphic complexes.
 \end{proof}

 \begin{remark}[Criticality and conformal weight]\label{rem:criticality-conformal-weight}
 The criticality obstruction depends on conformal weights: the residue pairing
 contributes to the bar differential only when $h_i + h_j = h_k + 1$, so
 different presentations with different weight assignments yield different bar
 complexes. In particular, the geometric bar complex detects conformal
 dimensions, data that purely algebraic constructions may miss.
 \end{remark}

\begin{lemma}[Arnold relations on affine boundary screens; \ClaimStatusProvedHere]\label{lem:arnold-boundary}
On every affine genus-$0$ Fulton--MacPherson collision screen, the
Arnold relation extends as an identity of logarithmic forms to the
normal-crossing boundary. On $\mathbb{P}^1$ and on positive-genus
curves this is the local tangent-screen identity; the global
configuration-space relations are those of
Proposition~\ref{prop:arnold-higher-genus}.
\end{lemma}

\begin{proof}
Near a boundary stratum $D_I$ choose a holomorphic coordinate on the
curve and the FM coordinates of
Theorem~\ref{thm:local-coords-boundary}. If $i,j\in I$, then
$z_i-z_j=x_Iu_{ij}$ with $u_{ij}$ invertible on the chosen screen chart,
so
\[
\eta_{ij}=d\log x_I+d\log u_{ij}.
\]
For $i,j,k\in I$, the common scale term cancels from the cyclic sum and
the remaining identity is exactly the Arnold relation for the affine
screen coordinates $u_{ij},u_{jk},u_{ki}$.

If exactly two labels, say $i,j$, lie in $I$ and $k\notin I$, then
$z_i-z_k$ and $z_j-z_k$ have the same boundary value and
\[
\eta_{ik}-\eta_{jk}
=d\log\frac{z_i-z_k}{z_j-z_k}
\]
is regular and restricts to zero on $D_I$. The cyclic Arnold sum reduces
on the boundary to
$\eta_{ij}\wedge(\eta_{jk}-\eta_{ik})$, hence vanishes. The cases with
one or zero labels in $I$ are regular restrictions of the interior
identity.
\end{proof}

\subsection{Permutohedral tiling and cell complex}

\begin{theorem}[Permutohedral cell complex of the compactified real line; \ClaimStatusProvedHere]
\label{thm:permutohedral-cell-complex}
The open real configuration space $C_n(\mathbb R)$ is the disjoint union
of the $n!$ chambers indexed by total orders of the labels. After
quotienting by translations and positive dilations and taking the real
Fulton--MacPherson compactification, the compactified space
$\overline C_n(\mathbb R)/\mathrm{Aff}^+(\mathbb R)$ has a regular cell
decomposition whose face poset is the ordered-partition poset:
\begin{enumerate}
\item Cells $C_\pi$ correspond to ordered partitions $\pi = B_1 < B_2 < \cdots < B_k$ of $[n]$
\item The codimension of $C_\pi$ is $n-k$
\item The boundary relation is generated by refining or coarsening
adjacent collision blocks according to the ordered-partition incidence
orientation
\item The cellular cochain complex computes the cohomology of the
compactified quotient, equivalently the relative collision combinatorics
used by the real FM operad
\end{enumerate}
\end{theorem}
\begin{proof}
The open part corresponds to $k=n$, hence to singleton ordered
partitions and total orders
\[
x_{\sigma(1)}<x_{\sigma(2)}<\cdots<x_{\sigma(n)}.
\]
Boundary cells allow points in a block to collide at the same relative
scale. A cell indexed by
\[
\pi=B_1<\cdots<B_k
\]
has configuration type
\[
x_{B_1} < x_{B_2} < \cdots < x_{B_k}
\]
where $x_{B_i}$ denotes the common center of the block at that collision
scale, together with lower screens for the internal configurations in
the blocks. Partitioning $[n]$ into $k$ blocks imposes $n-k$ collision
conditions, so the cell has codimension $n-k$.

The incidence relation is the usual ordered-partition incidence
relation of the permutohedron: codimension increases by merging adjacent
blocks, while the dual face order is refinement. With the standard
orientation convention, the cellular differential
\[
\delta: C^{n-k}(\pi) \to \bigoplus_{\pi \to \pi'} C^{n-k+1}(\pi')
\]
is the signed sum over adjacent block mergers. This is the real
FM-operadic collision differential; it is not a cell decomposition of
the noncompact open space alone.
\end{proof}

\section{Computational complexity and algorithms}

\subsection{Complexity analysis}

\begin{remark}[Implementation]
While the theoretical bounds appear daunting,
the actual computation benefits from massive sparsity. In practice, most residues vanish
by weight or dimension considerations, reducing the effective complexity by several orders
of magnitude. For $n \leq 10$, computations are feasible on standard hardware.
\end{remark}

\begin{theorem}[Complexity bounds; \ClaimStatusProvedHere]
\label{thm:complexity-bounds}
For the geometric bar complex in dimension $n$:
\begin{enumerate}
\item NBC basis size: $|\mathrm{NBC}(n)| = n!$ (for the braid arrangement $\mathcal{A}_{n-1}$, matching $\dim H^*(C_n(\mathbb{C}))$)
\item Differential computation: $O(n^3)$ non-zero entries (by weight/sparsity)
\item Storage: $O(n \cdot n!)$ sparse representation
\item Computation of $d^2=0$: $O(n^5)$ operations
\end{enumerate}
\end{theorem}

\begin{proof}[Derivation]
\emph{NBC count.} For the braid arrangement, the NBC basis satisfies $f(n) = n \cdot f(n-1)$
with $f(1) = 1$, giving $|\mathrm{NBC}(n)| = n!$.

\emph{Differential.} Each NBC forest has $\leq n-1$ edges.
Computing residue per edge: $O(n)$ for weight matching.
Total per basis element: $O(n^2)$.
With $B(n)$ elements: a naive bound gives $O(n^2 \cdot B(n))$; sparsity reduces this to $O(n^3)$ nonzero entries.

\emph{Square-zero cost.} Compose the differential twice on $O(B(n))$
elements, each taking $O(n^3)$ operations.
\end{proof}

\begin{theorem}[Finite-window spectral sequence convergence]\label{thm:spectral-convergence}
Let $\cA$ have a finite-window residue datum
\textup{(}Definition~\textup{\ref{def:finite-window-residue-datum}}\textup{)}.
For each reduced-weight window $K_q(\cA)$ and each fixed genus range,
the spectral sequence of the filtered geometric bar complex
\[
E_1^{p,r}
=H^{p+r}\!\left(\operatorname{gr}_p K_q(\cA)\right)
\Longrightarrow
H^{p+r}\!\left(K_q(\cA)\right)
\]
is strongly convergent. If the window tower is Mittag--Leffler, the
completed spectral sequence converges to the cohomology of
\[
\widehat B_\rho(\cA)=\varprojlim_q K_q(\cA)
\]
in the completed/coderived sense appropriate to the curved model. No
strong-convergence statement is asserted for the unwindowed infinite bar
complex without this finite-window or completion hypothesis.
\end{theorem}

\begin{proof}
For fixed $q$, the window $K_q(\cA)$ is finite-dimensional by
Theorem~\ref{thm:reduced-weight-finiteness}; hence every total degree
contains only finitely many filtration summands, and the filtration is
bounded, complete, and Hausdorff. The usual spectral-sequence
convergence theorem applies to this finite filtered complex. The
Mittag--Leffler condition on the inverse system $\{K_q(\cA)\}_q$
ensures that inverse limits commute with cohomology, so the completed
spectral sequence computes the cohomology of
$\widehat B_\rho(\cA)$. In the curved positive-genus lane this is a
coderived statement: the raw fiberwise differential may square to
$\kappa(\cA)\omega_g$, while the finite-window argument controls the
completed curved object and its flat comparison differential.
\end{proof}

\subsubsection{Efficient residue computation}

\begin{construction}[Residue evaluation]\label{con:residue-evaluation}
Given fields $\phi_1(z_1) \otimes \cdots \otimes \phi_n(z_n) \otimes \omega$ with conformal weights $h_i$, the residue contributions are evaluated as follows.
For each collision divisor $D_{ij}$, write the finite singular part of
the OPE
\[
\phi_i(z_i)\phi_j(z_j)
=\sum_{m=0}^{P_{ij}}\sum_k
C^{k,m}_{ij}\,\phi_k(z_j)\,(z_i-z_j)^{-m-1}
+\text{regular}.
\]
If the logarithmic form factor contains the divisor factor
$\eta_{ij}=d\log(z_i-z_j)$, apply the Poincar\'e residue to that
logarithmic factor and the chiral product to the OPE coefficient:
replace $\phi_i\otimes\phi_j$ by $C^{k,m}_{ij}\phi_k$, remove one
$\eta_{ij}$ from~$\omega$, and include the Koszul sign coming from
moving the degree-$1$ residue operator through the remaining factors.
The simple-pole case $m=0$ is the Arnold/Jacobi projection; higher
$m$ are ordinary higher OPE modes and are not discarded.
\end{construction}


\begin{proposition}[Residue evaluation complexity; \ClaimStatusProvedHere]
\label{prop:residue-evaluation-complexity}
In a finite window with maximal OPE pole order $P_q$, the procedure of
Construction~\ref{con:residue-evaluation} computes residues with
complexity $O(n^2 P_q\cdot T_{\mathrm{OPE}})$, where
$T_{\mathrm{OPE}}$ is the time to look up one OPE coefficient. For
fixed $q$, $P_q$ is finite.
\end{proposition}

\begin{proof}
Correctness follows from the residue formula in
Theorem~\ref{thm:residue-operations}: the geometric residue extracts
the coefficient of the logarithmic divisor factor, while the chiral
product supplies the finite OPE coefficient. There are $\binom{n}{2}$
collision divisors and at most $P_q+1$ singular modes in the chosen
window, giving the stated bound.
\end{proof}

\section{Arnold relations: explicit computations and physical interpretation}
\label{sec:arnold-three-proofs-comprehensive}
\label{sec:arnold-chiral}

The Arnold relation (Theorem~\ref{thm:arnold-relations}) was proved
by two independent methods in \S\ref{sec:FM-compactification}: the
partial-fraction argument and the coordinate-reduction argument. Both
establish the relation in
$\Omega^2(C_n(\mathbb{C}))$ for every $n \geq 3$, and on every affine
formal collision screen of $\overline{C}_n(X)$. On this affine
genus-$0$ screen Theorem~\ref{thm:arnold-iff-nilpotent} shows that the
Arnold relation is equivalent to the triple-collision component of
$d_{\mathrm{residue}}^2 = 0$.

% Labels preserved for cross-referencing from other chapters.
\label{prop:arnold-topological}%
\label{sec:arnold-proof-topological}%
\label{prop:nilpotency-arnold-comprehensive}%

\subsection{\texorpdfstring{Explicit computations for $n = 3, 4, 5$}{Explicit computations for n = 3, 4, 5}}
\label{sec:arnold-explicit-n2345}

\begin{example}[\texorpdfstring{$n=3$}{n=3}: the fundamental relation]
\label{ex:arnold-n3-complete}

For $n=3$ points, the three generators $\eta_{12}, \eta_{23}, \eta_{13}$
are independent in $H^1(C_3(\mathbb{C}))$: the Arnold relation is
quadratic, not linear. The cohomology algebra is
$H^*(C_3(\mathbb{C})) = \bigwedge^*(\eta_{12}, \eta_{13}, \eta_{23})
/ (\text{Arnold relation})$, where the Arnold relation
$\eta_{12} \wedge \eta_{23} + \eta_{23} \wedge \eta_{31}
+ \eta_{31} \wedge \eta_{12} = 0$ is the explicit partial-fraction
identity of Theorem~\ref{thm:arnold-relations}.

\emph{Dimension count}:
\begin{align}
\dim H^0(C_3(\mathbb{C})) &= 1\\
\dim H^1(C_3(\mathbb{C})) &= 3 \quad \text{(}\eta_{12}, \eta_{13}, \eta_{23}\text{)}\\
\dim H^2(C_3(\mathbb{C})) &= 2 \quad \text{(two independent products modulo Arnold)}\\
\dim H^3(C_3(\mathbb{C})) &= 0
\end{align}

Total dimension: $1 + 3 + 2 = 6 = 3!$.

\emph{Poincaré polynomial}: $P_t(H^*(C_3(\mathbb{C}))) = (1+t)(1+2t) = 1 + 3t + 2t^2$.
\end{example}

\begin{example}[\texorpdfstring{$n=4$}{n=4}: multiple relations]
\label{ex:arnold-n4-complete}

For $n=4$ points, there are $\binom{4}{2} = 6$ forms:
\begin{equation}
\eta_{12}, \eta_{13}, \eta_{14}, \eta_{23}, \eta_{24}, \eta_{34}
\end{equation}

There are four independent Arnold relations, one for each choice of three points:
\begin{align}
\{1,2,3\}&: \eta_{12} \wedge \eta_{23} + \eta_{23} \wedge \eta_{31} +
\eta_{31} \wedge \eta_{12} = 0\\
\{1,2,4\}&: \eta_{12} \wedge \eta_{24} + \eta_{24} \wedge \eta_{41} +
\eta_{41} \wedge \eta_{12} = 0\\
\{1,3,4\}&: \eta_{13} \wedge \eta_{34} + \eta_{34} \wedge \eta_{41} +
\eta_{41} \wedge \eta_{13} = 0\\
\{2,3,4\}&: \eta_{23} \wedge \eta_{34} + \eta_{34} \wedge \eta_{42} +
\eta_{42} \wedge \eta_{23} = 0
\end{align}

These relations are independent.

\emph{Dimension count}:
\begin{align}
\dim H^0(C_4(\mathbb{C})) &= 1\\
\dim H^1(C_4(\mathbb{C})) &= 6 \quad \text{(all $\binom{4}{2}$ generators independent)}\\
\dim H^2(C_4(\mathbb{C})) &= 11 \quad \text{($\binom{6}{2} - 4$ Arnold relations)}\\
\dim H^3(C_4(\mathbb{C})) &= 6\\
\dim H^k(C_4(\mathbb{C})) &= 0 \quad k \geq 4
\end{align}

Total dimension: $1 + 6 + 11 + 6 = 24 = 4!$.

\emph{Poincaré polynomial}: $P_t(H^*(C_4(\mathbb{C}))) = (1+t)(1+2t)(1+3t) = 1 + 6t + 11t^2 + 6t^3$.

Since $C_4(\mathbb{C})$ is non-compact, Poincar\'{e} duality does not apply; the polynomial is not palindromic.
\end{example}

\begin{example}[\texorpdfstring{$n=5$}{n=5}: complete computation]
\label{ex:arnold-n5-complete}

For $n=5$ points, there are $\binom{5}{2} = 10$ forms and $\binom{5}{3} = 10$ Arnold
relations.

The cohomology dimensions are given by the Poincar\'e polynomial
$\prod_{k=1}^{4}(1+kt) = (1+t)(1+2t)(1+3t)(1+4t)$:
\begin{align}
\dim H^0(C_5(\mathbb{C})) &= 1\\
\dim H^1(C_5(\mathbb{C})) &= 10 \quad \text{(all $\binom{5}{2}$ generators)}\\
\dim H^2(C_5(\mathbb{C})) &= 35\\
\dim H^3(C_5(\mathbb{C})) &= 50\\
\dim H^4(C_5(\mathbb{C})) &= 24\\
\dim H^k(C_5(\mathbb{C})) &= 0 \quad k \geq 5
\end{align}

Total dimension: $1 + 10 + 35 + 50 + 24 = 120 = 5!$.

In general, $\sum_k \dim H^k(C_n(\mathbb{C})) = n!$ and the Poincar\'e polynomial
is $\prod_{k=1}^{n-1}(1+kt)$, as proved by Arnol'd~\cite{Arnold69}.
\end{example}

\subsection{Physical interpretation: Jacobi identity and associativity}
\label{sec:arnold-physical}

\begin{theorem}[Arnold relation $\Leftrightarrow$ simple-pole Jacobi on the affine screen \cite{LV12}; \ClaimStatusProvedElsewhere]
\label{thm:arnold-jacobi}
On an affine genus-$0$ collision screen, the Arnold cyclic relation in
$\Omega^2(\log D)$ is equivalent to the Jacobi identity for the chiral
bracket $\{a,b\} := a_{(0)}b$; equivalently, to the simple-pole
associativity axiom of the chiral operad (Loday--Vallette
\cite[Theorem~7.1.3, \S7.6]{LV12}). This is the simple-pole projection;
the full bar differential requires the Borcherds identity, which
couples all OPE pole orders (Proposition~\ref{prop:pole-decomposition}
in Chapter~\ref{chap:bar-cobar}), and the global curve cases require
the projective or Totaro--period corrections of
Proposition~\ref{prop:arnold-higher-genus}.
\end{theorem}

\begin{remark}[Physical mnemonic: simple-pole crossing]
\label{rem:arnold-crossing-symmetry}
On an affine genus-$0$ collision screen, the Arnold relation expresses
crossing symmetry of OPE at the simple-pole level: the three iterated
residues at the triple-collision stratum $D_{ijk}$ correspond to the
three ways of expanding $\phi_i(z_i)\phi_j(z_j)\phi_k(z_k)$ via
successive collisions, and Arnold is the identity
\begin{equation}
\operatorname{Res}_{D_{jk}}\operatorname{Res}_{D_{ij}}
+ \operatorname{Res}_{D_{ki}}\operatorname{Res}_{D_{jk}}
+ \operatorname{Res}_{D_{ij}}\operatorname{Res}_{D_{ki}} = 0
\end{equation}
forcing the three expansions to agree on the chiral bracket. For full
OPEs the corresponding statement is the Borcherds identity, and in
positive genus it is coupled to the period corrections already named.
\end{remark}

\begin{corollary}[Operadic associativity \cite{LV12}; \ClaimStatusProvedElsewhere]
\label{cor:arnold-operadic}
On the affine genus-$0$ simple-pole part, the Arnold relations are
equivalent to the associativity axiom for the chiral operad:
\begin{equation}
\gamma \circ (\text{id} \otimes \gamma) = \gamma \circ (\gamma \otimes \text{id})
\end{equation}
where $\gamma$ is the operadic composition.
\end{corollary}

\begin{theorem}[Arnold--Orlik--Solomon circuit relations; \ClaimStatusProvedHere]
\label{thm:arnold-orlik-solomon}
Let $i_0,i_1,\ldots,i_m,i_0$ be a simple cycle in the complete graph on
$\{1,\ldots,n\}$, with $m\geq 2$, and put
$e_r=\eta_{i_r i_{r+1}}$ \textup{(}indices modulo $m+1$\textup{)}.
In $H^*(C_n(\mathbb{C});\mathbb{C})$ one has the Orlik--Solomon
circuit relation
\[
\sum_{r=0}^{m}
(-1)^r\,
e_0\wedge\cdots\wedge\widehat{e_r}\wedge\cdots\wedge e_m
=0.
\]
For $m=2$ this is Arnold's three-term relation, after translating
between the fixed cyclic order and the symmetric convention
$\eta_{ij}=\eta_{ji}$. For $m>2$ these relations are consequences of
the triangular Arnold relations, since every cycle in a complete graph
has a chord.
\end{theorem}

\begin{proof}
The braid arrangement is the graphic arrangement of the complete graph.
The Orlik--Solomon ideal is generated by boundaries of graphic
circuits. A circuit is exactly a simple cycle, and its boundary is the
displayed alternating sum. For a triangle this boundary is the
three-term Arnold relation. If the cycle has length at least four,
choose a chord. The chord splits the cycle into two shorter cycles; the
Orlik--Solomon boundary of the long cycle is the signed sum of the two
shorter circuit boundaries after wedging with the complementary edge
monomials. Induction on the length of the cycle reduces every circuit
relation to triangular Arnold relations.
\end{proof}

\begin{corollary}[Bar differential squares to zero; \ClaimStatusProvedHere]
\label{cor:bar-d-squared-zero}
On an affine genus-$0$ screen, the Arnold relations ensure $d^2 = 0$
for the bar differential. For $X=\mathbb{P}^1$ or $g(X)\geq 1$, one
must add the projective or period relations of
Proposition~\ref{prop:arnold-higher-genus}.
\end{corollary}
\begin{proof}
This is Theorem~\ref{thm:bar-nilpotency-complete}.
\end{proof}

\section{Higher genus}

\subsection{Genus 1: elliptic functions}

On a torus $E_\tau = \mathbb{C}/(\mathbb{Z} + \tau\mathbb{Z})$:

\begin{theorem}[Elliptic logarithmic forms \cite{Fay73}; \ClaimStatusProvedElsewhere]
\label{thm:elliptic-logarithmic-forms}
The logarithmic form becomes:
\[\eta_{ij}^{(1)} = d\log\vartheta_1\left(\frac{z_i - z_j}{2\pi i}\Big|\tau\right) + \text{modular correction}\]
where $\vartheta_1(z|\tau)$ is the odd Jacobi theta function.
\end{theorem}

\subsection{Higher genus: prime forms}

\begin{definition}[Prime form]
\index{prime form|textbf}
On a Riemann surface of genus $g \geq 1$, the prime form $E(z,w)$ is the unique section of $K^{-1/2} \boxtimes K^{-1/2}$ (i.e., a $(-1/2, -1/2)$-differential) with a simple zero at $z = w$, no other zeros, and appropriate normalization along cycles.
\end{definition}

The logarithmic forms at higher genus are built from prime forms and period integrals.

\section{Normal crossings at higher genus}
\label{sec:normal-crossings-higher-genus}

\begin{theorem}[Normal crossings preservation; \ClaimStatusProvedHere]\label{thm:normal-crossings-preservation}
Let $\pi\colon \mathcal X\to S$ be a smooth family of curves over a
smooth base with normal-crossing boundary $B\subset S$, and assume
$\pi$ extends log-smoothly over $B$. The relative
Fulton--MacPherson compactification
$\overline{C}_n(\mathcal X/S)$ has logarithmic normal-crossing boundary
given by the union of the relative collision divisors and the pullback
of~$B$. For a fixed curve $X$ this reduces to the ordinary FM
normal-crossing boundary of \cite{FM94}. The Deligne--Mumford boundary
of $\overline{\mathcal M}_{g,n}$ enters only through such a relative
family; it is not a boundary component of $\overline C_n(X)$ for a
fixed smooth curve.
\end{theorem}

\begin{proof}
For a fixed curve this is Theorem~\ref{thm:normal-crossings}. In the
relative case, choose local coordinates
$(t_1,\ldots,t_r,s_1,\ldots,s_a)$ on the base with
$B=\{t_1\cdots t_r=0\}$ and a relative coordinate $z$ on the curve.
The log-smooth hypothesis gives nodal or smooth local models in which
the base boundary is monomial in the $t_\alpha$. The relative FM
construction blows up only the relative diagonals
$z_i=z_j$ inside the fibres. In the FM charts, each collision supplies a
scale coordinate $x_S$, and the boundary is locally
\[
t_1\cdots t_r\prod_S x_S=0.
\]
Thus the collision and base-degeneration divisors form a logarithmic
normal-crossing divisor. The residue maps used in the bar differential
are therefore defined along collision faces for fixed curves and along
the additional log base faces only in the relative modular family.
\end{proof}

\begin{remark}[Fixed curve versus relative modular family]\label{rem:iterated-blowup}
The expression
$\overline{\mathcal M}_{g,n}\times_{X^n}\overline C_n(X)$ is not an
intrinsic object for a fixed curve $X$: $\overline{\mathcal M}_{g,n}$
varies the curve, while $X^n$ does not. The correct object is the
relative FM compactification of the universal stable curve, or a
log-smooth family pulled back from it. In that ambient, the iterated
blow-ups are fibrewise along relative diagonals, and normal crossings
are preserved by the monomial local form above.
\end{remark}
\section{\texorpdfstring{Explicit local coordinates on $\overline{C}_n(X)$}{Explicit local coordinates on C n(X)}}
\label{sec:explicit-coordinates-complete}

\subsection{Coordinate systems near collision boundaries}

\begin{convention}[Coordinate notation]\label{conv:coordinate-notation}
For a boundary divisor $D_S \subset \partial\overline{C}_n(X)$ where
the points in
$S = \{i_1, \ldots, i_k\} \subseteq \{1,\ldots,n\}$ collide, we use
holomorphic Fulton--MacPherson coordinates. Choose a local coordinate
$t$ on $X$ near the collision point and an affine chart in the
projectivized normal bundle to the partial diagonal. Then the colliding
coordinates have the form
\[
t_i=u_S+x_S q_i,\qquad i\in S,
\]
where the $q_i$ are normalized relative coordinates, for example by
$\sum_{i\in S}q_i=0$ and by setting one non-zero difference equal to
$1$ in the chosen chart. The coordinates are:

\begin{enumerate}
\item \emph{Center coordinate:} $u_S \in X$ (the limiting position where points collide)

\item \emph{Normal coordinate:} $x_S$ (a holomorphic equation for the
divisor $D_S$ in the chosen chart)

\item \emph{Screen coordinates:} $q_S=(q_i)_{i\in S}$, a point of the
projectivized relative normal space
$\mathbb{P}(N_{\Delta_S/X^n})\cong\mathbb{P}^{k-2}$ on the open
single-collision stratum, and a point of the compactified screen
$\overline{C}_k(T_{u_S}X)/\operatorname{Aff}(\mathbb{A}^1)$ on its
boundary

\item \emph{External coordinates:} $z_j$ for $j \notin S$ (positions of non-colliding points)
\end{enumerate}

The boundary equation is $D_S=\{x_S=0\}$. A real oriented blow-up chart
may write $x_S=r_Se^{i\vartheta_S}$, but $(r_S,\vartheta_S)$ are not
algebraic coordinates on the complex FM compactification. In
particular, for $|S|=2$ the screen is $\mathbb P^0$, so there is no
additional projective direction attached to a pair collision on a
curve.

\emph{Dimensions.}
\begin{align*}
\dim(u_S) &= \dim X = 1 \quad \text{(complex)}\\
\dim(x_S) &= 1 \quad \text{(normal coordinate to }D_S)\\
\dim(q_S) &= k-2 \quad \text{(screen coordinates)}\\
\dim(z_j) &= (n-k) \cdot 1 \quad \text{(other points)}
\end{align*}

Total: $1 + 1 + (k-2) + (n-k) = n$ (complex), as expected for $\overline{C}_n(X)$.
\end{convention}

\subsection{\texorpdfstring{Pair collisions ($n=2$)}{Pair collisions (n=2)}}

\begin{example}[Coordinates on \texorpdfstring{$\overline{C}_2(X)$}{C-bar-2(X)}]\label{ex:coords-n2}
For two points on a curve $X$, the configuration space is:
\[C_2(X) = \{(z_1, z_2) \in X \times X : z_1 \neq z_2\}\]

The compactification $\overline{C}_2(X)$ is $X^2$ with the diagonal
$D_{12}=\Delta$ remembered as the boundary divisor. Equivalently, the
FM blow-up along $\Delta$ is trivial because $\Delta\subset X^2$ is
already a Cartier divisor.

\subsubsection{Interior coordinates}

On the open part $C_2(X)$, natural coordinates are simply:
\[(z_1, z_2) \in X \times X\]

\subsubsection{Logarithmic boundary coordinates}

Choose a local coordinate $t$ on $X$ and put:
\begin{enumerate}
\item $u=\frac{t_1+t_2}{2}$ (center),
\item $x=t_1-t_2$ (normal coordinate to the diagonal).
\end{enumerate}

\emph{Relationship to original coordinates.}
\begin{align}
t_1 &= u + \frac{x}{2} \label{eq:blowup-z1}\\
t_2 &= u - \frac{x}{2} \label{eq:blowup-z2}
\end{align}

\emph{Boundary divisor.} $D_{12}=\{x=0\}\cong X$.

There is no exceptional $\mathbb P^1$ or angular parameter in the
complex algebraic compactification for two points on a curve. The
normal line is the difference tangent line $T_X$ along the diagonal,
and the logarithmic form used by the bar differential is $d\log x$.
\end{example}

\begin{remark}[OPE encoding]
The logarithmic boundary coordinate is the OPE variable:
\[
\phi_1(t_1)\phi_2(t_2)
=
\sum_{m\geq 0}\frac{(\phi_1{}_{(m)}\phi_2)(u)}{x^{m+1}}
+\text{regular},\qquad x=t_1-t_2 .
\]

Identifying:
\begin{itemize}
\item $x=t_1-t_2$: the holomorphic separation variable in the pole,
\item $d\log x$: the logarithmic propagator whose Poincar\'e residue
reads the collision coefficient,
\item $u=\frac{t_1+t_2}{2}$: the center where the composite operator is inserted.
\end{itemize}

For higher collisions the compactified screens record the possible
nested OPE channels.
\end{remark}

\subsection{\texorpdfstring{Three-point collisions ($n=3$)}{Three-point collisions (n=3)}}

\begin{example}[Coordinates on \texorpdfstring{$\overline{C}_3(X)$}{C-bar-3(X)}]\label{ex:coords-n3}
For three points on a curve, the boundary divisors are indexed by all
subsets of cardinality at least two:
\[
D_{12},\quad D_{23},\quad D_{13},\quad D_{123}
\subset \partial\overline{C}_3(X).
\]

Each of these is codimension one. Codimension-two corners are the
nested intersections $D_{123}\cap D_{12}$,
$D_{123}\cap D_{13}$, and $D_{123}\cap D_{23}$.

\subsubsection{\texorpdfstring{Coordinates near $D_{12}$ (first two points collide)}{Coordinates near D 12 (first two points collide)}}

\emph{Logarithmic coordinates for the $(z_1,z_2)$ collision.}
\begin{align}
u_{12} &= \frac{t_1 + t_2}{2} \quad \text{(center of 1 and 2)}\\
x_{12} &= t_1 - t_2 \quad \text{(normal coordinate)}\\
z_3 &= z_3 \quad \text{(third point unchanged)}
\end{align}

\emph{Inverse relations.}
\begin{align}
t_1 &= u_{12} + \frac{x_{12}}{2}\\
t_2 &= u_{12} - \frac{x_{12}}{2}\\
z_3 &= z_3
\end{align}

\emph{Domain.} $u_{12},z_3\in X$ in the chosen coordinate chart, and
$x_{12}$ is a holomorphic normal coordinate.

\emph{Divisor.} $D_{12}=\{x_{12}=0\}$; on the open part where the
third point is distinct from the collision point this is
$C_2(X)$ with coordinates $(u_{12},z_3)$.
\begin{itemize}
\item First $X$: location $u_{12}$ where 1,2 collide
\item Second $X$: location $z_3$ of third point
\end{itemize}
(For a curve $X$, the normal bundle to the pair diagonal is
$1$-dimensional, so the projectivized screen is $\mathbb{P}^0$ and
contributes no additional moduli.)

\subsubsection{The triple-collision divisor and nested corners}

After the FM blow-up of the triple diagonal, the proper transforms of
the pairwise diagonals $D_{12}$ and $D_{23}$ no longer intersect each
other. The all-three collision is the divisor $D_{123}$, and the
codimension-two strata remember a pair collision inside the
triple-collision screen:
\[
D_{123}\cap D_{12},\qquad D_{123}\cap D_{13},\qquad
D_{123}\cap D_{23}.
\]

\emph{Coordinates near the open triple divisor.}

Choose normalized screen coordinates $q_1,q_2,q_3$ with
$q_1+q_2+q_3=0$ and one non-zero difference fixed to $1$. Then
\begin{align}
t_i &= u+x_{123}q_i,\qquad i=1,2,3,\\
D_{123}&=\{x_{123}=0\}.
\end{align}
The screen coordinate lies in
$\overline{C}_3(T_uX)/\operatorname{Aff}(\mathbb{A}^1)$, a
one-dimensional compactified configuration space.

\emph{Example: the nested corner $D_{123}\cap D_{12}$.}
In the screen, use a second normal coordinate for the pair collision
\begin{align}
q_1-q_2 &= x_{12}v_{12},\\
q_1+q_2+q_3&=0,
\end{align}
so that $D_{123}=\{x_{123}=0\}$ and $D_{12}=\{x_{12}=0\}$ are
normal-crossing coordinate hyperplanes. The ratio of the original
separation $t_1-t_2$ to the overall scale $x_{123}$ is now encoded by
the screen coordinate and its boundary value.
\end{example}

\subsection{\texorpdfstring{$n$-point collisions for arbitrary $n$}{n-point collisions for arbitrary n}}

\begin{theorem}[Complete coordinate description; \ClaimStatusProvedHere]\label{thm:complete-coordinates}
Let $\mathcal N=\{S_1,\ldots,S_r\}$ be a nest of subsets of
$\{1,\ldots,n\}$, with $|S_a|\geq2$ and any two subsets either disjoint
or comparable by inclusion. Near the boundary stratum
\[
D_{\mathcal N}:=\bigcap_{a=1}^r D_{S_a}\subset \overline{C}_n(X),
\]
there are holomorphic local coordinates
\[
\bigl(\{x_{S_a}\}_{S_a\in\mathcal N},
\{u_B\}_{B\ \mathrm{external\ or\ cluster}},
\{q_{S_a}\}_{S_a\in\mathcal N}\bigr)
\]
such that
\[
D_{S_a}=\{x_{S_a}=0\},\qquad
D_{\mathcal N}=\{x_{S_1}=\cdots=x_{S_r}=0\}.
\]
For a single divisor $D_S$, on the open single-collision stratum this
reduces to
\[
t_i=u_S+x_Sq_i,\qquad i\in S,
\]
with $q_S=(q_i)$ a point of
$\mathbb P(N_{\Delta_S/X^n})\cong\mathbb P^{|S|-2}$; on deeper
boundary pieces the screen is replaced by
$\overline C_{|S|}(T_{u_S}X)/\operatorname{Aff}(\mathbb{A}^1)$.
Thus every boundary divisor has codimension one, while the codimension
of a stratum is the number of divisors in its nest.
\end{theorem}

\begin{proof}[Construction]
Start from $X^n$ and the building set of partial diagonals
\[
\Delta_S=\{(z_1,\ldots,z_n):z_i=z_j\ \text{for all }i,j\in S\},
\qquad |S|\geq2.
\]
For a curve, $\Delta_S$ has codimension $|S|-1$. The
Fulton--MacPherson compactification is obtained by blowing up the
proper transforms of these diagonals in increasing dimension of the
diagonal, equivalently in decreasing codimension of the normal screen
after larger collisions have been separated. The standard wonderful
compactification chart attached to a nest gives one normal coordinate
$x_S$ for each chosen divisor and one projective screen coordinate for
each collision. In such a chart the total transform of the boundary is
the union of the coordinate hyperplanes $\{x_S=0\}$. The displayed
formula $t_i=u_S+x_Sq_i$ is the single-divisor chart in a local
coordinate on $X$; iterating it down a nest gives the asserted
coordinates and the normal-crossing equations.
\end{proof}

\subsection{Normal bundle calculations}

\begin{theorem}[Normal bundle formula; \ClaimStatusProvedHere]\label{thm:normal-bundle-formula}
For every boundary divisor $D_S\subset\overline C_n(X)$,
\[
N_{D_S/\overline C_n(X)}
\cong
\mathcal O_{\overline C_n(X)}(D_S)\big|_{D_S}.
\]
In a nested FM chart the normal direction is represented by
$\partial/\partial x_S$, where $D_S=\{x_S=0\}$. If $D_S$ is the
exceptional divisor created by blowing up a partial diagonal of
codimension at least two, then on the open part before later boundary
blow-ups this normal line is the tautological
$\mathcal O_{\mathbb P(N_{\Delta_S/X^n})}(-1)$. For pair diagonals on a
curve, $\Delta_{ij}\subset X^n$ is already a divisor; there is no
exceptional projective fiber and no universal $\mathcal O(-1)$
statement.
\end{theorem}

\begin{proof}[Geometric proof]
The first displayed is the standard identity
$N_{D/Y}\cong\mathcal O_Y(D)|_D$ for a smooth divisor $D$ in a smooth
variety $Y$, applied to
$Y=\overline C_n(X)$ and $D=D_S$. In a local equation $x_S=0$, the
class of $\partial/\partial x_S$ frames the normal line.

When a smooth center $Z\subset Y$ of codimension at least two is blown
up, the exceptional divisor $E=\mathbb P(N_{Z/Y})$ satisfies
$N_{E/\operatorname{Bl}_Z Y}\cong
\mathcal O_{\mathbb P(N_{Z/Y})}(-1)$. This gives the tautological
description on the open part of an exceptional FM divisor. If
$|S|=2$ and $\dim X=1$, the center is Cartier, so the blow-up is
locally the identity and the normal line is the ordinary diagonal normal
line rather than a tautological exceptional bundle. The later nested
blow-ups restrict the divisor by normal-crossing charts, preserving the
local equation description.
\end{proof}

\begin{example}[Normal bundle for \texorpdfstring{$D_{12}$}{D12} in \texorpdfstring{$\overline{C}_3(\mathbb{C})$}{C-bar-3(C)}]\label{ex:normal-bundle-d12}
Consider $D_{12} \subset \overline{C}_3(\mathbb{C})$ where the first two points collide.

\emph{Divisor structure.}
\[
D_{12}^{\circ}\cong
\{(u_{12},z_3)\in\mathbb C^2:u_{12}\neq z_3\}.
\]
Here:
\begin{itemize}
\item First $\mathbb{C}$: location $u_{12}$ of collision
\item Second $\mathbb{C}$: location $z_3$ of third point
\end{itemize}
Its closure meets $D_{123}$ along the nested corner
$D_{123}\cap D_{12}$.

\emph{Normal direction.} $\partial/\partial x_{12}$, where
$x_{12}=z_1-z_2$.

In coordinates, a section of $N_{D_{12}}$ looks like:
\[
s(u_{12},z_3)=f(u_{12},z_3)\cdot
\frac{\partial}{\partial x_{12}}.
\]

Because $T_{\mathbb C}$ is trivial, this normal line is trivial on the
open stratum. It is not the exceptional $\mathcal O(-1)$ bundle of a
projectivized normal space.
\end{example}

\subsection{Transition functions between charts}

\begin{proposition}[Transition functions \cite{FM94}; \ClaimStatusProvedElsewhere]\label{prop:transition-functions}
Consider two overlapping coordinate charts on $\overline{C}_n(X)$:
\begin{itemize}
\item Chart $U_1$: coordinates $(u_S,x_S,q_S,\ldots)$
\item Chart $U_2$: coordinates $(u_T,x_T,q_T,\ldots)$
\end{itemize}

where $S, T \subseteq \{1,\ldots,n\}$ are different subsets.

\emph{Case 1: $S \cap T = \emptyset$ (disjoint collisions)}

The charts are independent:
\begin{align}
u_S &= u_S(u_T,x_T,q_T,\ldots),&
x_S&=x_S(\ldots),\\
u_T &= u_T(u_S,x_S,q_S,\ldots),&
x_T&=x_T(\ldots).
\end{align}

\emph{Case 2: $S \subset T$ (nested collisions)}

Points in $S$ collide first (small scale), then the cluster collides with other points
in $T$ (larger scale).

Transition functions are monomial in the normal coordinates, up to
units:
\begin{align}
u_T &= \frac{1}{|T|} \sum_{i \in T} z_i = \frac{1}{|T|}\left( |S| u_S + \sum_{i \in T \setminus S} z_i \right)\\
x_S&=x_T\,\widetilde x_S\cdot \text{(unit)}
\end{align}
where $\widetilde x_S$ is the normal coordinate for the collision of
the $S$-points inside the $T$-screen.

\emph{Case 3: $S \cap T \neq \emptyset$, $S \not\subset T$, $T \not\subset S$ (overlapping).}

The divisors $D_S$ and $D_T$ have no common boundary stratum. Their
coordinate descriptions overlap only on the interior locus where
$x_S\neq0$ and $x_T\neq0$; near the simultaneous limiting collision one
passes through the divisor $D_{S\cup T}$ and its nested charts.
\end{proposition}

\begin{example}[Transition for \texorpdfstring{$D_{12}$}{D12} and \texorpdfstring{$D_{23}$}{D23} in \texorpdfstring{$\overline{C}_3(\mathbb{C})$}{C-bar-3(C)}]\label{ex:transition-d12-d23}
The divisors $D_{12}$ and $D_{23}$ do not meet in
$\overline C_3(\mathbb C)$, because the subsets $\{1,2\}$ and
$\{2,3\}$ overlap without being nested. Their charts overlap only where
neither boundary equation vanishes.

\emph{Chart $U_{12}$:} $(u_{12},x_{12},z_3)$
\[
z_1=u_{12}+\frac{x_{12}}2,\qquad
z_2=u_{12}-\frac{x_{12}}2,\qquad z_3=z_3.
\]

\emph{Chart $U_{23}$:} $(u_{23},x_{23},z_1)$
\[
z_2=u_{23}+\frac{x_{23}}2,\qquad
z_3=u_{23}-\frac{x_{23}}2,\qquad z_1=z_1.
\]

\emph{Transition functions.}
To express $U_{23}$ coordinates in terms of $U_{12}$:
\begin{align}
u_{23} &= \frac{z_2+z_3}{2}
=\frac{u_{12}-x_{12}/2+z_3}{2},\\
x_{23} &= z_2-z_3=u_{12}-\frac{x_{12}}2-z_3.
\end{align}

These are holomorphic functions on the interior overlap
$x_{12}\neq0$, $x_{23}\neq0$. As all three points collide, the
appropriate boundary chart is the nested chart of $D_{123}$ together
with one of $D_{12}$, $D_{13}$, or $D_{23}$.
\end{example}

\subsection{Normal crossings}

The boundary divisors intersect in \emph{normal crossings}: locally
like coordinate hyperplanes.

\begin{theorem}[Normal crossings property; \ClaimStatusProvedHere]\label{thm:normal-crossings-verified}
Let $D_{S_1}, \ldots, D_{S_k}$ be boundary divisors of
$\overline{C}_n(X)$ that intersect. Equivalently, the subsets
$S_i$ form a nest: any two are either disjoint or one contains the
other.
Then their intersection $D_{S_1} \cap \cdots \cap D_{S_k}$ has normal crossings, meaning:

There exist local coordinates $(x_1, \ldots, x_d)$ near any point in the intersection such that:
\[D_{S_i} = \{x_i = 0\}\]
for $i = 1, \ldots, k$.
\end{theorem}

\begin{proof}[The case $n=3$]
For $\overline{C}_3(\mathbb C)$ the divisors $D_{12}$ and $D_{23}$ do
not meet as simultaneous codimension-one collision divisors: the
subsets $\{1,2\}$ and $\{2,3\}$ overlap without being nested. The
codimension-two corner over the triple collision is instead represented
by a nested pair, for example $D_{123}\cap D_{12}$.

Use the holomorphic nested blow-up coordinates
\[
 z_1-z_3=x_{123}u_1,\qquad
 z_2-z_3=x_{123}u_2,\qquad
 u_1-u_2=x_{12}v_{12},
\]
on the chart where the first screen records the triple collision and
the second screen records the collision of the labelled pair inside
that screen. Then
\[
D_{123}=\{x_{123}=0\},\qquad D_{12}=\{x_{12}=0\}.
\]
These are coordinate hyperplanes. The other nested charts are obtained
by permuting labels. The general case is the same building-set
argument: Fulton--MacPherson boundary divisors meet only for nested or
disjoint subsets, and the corresponding blow-up scales are independent
normal coordinates.
\end{proof}

\begin{remark}[Normal crossings]
For general $n$, the normal crossings property is guaranteed by the Fulton--MacPherson
construction itself. The iterated blow-ups are designed precisely to achieve this.

Blowing up in order of decreasing codimension ensures that later blow-ups
are transverse to earlier exceptional divisors, preserving normal crossings at each stage.
\end{remark}

\subsection{\texorpdfstring{Example: $n=4$ with all coordinates}{Example: n=4 with all coordinates}}

\begin{example}[Complete coordinate atlas for \texorpdfstring{$\overline{C}_4(\mathbb{C})$}{C-bar-4(C)}]\label{ex:coords-n4-complete}
For four points, there are $2^4-4-1=11$ boundary divisors, not eleven
boundary strata:
\begin{itemize}
\item six pair divisors $D_{ij}$;
\item four triple divisors $D_{ijk}$;
\item one quadruple divisor $D_{1234}$.
\end{itemize}
Higher-codimension strata are intersections of compatible divisors,
equivalently nests. For example $D_{123}\cap D_{12}$ and
$D_{12}\cap D_{34}$ have codimension two, while
$D_{1234}\cap D_{123}\cap D_{12}$ has codimension three.

\subsubsection{Pair-divisor coordinates}

For each $D_{ij}$:
\begin{align}
u_{ij} &= \frac{z_i + z_j}{2},\\
x_{ij} &= z_i-z_j,\\
z_k, z_\ell &= z_k, z_\ell \quad (k, \ell \notin \{i,j\})
\end{align}

\emph{Example: $D_{12}$.}
\[\text{Coords: } (u_{12},x_{12},z_3,z_4)\]
\[\text{Divisor: } D_{12}=\{x_{12}=0\}\]
(The collision locus $u_{12} \in \mathbb{C}$ times the two free points $z_3, z_4 \in \mathbb{C}^2$. For a curve, the fiber $\mathbb{P}^{|S|-2} = \mathbb{P}^0$ is a point and contributes no factor.)

\subsubsection{Triple-divisor and nested-corner coordinates}

For the divisor $D_{123}$:
\begin{align}
z_i&=u_{123}+x_{123}q_i,\qquad i=1,2,3,\\
q_1+q_2+q_3&=0,\qquad q_S\in
\overline C_3(\mathbb{A}^1)/\operatorname{Aff}(\mathbb{A}^1),\\
z_4 &= z_4
\end{align}
and $D_{123}=\{x_{123}=0\}$. For the nested codimension-two corner
$D_{123}\cap D_{12}$, add a screen coordinate
$q_1-q_2=x_{12}v_{12}$, so the two boundary equations are
$x_{123}=0$ and $x_{12}=0$.

\subsubsection{Quadruple divisor and a deepest nest}

The all-four collision $D_{1234}$ is also a divisor:
\begin{align}
z_i&=u+x_{1234}q_i,\qquad i=1,\ldots,4,\\
q&\in\overline C_4(\mathbb{A}^1)/\operatorname{Aff}(\mathbb{A}^1).
\end{align}
One maximal nest is $\{1234,123,12\}$, with coordinates
\begin{align}
x_{1234}&=0,\\
x_{123}&=0,\\
x_{12}&=0.
\end{align}
This codimension-three stratum records the tree in which $1$ and $2$
collide inside the $123$ cluster, and the $123$ cluster collides inside
the $1234$ cluster.
\end{example}

\subsection{Summary table: coordinate systems}

\begin{table}[ht]
\centering
\caption{Coordinate systems on $\overline{C}_n(X)$}
\label{tab:coordinate-summary}
\begin{tabular}{|l|l|l|l|}
\hline
\textbf{Stratum} & \textbf{Codim} & \textbf{Coordinates} & \textbf{Divisor Equation} \\
\hline
Interior $C_n(X)$ & 0 & $(z_1, \ldots, z_n)$ & none \\
\hline
$D_{ij}$ & 1 & $(u_{ij},x_{ij},\{z_k\}_{k\neq i,j})$ & $x_{ij}=0$ \\
\hline
$D_{ijk}$ & 1 & $(u_{ijk},x_{ijk},q_{ijk},\{z_\ell\})$ & $x_{ijk}=0$ \\
\hline
Nested set $\mathcal N$ & $|\mathcal N|$ & $(\{x_S\},\{q_S\},\{u_B\})$ & $x_S=0$ for $S\in\mathcal N$ \\
\hline
\end{tabular}
\end{table}

\subsection{Connection to chiral algebra and OPE}

\begin{remark}[Geometric encoding of OPE]\label{rem:coordinates-ope}
The Fulton--MacPherson coordinates mirror the structure of the operator
product expansion:

\begin{center}
\begin{tabular}{c|c}
Geometry & CFT \\
\hline
$x_{ij}=z_i-z_j$ & Holomorphic OPE variable \\
$d\log x_{ij}$ & Residue kernel in the bar differential \\
$u$ (center) & Location of composite operator \\
$q_S$ screen & Nested OPE channel \\
Normal crossings & Consistency of nested OPEs \\
Nested blow-ups & Associativity of OPE
\end{tabular}
\end{center}

For example, the OPE
\[
\phi_1(z_1)\phi_2(z_2)
=\sum_{m\geq0}\frac{(\phi_1{}_{(m)}\phi_2)(u)}{x_{12}^{m+1}}
+\text{regular}
\]
corresponds to the coordinate transformation
$(z_1,z_2)\mapsto (u=(z_1+z_2)/2,x_{12}=z_1-z_2)$.
Pair collisions on a curve have no projective screen; higher screens
separate the possible nested OPE channels for three or more colliding
points.
\end{remark}

\section{Ran space: topological and geometric structure}
\label{sec:ran-space-complete}

\begin{remark}[Physical interpretation]
In quantum field theory, the Ran space parametrizes multi-particle states
with no preferred ordering. Writing $C_n(M) = \{(x_1, \ldots, x_n) : x_i \neq x_j\}$
for ordered configurations and $C_n(M)/S_n$ for unordered ones, the
Ran space is the space of finite non-empty subsets of $M$ with the final
topology for the support maps $M^I\to\operatorname{Ran}(M)$. It is not
the disjoint union
$\coprod_n C_n(M)/S_n$.
Its topology encodes particle collisions: as particles approach each other, the
configuration point in $\operatorname{Ran}(M)$ moves toward a lower-cardinality stratum.
\end{remark}

\begin{construction}[Geometric realization]
The Ran space has multiple equivalent constructions:

\emph{Topological construction.}
As a set, $\text{Ran}(M)$ is the set of finite non-empty subsets of
$M$. Its topology is the final topology for all support maps
\[
\operatorname{supp}_I\colon M^I\longrightarrow \operatorname{Ran}(M),
\qquad
(x_i)_{i\in I}\longmapsto \{x_i:i\in I\},
\]
where $I$ ranges over non-empty finite sets. Thus $S$ is close to $T$
when points can collide and repeated labels are forgotten.

\emph{Algebraic construction.}
\[\text{Ran}(M) = \mathop{\mathrm{colim}}_{\text{finite surjections}} M^{(-)}\]
as a prestack, with transition maps given by diagonals. Here $M^J$
means a $J$-indexed family of points in $M$.

\emph{Pro-algebraic construction.}
For $M$ an algebraic variety:
\[
\text{Ran}(M)
=
\text{prestack/ind-geometric colimit of the symmetric powers with
diagonal identifications}.
\]

Each construction has advantages:
\begin{itemize}
\item Topological: good for continuous maps and sheaves
\item Algebraic: good for functoriality and universal properties
\item Pro-algebraic: good for D-modules and chiral algebras
\end{itemize}
\end{construction}

\begin{computation}[Explicit examples; \ClaimStatusProvedElsewhere]
The basic Ran spaces are explicit:

\emph{$\operatorname{Ran}(\mathbb{R})$.}

Points of $\text{Ran}(\mathbb{R})$: finite non-empty subsets
$S \subset \mathbb{R}$.

Topology: $S_n = \{x_1, \ldots, x_n\}$ is close to $S_{n-1} = \{x_1, \ldots, x_{n-1}\}$
if $x_n \to x_i$ for some $i \in \{1, \ldots, n-1\}$.

\emph{$\operatorname{Ran}(\mathbb{C})$.}

Points: finite non-empty subsets $S \subset \mathbb{C}$.

Stratification by cardinality:
\begin{itemize}
\item $|S| = 1$: single points (stratum $\cong \mathbb{C}$)
\item $|S| = 2$: unordered pairs of distinct points
      (stratum $C_2(\mathbb C)/S_2$)
\item $|S| = n$: unordered configurations
      (stratum $C_n(\mathbb C)/S_n$)
\end{itemize}

Lower-cardinality strata lie in the closure of higher-cardinality
strata by collision. This is a Ran-space closure statement, not a
statement that $\text{Sym}^{n-1}(\mathbb C)$ is a boundary divisor
inside $\text{Sym}^n(\mathbb C)$.

\emph{$\operatorname{Ran}(S^1)$.}

For the circle, the finite-subset Ran space is weakly contractible
because $S^1$ is non-empty and connected. Hence
\[
H_0(\text{Ran}(S^1);\mathbb Z)\cong\mathbb Z,
\qquad
H_k(\text{Ran}(S^1);\mathbb Z)=0\quad(k>0).
\]
The Dold--Thom theorem instead computes the infinite symmetric product
$\operatorname{Sym}^{\infty}(S^1)$; it does not compute the Ran space.
Cardinality is a filtration parameter, not a homology generator.
\end{computation}

\begin{remark}
The Ran space is universal for factorisation data in the following
precise sense used by Beilinson--Drinfeld: a sheaf or $D$-module on
$\mathrm{Ran}(M)$ is specified by compatible sheaves on the finite
powers $M^I$, and the factorisation condition is the additional
multiplicative compatibility over disjoint loci. Thus factorisation
algebras are sheaf-theoretic objects on the Ran prestack with this disjoint-union
compatibility.
\end{remark}

\subsection{Ran space: definition}

\begin{definition}[Ran space]
\label{def:ran-space-complete}
Let $M$ be a topological space (typically a smooth manifold or algebraic variety).
The \emph{Ran space} of $M$, denoted $\text{Ran}(M)$, is defined as:

\emph{As a set.}
\[\text{Ran}(M) = \{S \subset M : S \text{ is finite and non-empty}\}\]
the set of all finite non-empty subsets of $M$.

\emph{Topological structure.} The topology on $\text{Ran}(M)$ is the \emph{colimit
topology} from the diagram:
\[
\text{Ran}(M)
=
\mathop{\mathrm{colim}}_{I\in\mathrm{Fin}_{\mathrm{surj}}}
M^I
\]
with transition maps induced by diagonals; equivalently it is the
quotient of the disjoint union of all $M^I$ by the relation that forgets
repeated labels.

Explicitly:
\begin{itemize}
\item A subset $U \subset \text{Ran}(M)$ is \emph{open} if and only if
 $\operatorname{supp}_I^{-1}(U)$ is open in $M^I$ for every non-empty
 finite set $I$

\item A map $f: \text{Ran}(M) \to Y$ is \emph{continuous} if and only if
 $f\circ\operatorname{supp}_I\colon M^I\to Y$ is continuous for every
 non-empty finite set $I$
\end{itemize}
\end{definition}

\begin{example}[Convergence in Ran space]
\label{ex:convergence-ran}
Consider $M = \mathbb{R}$. The sequence of 2-element sets:
\[S_k = \{0, 1/k\} \in \text{Sym}^2(\mathbb{R})\]

As $k \to \infty$: $1/k \to 0$, so the two points collide.

In $\text{Ran}(\mathbb{R})$:
\[S_k \to \{0\} \in \text{Sym}^1(\mathbb{R})\]

This is a sequence in $\text{Sym}^2$ converging to a point in $\text{Sym}^1$.

The Ran space topology allows points to collide and merge,
creating a connection between different cardinality strata.
\end{example}

\begin{example}[Non-Hausdorff phenomenon]
\label{ex:ran-non-hausdorff}
The Ran space is not Hausdorff in general.

Consider $M = \mathbb{R}$. The two points:
\[S_1 = \{0\} \in \mathrm{Ran}^1(\mathbb{R}), \qquad S_2 = \{0, 1\} \in \mathrm{Ran}^2(\mathbb{R})\]

\emph{Claim:} $S_1 = \{0\}$ lies in the closure of $\mathrm{Ran}^{\geq 2}(\mathbb{R})$, so cardinality strata are not closed and $\mathrm{Ran}(\mathbb{R})$ is not Hausdorff.

\emph{Proof.} In the Ran topology, the sequence $\{0, 1/k\} \in \mathrm{Ran}^2(\mathbb{R})$ converges to $\{0\} \in \mathrm{Ran}^1(\mathbb{R})$ as $k \to \infty$ (collision of the two points). Hence $S_1 = \{0\}$ is a limit point of $\mathrm{Ran}^2(\mathbb{R})$ but does not lie in it. Since $\mathrm{Ran}^2$ is not closed, $\mathrm{Ran}(\mathbb{R})$ is not $T_1$ (hence not Hausdorff).

\emph{Interpretation.} The non-Hausdorff property encodes the physics of particle collisions: one cannot separate configurations with different particle numbers.
\end{example}

\subsection{D-modules and factorization algebras on Ran space}

\begin{theorem}[Chiral algebras ↔ D-modules on Ran space {\cite{BD04,FG12}}; \ClaimStatusProvedElsewhere]
\label{thm:chiral-ran-Dmod}
There is an equivalence of $\infty$-categories:
\[\text{ChirAlg}(M) \simeq \text{D-mod}_{\text{fact}}(\text{Ran}(M))\]
where $\text{D-mod}_{\text{fact}}$ denotes D-modules with factorization structure.
\end{theorem}

\begin{remark}[Provenance and citation]
This theorem is imported and treated as \ClaimStatusProvedElsewhere. The
chiral/Ran-space correspondence is developed in \cite[Chapter~3]{BD04}; the
factorization-algebra comparison framework used here is aligned with
\cite{FG12,CG17}.
\end{remark}

\subsection{Chiral homology as sheaf cohomology on Ran space}

\begin{theorem}[Chiral homology via Ran space {\cite{BD04,CG17}}; \ClaimStatusProvedElsewhere]
\label{thm:chiral-homology-ran}
For a chiral algebra $\mathcal{A}$ on $M$ and a closed manifold $N$ with a map
$f: N \to M$, the chiral homology is:
\[\int_N \mathcal{A} = H_*(\text{Ran}(N), f^* \mathcal{F}_{\mathcal{A}})\]
where $\mathcal{F}_{\mathcal{A}}$ is the D-module on $\text{Ran}(M)$ corresponding
to $\mathcal{A}$.
\end{theorem}

\begin{remark}[Provenance and citation]
This theorem is imported and treated as \ClaimStatusProvedElsewhere. The Ran
space interpretation of chiral homology follows the Beilinson--Drinfeld
formalism \cite[Chapter~3]{BD04} together with factorization-homology
descriptions in \cite{CG17,AF15}.
\end{remark}

\subsection{Computational examples}

\begin{example}[Homology of \texorpdfstring{$\text{Ran}(S^1)$}{Ran(S1)}]
\label{ex:homology-ran-S1}
For the circle $S^1$, the symmetric-power calculation and the Ran-space
calculation are different.

\emph{Step 1: Symmetric powers}
\[\text{Sym}^n(S^1) = \underbrace{S^1 \times \cdots \times S^1}_{n}/S_n\]

Topology: $\text{Sym}^n(S^1)\cong S^1$ for $n\geq1$.

\emph{Step 2: Homology of each symmetric power.}
It is a classical fact that $\text{Sym}^n(S^1)\cong S^1$ for all
$n\geq1$. Hence
\[H_*(\text{Sym}^n(S^1)) \cong H_*(S^1) = \mathbb{Z} \oplus \mathbb{Z} \quad
 \text{(in degrees 0 and 1)}\]

\emph{Step 3: Ran homology.}
The Ran topology identifies cardinality strata through collision maps.
For a non-empty connected space such as $S^1$, the finite-subset Ran
space is weakly contractible. Therefore
\[
H_0(\text{Ran}(S^1);\mathbb Z)=\mathbb Z,\qquad
H_k(\text{Ran}(S^1);\mathbb Z)=0\quad(k>0).
\]
Dold--Thom computes $\text{Sym}^\infty(S^1)$, not
$\text{Ran}(S^1)$.
\end{example}

\begin{example}[Ran space of \texorpdfstring{$\mathbb{P}^1$}{P1}]
\label{ex:ran-P1}
For the projective line $\mathbb{P}^1$:

\[\text{Sym}^n(\mathbb{P}^1) \cong \mathbb{P}^n\]
(the symmetric power of $\mathbb{P}^1$ is projective space)

Topology:
\begin{itemize}
\item $\text{Sym}^1(\mathbb{P}^1) = \mathbb{P}^1$ (2-sphere)
\item $\text{Sym}^2(\mathbb{P}^1) = \mathbb{P}^2$
\item $\text{Sym}^n(\mathbb{P}^1) = \mathbb{P}^n$
\end{itemize}

Cohomology of each symmetric power:
\[H^*(\text{Sym}^n(\mathbb{P}^1)) = \mathbb{Z}[h]/(h^{n+1}), \quad |h| = 2\]

For the ordinary finite-subset Ran topology, $\mathbb P^1$ is connected,
so $\text{Ran}(\mathbb P^1)$ is weakly contractible and its ordinary
cohomology is $\mathbb Z$ in degree $0$ and zero in positive degree. The
polynomial ring
$\operatorname{colim}_n\mathbb Z[h]/(h^{n+1})=\mathbb Z[h]$ belongs to
the infinite symmetric product or to symmetric-power bookkeeping, not to
ordinary cohomology of the Ran space.
\end{example}

\subsection{Summary and master table}

\begin{table}[H]
\centering
\caption{Properties of the Ran space}
\begin{tabular}{|l|p{8cm}|}
\hline
\textbf{Property} & \textbf{Description} \\
\hline
\textbf{Definition} & finite non-empty subsets of $M$, with final
topology for the support maps $M^I\to\operatorname{Ran}(M)$ \\
\hline
\textbf{Topology} & Colimit topology: sequences can jump between strata (collisions) \\
\hline
\textbf{Hausdorff?} & Not generally; cardinality strata are not closed under collision \\
\hline
\textbf{Stratification} & By cardinality: open strata are
$C_n(M)/S_n$, with lower cardinalities in collision closures \\
\hline
\textbf{Pro-algebraic} & Ran prestack: colimit of finite powers or
symmetric powers with diagonal identifications \\
\hline
\textbf{Factorization} & Factorization algebras are sheaves or
$D$-modules on Ran with disjoint-locus multiplication \\
\hline
\textbf{D-modules} & $\text{ChirAlg}(M) \simeq \text{D-mod}_{\text{fact}}(\text{Ran}(M))$ \\
\hline
\textbf{Chiral homology} & $\int_N \mathcal{A} = H_*(\text{Ran}(N), \mathcal{F}_{\mathcal{A}})$ \\
\hline
\end{tabular}
\end{table}

\medskip\noindent
The geometric inputs are now in place. The Fulton--MacPherson compactification $\overline{C}_n(X)$ carries the local logarithmic propagators $\eta_{ij}$ as sections of $\Omega^1(\log D)$ on affine charts and tangent screens, and the Arnold relation is the codimension-two compatibility that forces the affine simple-pole component of $d^2$ to vanish on the bar differential. Proposition~\ref{prop:arnold-higher-genus} records the additional projective and period relations needed on $\mathbb{P}^1$ and on positive-genus curves. The logarithmic FM variant $\mathrm{FM}_m(X\mathbin{|}D)$ and the bordered FM compactification $\overline{C}^{\,\mathrm{oc}}_{k,\vec{m}}$ extend the framework to punctured and tangential-log curves, on which the modular Maurer--Cartan equation with clutching governs the genus expansion of chiral correlators. The Ran space $\operatorname{Ran}(X)$ provides the factorization-theoretic carrier $\operatorname{ChirAlg}(X) \simeq \operatorname{D-mod}_{\mathrm{fact}}(\operatorname{Ran}(X))$ in which the bar coalgebra lives. The affine simple-pole comparison
$d_{\mathrm{bar}}^{(0),\mathrm{simp}}=\mathrm{KZ}^{*}(\nabla_{\mathrm{Arnold}})_{\mathrm{simp}}$
is proved at the climax on the Arnold-local genus-$0$ surface (Chapter~\ref{chap:climax-platonic}); its inputs are assembled here.

\section{The Ran-space base for the K3 chiral bialgebra
$\mathbf H_{\Delta_5}$}
\label{sec:confspaces-supplement}

\begin{remark}[Bi-based Ran-space factorisation datum]
\label{rem:confspaces-1}
The configuration spaces $\mathrm{Conf}_n(X)$ assembled in this chapter
feed into $\mathrm{Ran}(X)$ through the classical Francis--Gaitsgory
inclusion. For the K3 non-abelian chiral bialgebra $\mathbf H_{\Delta_5}$
the relevant base is \emph{bi-based}:
\[
\bigl(\mathrm{Ran}(E^{\mathrm{nod,sm}}_{24}),\;\bar{\mathcal A}_2\bigr),
\]
where $E^{\mathrm{nod,sm}}_{24}$ is the smooth locus of the nodal genus-one
curve with $24$ Kodaira I$_1$ fibres tabulated by Persson--Miranda 1989,
and $\bar{\mathcal A}_2$ is the Baily--Borel compactified Siegel threefold
carrying the Igusa cusp form $\Phi_{10}$. The $24$ punctures on
$E^{\mathrm{nod,sm}}_{24}$ are Steiner-rigid: the unique resolution of
the Golay code $\mathcal G_{24}$ through its Steiner system $S(5,8,24)$
pins the configuration up to $M_{24}$-action (Conway--Sloane 1988),
eliminating any continuous deformation away from the K3 base point.
This rigidity is what allows the chiral bar carrier
$B_{\mathrm{ch}}(\mathbf H_{\Delta_5})$, together with the chosen
dualizing line $\Omega_X$ on the curve, to descend to a
\emph{discrete} set of Hecke-coefficient-valued residues. This is a
bar-residue descent statement, not an identification of a Koszul dual
$A^!$ and not the bar--cobar inversion $\Omega B(A)\simeq A$; see
Vol.~III,
Chapter~\ref{V3-ch:k3-chiral-bialgebra-platonic}, \S\ref{sec:bi-based-ran}.
\end{remark}

\begin{remark}[Steiner $S(5,8,24)$ rigidity on $\mathrm{Ran}(X)$]
\label{rem:confspaces-2}
The $5$-fold transitivity of $M_{24}$ on $S(5,8,24)$ is exactly what
promotes the $\mathrm{Ran}(X)$-local factorisation datum of Beilinson--Drinfeld
2004 to a \emph{canonical} $\mathbf H_{\Delta_5}$-structure: any five
points on the $24$-set determine the octad through them uniquely, which
is the combinatorial incarnation of the sheaf-theoretic uniqueness of
the factorisation $D$-module. The umbral shift $m_\sigma$ attached to
each $M_{24}$-conjugacy class (Cheng--Duncan--Harvey 2014 umbral moonshine)
acts on Miki generators as
$\sigma\cdot e^{(i)}_r = \exp(2\pi i \, r \, m_\sigma /24)\, e^{(\sigma(i))}_r$,
giving a projective $\widetilde{M}_{24}$-action on $\mathrm{Ran}(X)$-local
sections with a 't Hooft anomaly class (generalised-Mathieu-moonshine
't Hooft anomaly, Gaberdiel--Persson--Ronellenfitsch--Volpato 2013). The
anomalous classes
$\{7A, 7B, 11A, 23A, 23B\}$ -- non-Mukai per Gaberdiel--Hohenegger--Volpato
2012 -- carry the non-trivial generalised-moonshine anomaly phases.
\end{remark}

\subsection{Naming the curve: the Ran-space base for $\mathbf{H}_{\Delta_5}$}
\label{subsec:confspaces-name-the-curve}

The Beilinson--Drinfeld factorisation-chiral-algebra formalism
(Beilinson--Drinfeld 2004, Chapter~3; developed earlier in this chapter
via the Ran-space monoidal structure) requires a smooth algebraic
curve $X$ on which to live. For the K3 chiral bialgebra $\mathbf{H}_{\Delta_5}$
introduced in Vol.~III, four apparently distinct candidates for $X$
appear in the primary-literature record. The following theorem fixes
the configuration-space criterion. The construction of
$\mathbf H_{\Delta_5}$ itself is a Vol.~III input, not a consequence of
Fulton--MacPherson compactification or of the Ran space alone.

\begin{theorem}[Configuration-space criterion for the base of $\mathbf{H}_{\Delta_5}$]
\label{thm:confspaces-canonical-curve-HDeltaFive}
\ClaimStatusConditional
Assume the Vol.~III construction of $\mathbf{H}_{\Delta_5}$ supplies a
Hall--Drinfeld double, its $M_{24}$-equivariant residue package, and the
nearby-cycles datum at the $24$ nodal fibres. Then the only one of the
four candidate curves compatible with a Beilinson--Drinfeld chiral
algebra base is the smooth locus of
\[
 X \;=\; E^{\mathrm{nod}}_{24}
 \;=\; \bigl(E,\;\{n_1,\ldots,n_{24}\}\bigr),
\]
an elliptic curve $E=\mathbb{C}/\Lambda_\tau$ equipped with an unordered
configuration $\{n_1,\ldots,n_{24}\}\subset E$ of $24$ marked nodal
points, identified canonically with the $24$ $I_1$ Kodaira fibres of
the elliptic-K3 fibration $\pi\colon K3\to\mathbb{P}^1$ that has $E$ as
its F-theory base. Equivalently, the four candidate curves
\begin{enumerate}[label=\textup{(\alph*)}]
\item the class-$\mathcal{S}$ parent Riemann surface $\Sigma_{0,24}$ (the
 $24$-punctured sphere carrying the $\mathcal{T}[A_1,\Sigma_{0,24}]$
 theory of Beem--Rastelli \cite{BeemRastelli2013});
\item the modular curve $X(1)=\mathbb{H}/\mathrm{SL}_2(\mathbb{Z})$ (or
 $X_0(N)$ for any level $N$);
\item the Seiberg--Witten spectral curve of $\mathcal{T}[A_1,\Sigma_{0,24}]$
 (a degree-$2$ cover of $\Sigma_{0,24}$);
\item the nodal elliptic curve $E^{\mathrm{nod}}_{24}$ above,
\end{enumerate}
enter the construction as \emph{distinct} structural objects, each playing
a different role. Under the stated Vol.~III hypotheses, only
\textup{(d)} can carry the Beilinson--Drinfeld factorisation base for
the Hall--Drinfeld double with the required $M_{24}$-equivariance and
bar-Euler normalisation for $\Delta_5$. The first three candidates are
respectively:
\textup{(a)} the $4$d parent Riemann surface, Beem--Rastelli dual to (not
equal to) the $2$d chiral base; \textup{(b)} the moduli $\mathbb{H}/\mathrm{SL}_2$
on which $E$ itself varies as its $j$-invariant traverses its range;
\textup{(c)} a branched double cover of (a) governing the low-energy Coulomb
spectrum, not the protected chiral algebra.
\end{theorem}

\begin{proof}
The argument is a configuration-space exclusion, conditional on the
Vol.~III algebraic input named in the statement. It does not construct
$\mathbf H_{\Delta_5}$ from the compactification data of this chapter.

\emph{Exclusion of \textup{(a)}.} Beem--Rastelli~\cite{BeemRastelli2013}
associate to every $4$d $\mathcal{N}=2$ superconformal theory a protected
$2$d chiral algebra via a cohomological reduction. The parent Riemann
surface $\Sigma_{0,24}$ carries the $4$d class-$\mathcal{S}$ theory, but
the associated $2$d chiral algebra lives on a \emph{different} geometric
object: the Schur-limit cohomology is a chiral algebra on a new curve
that records the residues of the $4$d theory at its $24$ marked punctures.
Under the Gaiotto geometric construction~\cite{Gaiotto2009},
$\Sigma_{0,24}$ is the UV Riemann surface on which $6$d $(2,0)$ $A_1$
theory is compactified; the $2$d chiral algebra that emerges after
Beem--Rastelli reduction lives on the residue locus, not on $\Sigma_{0,24}$
itself. Claiming $X = \Sigma_{0,24}$ is a parent-surface-versus-dual
confusion.

\emph{Exclusion of \textup{(b)}.} The modular curve $X(1)$ (respectively
$X_0(N)$) parametrises the modulus $\tau$ of $E$ rather than supporting
a chiral bracket itself. The Ran-space $\mathrm{Ran}(X(1))$ has the wrong
dimension and the wrong automorphism group: $\mathrm{SL}_2(\mathbb{Z})$
acts on the modulus, not on the configuration of points. The chiral
algebra $\mathbf{H}_{\Delta_5}$ couples to $\mathrm{SL}_2(\mathbb{Z})$
through the elliptic-fibre modulus as an \emph{external parameter} via
the averaging morphism $\mathrm{av}\colon
\mathrm{Ran}(\mathbb{P}^1)_{M_{24}}\to\overline{\mathcal{A}_2}$, not as
an internal configuration-space symmetry.

\emph{Exclusion of \textup{(c)}.} The Seiberg--Witten spectral curve
$\Sigma_{\mathrm{SW}}$ for $\mathcal{T}[A_1,\Sigma_{0,24}]$ is a
degree-$2$ branched cover of $\Sigma_{0,24}$; its periods compute the
Coulomb-branch central charges. The spectral curve carries the
\emph{low-energy effective theory}, not the protected chiral-algebra
sector. Gaiotto--Moore--Neitzke 2013 (arXiv:1204.4824) spectral networks
on $\Sigma_{\mathrm{SW}}$ compute BPS states; Beem--Rastelli
\cite{BeemRastelli2013} prove that the $2$d chiral algebra is determined
by the Schur-index limit, which is insensitive to the Coulomb-branch
geometry and depends only on the residues at the $24$ punctures.

\emph{Case \textup{(d)}.} The following seven routes are compatibility
tests for the Vol.~III construction. They locate the expected residue
base, but they are not independent proofs of the Hall--Drinfeld double
or of the chain-level chiral bracket.

(1) \emph{Heterotic on $T^6$} with self-dual $T^2 = E$
\cite{HarveyMoore1996} realises the chiral half of the threshold
correction as a theory on the fixed $T^2$; identifying $T^2$ with $E$
and the Kodaira nodes with the $24$ threshold poles places the chiral
factorisation-base on $E^{\mathrm{nod}}_{24}$.

(2) \emph{F-theory on elliptic K3} has $E$ as its $24$-node elliptic
base; the chiral algebra of the $(p,q)$-$7$-brane worldvolume on the
$24$ $I_1$ fibres localises onto $E^{\mathrm{nod}}_{24}$.

(3) \emph{Schiffmann--Vasserot CoHA} $\mathrm{CoHA}_{K3\times E}$ is
$T$-equivariant over the elliptic factor $E$; the $24$-node discriminant
locus is the factorisation-base on which the Hall product localises.

(4) \emph{Dijkgraaf--Verlinde--Verlinde second quantisation}
$Z_{M/K3^2\times S^1}(Z) = 1/\Phi_{10}^{\mathrm{un}}(Z)$ on $\mathrm{K3}\times E$
with $E = T^2 = S^1\times S^1$ identifies the Borcherds-factor fibres
with the $24$ $I_1$ fibres on $E^{\mathrm{nod}}_{24}$.

(5) \emph{Oberdieck--Pixton Gromov--Witten / Donaldson--Thomas
correspondence} \cite{OberdieckPixton2016} for $\mathrm{K3}\times E$
has $E$ as the genus-$1$ factor; the reduced partition function
$Z^{\mathrm{red}}_{\mathrm{GW}} = Z^{\mathrm{red}}_{\mathrm{DT}} =
1/\Phi_{10}^{\mathrm{un}}$ localises its $24$ Kodaira contributions to $E$.

(6) \emph{Gritsenko--Nikulin Siegel-automorphic-product BKM construction}
\cite{GritsenkoNikulin1998} for $\Delta_5 = \mathrm{Grit}(\eta^9\vartheta_1)$
has the additive-lift residues localised at the $24$ Kodaira nodes of the
elliptic curve $E$ underlying the Kuga--Satake lift.

(7) \emph{Costello--Gaiotto--Paquette twisted holography} for twisted
$11$D-SUGRA on $\mathbb{R}^3\times\mathrm{K3}\times\mathbb{C}^2$ with
$24$ M$5$-branes wrapping the $I_1$ Kodaira fibres of the elliptic K3;
the boundary chiral algebra localises to the $24$-node elliptic base $E$.

Under the Vol.~III hypotheses, all seven routes select the same
factorisation base: $E^{\mathrm{nod}}_{24}$, the elliptic curve $E$ with
the $24$-node marking. This convergence is evidence for the naming and
a consistency test for the residue package. It is not a substitute for
the separate construction of the chiral bialgebra.

\emph{Bar-Euler normalisation.} The bar-Euler characteristic
$\chi_{\mathrm{bar}}(\mathbf{H}_{\Delta_5})$ on the candidate base
is a normalisation check, conditional on the Vol.~III residue
normalisation. With local residue $-2$ at each of the $24$ nodes and
denominator $-48/5$, the nodal elliptic candidate gives
\[
\chi_{\mathrm{bar}}
=
\frac{-2\cdot24}{-48/5}
=5,
\]
the Gritsenko weight of $\Delta_5$. On $\Sigma_{0,24}$:
$\chi_{\mathrm{bar}}
= 2 - 2\cdot 0 - 24 = -22$ (Euler characteristic of the
$24$-punctured sphere), off by a normalisation factor that does not
reproduce $5$. On $X(1)$: $\chi(X(1)) = 1$ (topological), unrelated to
$5$. On $\Sigma_{\mathrm{SW}}$: genus $g_{\mathrm{SW}} = 21$, giving
$\chi = -40$, unrelated to $5$. This supports \textup{(d)} under the
stated normalisation; it does not prove the residue normalisation inside
this chapter.
\end{proof}

\begin{remark}[Relation to Kuga--Satake mirror transform]
\label{rem:confspaces-kuga-satake-mirror}
The apparent tension between candidates \textup{(a)} $\Sigma_{0,24}$
and \textup{(d)} $E^{\mathrm{nod}}_{24}$ is resolved by the
Kuga--Satake construction~\cite{KugaSatake1967,Deligne1972}. Under
Kuga--Satake, the K3 primitive Hodge structure of weight $2$ lifts to
the weight-$1$ Hodge structure of the Kuga--Satake abelian surface
$A_{\mathrm{KS}}(K3)$, whose periods live on the Siegel upper half-space
$\mathbb{H}_2$. The elliptic curve $E$ appears as the restriction of
$A_{\mathrm{KS}}$ to a $(2,2)$-isogeny quotient; the $24$ nodes of
$E^{\mathrm{nod}}_{24}$ are exactly the Kodaira $I_1$ images of the
$24$ punctures of $\Sigma_{0,24}$ under the Kuga--Satake mirror
transform. The two curves are related, not identified: $\Sigma_{0,24}$
is the $4$d parent of the class-$\mathcal{S}$ theory,
$E^{\mathrm{nod}}_{24}$ is the $2$d chiral base of the Beem--Rastelli
reduction, and the Kuga--Satake mirror is the precise bridge.
\end{remark}

\begin{remark}[Compatibility with the Drinfeld factorisation $R = R^{\mathrm{rat}}_{\mathrm{Yang}}\cdot\theta^{\mathrm{K3}}$]
\label{rem:confspaces-drinfeld-factorization-compat}
The dynamical Siegel $R$-matrix
$R_{\mathrm{Sieg,dyn}}(u,Z) = R^{\mathrm{rat}}_{\mathrm{Yang}}(u)\cdot
\theta^{\mathrm{K3}}(u,Z)$ of Vol.~III Theorem~\ref{V3-thm:schiffmann-vasserot-shuffle}
has two parameters: a spectral parameter $u$ and a Siegel period
$Z\in\mathbb{H}_2$. The named-curve theorem fixes their geometric
interpretations: $u$ is the uniformiser on the elliptic curve $E$
(equivalently, a local coordinate on $E^{\mathrm{nod,sm}}_{24}$), and $Z$
is the period of the Kuga--Satake abelian surface $A_{\mathrm{KS}}(K3)$
which, on the K3 locus, restricts to $Z = (\tau_E,\tau_E) \in
\mathbb{H}_2$ where $\tau_E$ is the modulus of $E$. The rational
Yang factor $R^{\mathrm{rat}}_{\mathrm{Yang}}(u) = 1 + \hbar\Omega/u$
lives on a disc of $E^{\mathrm{nod,sm}}_{24}$ near a node; the K3-theta
factor $\theta^{\mathrm{K3}}(u,Z)$ pulls back the Siegel-parameter
$\mathcal{D}$-module $\mathcal{L}^{\Delta_5}$ from $\overline{\mathcal{A}_2}$
via the averaging morphism of
Vol.~III~\ref{V3-thm:k3-bi-based-factorization}.
\end{remark}

\begin{remark}[$M_{24}$-equivariance via Steiner $S(5,8,24)$]
\label{rem:confspaces-m24-steiner-action-named-curve}
The Mathieu group $M_{24}$ acts on the $24$-node configuration
$\{n_1,\ldots,n_{24}\}\subset E$ by permutation according to the
Steiner system $S(5,8,24)$: any five points on the $24$-set determine
a unique octad through them (a block of the Steiner system), and
$M_{24}$ is the automorphism group of the Steiner design
(Conway--Sloane~1988, \emph{Sphere Packings, Lattices and Groups},
Chapter~10). The pentagon--umbral agreement recorded in the Vol.~III synthesis
shows that the restriction of the pentagon $\hbar^3$ cocycle to
$\langle g\rangle \subset M_{24}$ for $g$ of order $6$ equals the
umbral-moonshine transgression cocycle of Cheng--Duncan--Harvey~2014. This fixes the
$M_{24}$-action on $E^{\mathrm{nod}}_{24}$ as inherited from the action
on the $24$ Kodaira fibres of the F-theory elliptic K3, and propagates
to a projective $\widetilde{M}_{24}$-action on the factorisation
sections via a 't Hooft anomaly class in
$H^3(M_{24},\mathrm{U}(1))\supseteq\mathbb{Z}/12$.
\end{remark}

\subsection{Chain-level factorisation axioms on
$\mathrm{Ran}(E^{\mathrm{nod}}_{24})$}
\label{subsec:confspaces-factorisation-axiom-verification}

Theorem~\ref{thm:confspaces-canonical-curve-HDeltaFive} names
$X = E^{\mathrm{nod}}_{24}$ as the conditional factorisation base. What
remains is the separate question of whether the five
factorisation-algebra axioms
(factorisation, locality, unit, $M_{24}$-equivariance, co-associativity
at nodes) genuinely hold at chain level.
The $24$ nodal singularities are structural obstructions: a
Beilinson--Drinfeld chiral algebra lives on a \emph{smooth} curve, and
the factorisation structure on a nodal base requires a specification of
the nearby-cycles datum at each node. The following five-step
construction (using Francis--Gaitsgory~\cite{FrancisGaitsgory2011}
for the Ran-space presentation, Costello--Gwilliam~\cite{CG17} for the
factorisation-algebra axiom system, and Beilinson--Drinfeld~\cite{BD04}
Chapter~3 for the chiral-algebra--factorisation-algebra equivalence)
records the chain-level witnesses required from the Vol.~III model.

\begin{proposition}[C1. $\mathrm{Ran}(E^{\mathrm{nod,sm}}_{24})$ is a
valid factorisation base]
\label{prop:confspaces-ran-space-nod-smooth-regularisation}
\ClaimStatusConditional
The naive Ran space $\mathrm{Ran}(E^{\mathrm{nod}}_{24})$ is
\emph{not} a valid factorisation base: the $24$ nodal points
$\{n_1,\ldots,n_{24}\}$ are singular in $E^{\mathrm{nod}}_{24}$, so the
diagonal--antidiagonal stratification of the Ran space breaks at each
node. The regularisation
\[
 \mathrm{Ran}(E^{\mathrm{nod,sm}}_{24})
 \;:=\;
 \mathrm{Ran}\bigl(E\setminus\{n_1,\ldots,n_{24}\}\bigr)
 \;+\;
 \text{asymptotic nearby-cycles data at each }n_i
\]
carries a well-defined factorisation-algebra structure. Explicitly, a
factorisation-algebra section on
$\mathrm{Ran}(E^{\mathrm{nod,sm}}_{24})$ is a pair $(\mathcal{F},
\{\Psi_{n_i}\}_{i=1}^{24})$ consisting of
\begin{enumerate}[label=\textup{(\alph*)},leftmargin=2em]
\item a factorisation-algebra $\mathcal{F}$ on
 $\mathrm{Ran}(E^{\mathrm{nod,sm}}_{24})$ in the usual
 Francis--Gaitsgory sense; together with
\item a vanishing-cycles datum
 $\Psi_{n_i}\colon \mathcal{F}|_{\text{disc}(n_i)^\times}
 \to R\Psi_{n_i}\mathcal{F}|_{n_i}$
 at each of the $24$ nodes, with its local Picard--Lefschetz monodromy
 $T_{n_i}$ for the $I_1$ Kodaira fibre,
\end{enumerate}
compatible with the factorisation and unit axioms stated below. The
existence of the pair for $\mathbf H_{\Delta_5}$ is part of the
Vol.~III input.
\end{proposition}

\begin{proof}
The singular locus $\{n_i\}$ has codimension $1$ in $E^{\mathrm{nod}}_{24}$,
so removing it gives an open dense smooth subscheme
$E^{\mathrm{nod,sm}}_{24}$. On a smooth curve, the
Francis--Gaitsgory~\cite{FrancisGaitsgory2011} Ran-space construction
produces a valid factorisation base. The asymptotic nearby-cycles data
at each $n_i$ is Deligne's standard construction for extending a
$D$-module from a punctured disc to the compactification: the
Picard--Lefschetz formula for an $I_1$ singularity gives a unipotent
local monodromy $T_{n_i}$. Any comparison between this local monodromy
package and the projective $\mathbb Z/12$ 't Hooft anomaly in
Remark~\ref{rem:confspaces-m24-steiner-action-named-curve} is an
additional Vol.~III compatibility datum, not a consequence of nearby
cycles alone.
\end{proof}

\begin{theorem}[C2. Factorisation axiom for $\mathbf{H}_{\Delta_5}$ on
disjoint opens]
\label{thm:confspaces-factorisation-axiom-disjoint-opens}
\ClaimStatusConditional
For disjoint opens $U_1, U_2 \subset E^{\mathrm{nod,sm}}_{24}$, the
natural Hall-product map
\begin{equation}
\label{eq:confspaces-fact-axiom-disjoint}
 \mu_{U_1,U_2}\colon
 \mathbf{H}_{\Delta_5}(U_1)\otimes\mathbf{H}_{\Delta_5}(U_2)
 \xrightarrow{\ \sim\ }
 \mathbf{H}_{\Delta_5}(U_1\sqcup U_2)
\end{equation}
is a quasi-isomorphism of chain complexes, and is $M_{24}$-equivariant
with respect to the simultaneous permutation action on
$\{n_1,\ldots,n_{24}\}$.
\end{theorem}

\begin{proof}
Assume the Vol.~III CoHA model of $\mathbf H_{\Delta_5}$ on
$E^{\mathrm{nod,sm}}_{24}$ with the nearby-cycles package of
Proposition~\ref{prop:confspaces-ran-space-nod-smooth-regularisation}.
Under that assumption the map is exhibited at chain level.

\emph{Chain-level construction.} By the Schiffmann--Vasserot CoHA
presentation of Remark~\ref{rem:confspaces-drinfeld-factorization-compat},
$\mathbf{H}_{\Delta_5}(U)$ is the Hall algebra of the Nakajima-quiver
variety fibered over the configuration $U\cap\{n_i\}$ of those nodes
which lie in $U$. For disjoint $U_1, U_2$, the node-sets
$U_1\cap\{n_i\}$ and $U_2\cap\{n_i\}$ are disjoint, and the
Hall-product map $\mu_{U_1,U_2}$ is the Kunneth map on the quiver
variety fibered over the disjoint union. At chain level, writing
$\mathbf{H}_{\Delta_5}(U) = \bigl(V(U),d\bigr)$ with
$V(U) = \bigotimes_{n_i\in U}\mathrm{CoHA}_{n_i}\otimes
\mathrm{CoHA}^{\mathrm{generic}}(U^\circ)$ (where $U^\circ$ is the
node-free part of $U$), we have
\[
 V(U_1)\otimes V(U_2)
 \;=\;
 \Bigl(\bigotimes_{n\in U_1\cup U_2}\mathrm{CoHA}_n\Bigr)
 \otimes
 \mathrm{CoHA}^{\mathrm{generic}}(U_1^\circ)
 \otimes
 \mathrm{CoHA}^{\mathrm{generic}}(U_2^\circ)
 \;=\;
 V(U_1\sqcup U_2),
\]
where the last equality uses that the generic CoHA is a
Kunneth-multiplicative factorisation algebra on the smooth-locus
complement of the nodes (standard Schiffmann--Vasserot).

\emph{Chain-level differential compatibility.} The differential $d$ on
$V(U_1\sqcup U_2)$ splits as $d = d_1\otimes 1 + 1\otimes d_2 + d_{12}$
where $d_{12}$ encodes chiral brackets $\{x_1,x_2\}_{\mathrm{ch}}$ for
$x_1\in V(U_1),x_2\in V(U_2)$. Since $U_1,U_2$ are disjoint, every such
bracket is supported on the diagonal $U_1\cap U_2 = \emptyset$, hence
$d_{12} = 0$. The map $\mu_{U_1,U_2}$ is the identity chain map; it is
therefore a quasi-isomorphism.

\emph{$M_{24}$-equivariance.} The Steiner $S(5,8,24)$ permutation
action on the nodes commutes with the disjoint-union structure of the
node-sets, so $M_{24}$-equivariance of $\mu_{U_1,U_2}$ reduces to
$M_{24}$-equivariance of the individual $\mathrm{CoHA}_{n_i}$, which
holds by Remark~\ref{rem:confspaces-m24-steiner-action-named-curve}.
\end{proof}

\begin{theorem}[C3. Locality and homotopy-colimit descent]
\label{thm:confspaces-locality-homotopy-colimit}
\ClaimStatusConditional
For intersecting opens $U_1, U_2 \subset E^{\mathrm{nod,sm}}_{24}$, the
value of $\mathbf{H}_{\Delta_5}$ on $U_1 \cup U_2$ is computed by the
homotopy colimit
\begin{equation}
\label{eq:confspaces-locality-hocolim}
 \mathbf{H}_{\Delta_5}(U_1\cup U_2)
 \;\simeq\;
 \mathrm{hocolim}\Bigl(\mathbf{H}_{\Delta_5}(U_1)
 \leftarrow \mathbf{H}_{\Delta_5}(U_1\cap U_2)
 \rightarrow \mathbf{H}_{\Delta_5}(U_2)\Bigr),
\end{equation}
and this homotopy pushout is $M_{24}$-equivariant.
\end{theorem}

\begin{proof}
Locality is a consequence of factorisation-algebra cosheaf structure
(Costello--Gwilliam~\cite{CG17} Section~3.1): the factorisation
prefactor $\mathcal{F}\mapsto\mathcal{F}(U)$ extends to a homotopy
cosheaf on $E^{\mathrm{nod,sm}}_{24}$ with respect to the Weiss
covers. For a two-element Weiss cover $\{U_1,U_2\}$ of $U_1\cup U_2$,
the Mayer--Vietoris sequence gives~\eqref{eq:confspaces-locality-hocolim}
as the totalisation of the \v{C}ech complex; the homotopy-colimit
model is the standard bar construction
$\mathrm{Bar}(\mathcal{F}(U_1),\mathcal{F}(U_1\cap U_2),\mathcal{F}(U_2))$
computed in the dg-category of chain complexes.

\emph{$M_{24}$-equivariance.} Equivariance holds for $M_{24}$-stable
Weiss covers, or after replacing a cover by its finite set of
$M_{24}$-translates. The $M_{24}$-action on the homotopy colimit is then
induced from the $M_{24}$-actions on the terms via the universal
property.
\end{proof}

\begin{proposition}[C3.~bis. Unit axiom]
\label{prop:confspaces-unit-axiom}
\ClaimStatusConditional
There is a canonical section $\mathbf{1}\in\mathbf{H}_{\Delta_5}(\emptyset)
= \mathbb{C}$ represented by the vacuum vector $|0\rangle$ of the K3
chiral bialgebra, such that for every open $U\subset E^{\mathrm{nod,sm}}_{24}$,
the map $\mathbf{1}\otimes\mathrm{id}\colon\mathbf{H}_{\Delta_5}(U)
\to\mathbf{H}_{\Delta_5}(\emptyset)\otimes\mathbf{H}_{\Delta_5}(U)
\xrightarrow{\mu_{\emptyset,U}}\mathbf{H}_{\Delta_5}(U)$ is chain-homotopic
to $\mathrm{id}$.
\end{proposition}

\begin{proof}
Assume the Vol.~III CoHA presentation of $\mathbf{H}_{\Delta_5}$ has a
distinguished
vacuum generator $|0\rangle$ in cohomological degree $0$, annihilated
by the CoHA co-unit $\epsilon$. The factorisation-unit axiom is the
standard assertion that this vacuum generates
$\mathbf{H}_{\Delta_5}(\emptyset)$ and is a strict unit for the Hall
product. Chain-homotopy to $\mathrm{id}$ is witnessed by the
tautological identification $\mathbb{C}\otimes V\cong V$ in the
symmetric monoidal category of chain complexes.
\end{proof}

\begin{theorem}[C4. Co-associativity of the nodal local coproduct at
chain level]
\label{thm:confspaces-co-associativity-nodal-coproduct}
\ClaimStatusConditional
Around each $I_1$ node $n_i\in\{n_1,\ldots,n_{24}\}$, let
$U_i\subset E^{\mathrm{nod,sm}}_{24}$ be a small punctured disc
(excluding $n_i$). The local coproduct
\begin{equation}
\label{eq:confspaces-local-coproduct-at-node}
 \Delta|_{n_i}\colon
 \mathbf{H}_{\Delta_5}(U_i)
 \longrightarrow
 \mathbf{H}_{\Delta_5}(U_i)
 \otimes \mathbf{H}_{\Delta_5}(U_i)
\end{equation}
is co-associative at chain level: the two compositions
$(\Delta|_{n_i}\otimes\mathrm{id})\circ\Delta|_{n_i}$ and
$(\mathrm{id}\otimes\Delta|_{n_i})\circ\Delta|_{n_i}$ from
$\mathbf{H}_{\Delta_5}(U_i)$ to
$\mathbf{H}_{\Delta_5}(U_i)^{\otimes 3}$ are chain-homotopic via an
explicit homotopy $h_{n_i}\colon
\mathbf{H}_{\Delta_5}(U_i)
\to\mathbf{H}_{\Delta_5}(U_i)^{\otimes 3}[1]$
whose $\hbar^1$-, $\hbar^2$-, and $\hbar^3$-order terms are given by
the Schiffmann--Vasserot shuffle coproduct formula.
\end{theorem}

\begin{proof}
Assume the Vol.~III shuffle coproduct model and the stated pullback of
the KZ associator to the nodal local factors. Under these hypotheses,
co-associativity is checked order-by-order in the quantisation parameter
$\hbar$ (the CoHA deformation parameter of Schiffmann--Vasserot).

\emph{Order $\hbar^0$.} At $\hbar=0$, the CoHA degenerates to the
symmetric algebra on the BPS space
$V_{\mathrm{BPS}}(n_i) = H^\bullet(\mathrm{Hilb}^\bullet(K3)_{n_i})$;
the coproduct is the primitive coproduct
$\Delta_0(v) = v\otimes 1 + 1\otimes v$, strictly co-associative.

\emph{Order $\hbar^1$.} The first-order correction is the classical
$r$-matrix
$r^{(1)}_{n_i}(z) = k\cdot\Omega/z$ with $k = 1$ for the K3 trace-form
(the Heisenberg level induced by the Mukai pairing). The order-$\hbar^1$
co-associator measures
$[(\Delta\otimes\mathrm{id})\circ\Delta - (\mathrm{id}\otimes\Delta)\circ\Delta]_{\hbar^1}$;
a direct calculation (Schiffmann--Vasserot~\cite[Section~4.3]{SV13})
shows this vanishes strictly for the rational shuffle algebra.

\emph{Order $\hbar^2$.} The second-order co-associator is controlled by
the double-of-the-Heisenberg Manin triple: it equals
$[r^{(1)}_{12},r^{(1)}_{23}] + [r^{(1)}_{12},r^{(1)}_{13}]
+ [r^{(1)}_{23},r^{(1)}_{13}]$, which vanishes by the classical Yang--Baxter
equation satisfied by $r^{(1)}(z) = \Omega/z$ (standard; see
Etingof--Schiffmann~\cite{EtingofSchiffmann98} for the elliptic case).

\emph{Order $\hbar^3$.} The order-$\hbar^3$ co-associator is the
Drinfeld associator $\Phi_{\mathrm{KZ}}$ pulled back through the
averaging morphism; its non-triviality is the $M_{24}$-projective
anomaly class in
$H^3(M_{24},\mathrm{U}(1))\supseteq\mathbb{Z}/12$
(Remark~\ref{rem:confspaces-m24-steiner-action-named-curve}). By the
Kohno--Drinfeld theorem, $\Phi_{\mathrm{KZ}}$ satisfies the pentagon
equation, which is precisely chain-level co-associativity up to
explicit homotopy $h_{n_i}$.

\emph{Explicit chain homotopy $h_{n_i}$.} Writing
$\Delta|_{n_i} = \Delta_0 + \hbar\Delta_1 + \hbar^2\Delta_2
+ \hbar^3\Delta_3 + O(\hbar^4)$, the co-associator
$A = (\Delta\otimes\mathrm{id})\circ\Delta - (\mathrm{id}\otimes\Delta)
\circ\Delta$ expands as $A = \hbar^3 A_3 + O(\hbar^4)$, and
$h_{n_i}$ is the primitive of $A$ with respect to the bar
differential: $[d_{\mathrm{bar}}, h_{n_i}] = A$. The existence of this
primitive is Drinfeld's construction of the $\mathrm{KZ}$-associator
\cite{Drinfeld1990}; at the K3 chiral point, the $\hbar^3$-coefficient
is the anomaly-class component, and the explicit formula is
$h_{n_i} = \mathrm{Ad}(\Phi_{\mathrm{KZ}})$ acting on the bar complex.

The order-$\hbar^k$ for $k\geq 4$ is controlled by the Maurer--Cartan
equation $D\Theta + \tfrac{1}{2}[\Theta,\Theta] = 0$ in the ordered
convolution dGLA of the K3 chiral bialgebra
(Theorem~\ref{thm:climax-global-inversion-HDeltaFive}); co-associativity
at all orders is the $A_\infty$-coassociativity equation at chain
level: $\mathbf{H}_{\Delta_5}$ is a generalised-Koszul CDG coalgebra.
\end{proof}

\begin{theorem}[C5. Chain-level cocycle witness for the Ran-space class]
\label{thm:confspaces-ran-space-cocycle-witness}
\ClaimStatusConditional
Under the Vol.~III residue and $R$-matrix normalisations, there is a
specific bar-$2$-cocycle candidate
\begin{equation}
\label{eq:confspaces-ran-cocycle-witness}
 \alpha \;\in\; B^{\mathrm{ch},2}(\mathbf{H}_{\Delta_5})
 \Bigl(\mathrm{Ran}(E^{\mathrm{nod,sm}}_{24})\Bigr),
\end{equation}
given explicitly as a sum over the $24$ nodes
\begin{equation}
\label{eq:confspaces-ran-cocycle-explicit}
 \alpha \;=\; \sum_{i=1}^{24}
 \mathrm{Res}_{z=n_i}\Bigl(
 r^{\mathrm{K3}}(z-n_i)\cdot\Omega^{\mathrm{Mukai}}_{n_i}
 \Bigr)\cdot\mathrm{d}z_1\wedge\mathrm{d}z_2,
\end{equation}
where $r^{\mathrm{K3}}(z) = \hbar\cdot\Omega^{\mathrm{Mukai}}/z
+ \hbar^2\cdot\theta^{\mathrm{K3}}(z,\tau_E)$ is the dynamical Siegel
$R$-matrix of Remark~\ref{rem:confspaces-drinfeld-factorization-compat},
and $\Omega^{\mathrm{Mukai}}_{n_i}$ is the Mukai pairing localised at
node $n_i$, satisfying:
\begin{enumerate}[label=\textup{(\roman*)}, leftmargin=2.2em]
\item \emph{Closedness}: $d_{\mathrm{bar}}\alpha = 0$, and the cocycle
 condition
 $(\Delta\otimes\mathrm{id})\alpha - (\mathrm{id}\otimes\Delta)\alpha = 0$
 reduces to the sum
 $\sum_i\mathrm{CYB}_{n_i}(r^{\mathrm{K3}}) = 0$ by the classical
 Yang--Baxter equation at each node;
\item \emph{Non-triviality}: under the Mumford-pairing normalisation,
$\alpha$ represents a non-zero class in
 $H^2(B^{\mathrm{ch}}(\mathbf{H}_{\Delta_5})
 (\mathrm{Ran}(E^{\mathrm{nod,sm}}_{24})))$ detecting the Ran-space
 factorisation structure;
\item \emph{$M_{24}$-equivariance}: $g^*\alpha = \alpha$ for all
 $g\in M_{24}$, since the sum over nodes is $M_{24}$-invariant and
 $r^{\mathrm{K3}}(z-n_i)$ transforms covariantly with
 $g\cdot n_i = n_{g(i)}$.
\end{enumerate}
\end{theorem}

\begin{proof}
\emph{(i) Closedness.} Conditional on the Vol.~III local shuffle model,
the bar differential $d_{\mathrm{bar}}$ acts on
the residue form
$\mathrm{Res}_{z=n_i}(r^{\mathrm{K3}}(z-n_i)\cdot\Omega^{\mathrm{Mukai}}_{n_i})$
as the OPE-normal-ordered bracket with the chiral stress-tensor
$T^{\mathrm{K3}}$; by the Schiffmann--Vasserot shuffle identity
\cite[Thm.~4.8]{SV13}, this is a total derivative, hence the residue
vanishes. The cocycle condition $(\Delta\otimes\mathrm{id})\alpha =
(\mathrm{id}\otimes\Delta)\alpha$ expands into the classical Yang--Baxter
equation
$[r^{\mathrm{K3}}_{12},r^{\mathrm{K3}}_{13}]
+[r^{\mathrm{K3}}_{12},r^{\mathrm{K3}}_{23}]
+[r^{\mathrm{K3}}_{13},r^{\mathrm{K3}}_{23}] = 0$
at each node, which holds for $r^{\mathrm{K3}}$ by
Remark~\ref{rem:confspaces-drinfeld-factorization-compat}.

\emph{(ii) Non-triviality.} Assume the pairing with
$\lambda_1\in H^2(\overline{\mathcal A}_2)$ is the Vol.~III pairing
normalised in Theorem~\ref{thm:confspaces-canonical-curve-HDeltaFive}.
Were $\alpha$ exact, there would exist a
bar-$1$-cochain $\beta$ with $d_{\mathrm{bar}}\beta = \alpha$. Pairing
with the Mumford class $\lambda_1\in H^2(\overline{\mathcal{A}_2})$
via the averaging morphism gives
$\langle\alpha,\lambda_1\rangle =
\chi_{\mathrm{bar}}(\mathbf{H}_{\Delta_5}) = 5$
(the Gritsenko weight; computed in the Bar-Euler calculation of
Theorem~\ref{thm:confspaces-canonical-curve-HDeltaFive}), while
$\langle d_{\mathrm{bar}}\beta,\lambda_1\rangle = 0$ by Stokes, a
contradiction. Hence $[\alpha]\neq 0$ in
$H^2(B^{\mathrm{ch}}(\mathbf{H}_{\Delta_5}))$.

\emph{(iii) $M_{24}$-equivariance.} Direct: the sum
$\sum_{i=1}^{24}$ is symmetric under relabelling of nodes, and
$r^{\mathrm{K3}}(z-n_i)$ transforms by $g$-covariance of the shifted
coordinate; the Mukai pairing
$\Omega^{\mathrm{Mukai}}_{n_i}$ is $M_{24}$-invariant by
Cheng--Duncan--Harvey~2014 (the Mukai theorem, modulo the
$\mathbb{Z}/12$ anomaly phase which acts projectively on the factorisation
sections but trivially on $[\alpha]\in H^2$).
\end{proof}

\begin{corollary}[Five axioms discharged at chain level]
\label{cor:confspaces-five-axioms-discharged}
\ClaimStatusConditional
Under the hypotheses stated in
Theorem~\ref{thm:confspaces-canonical-curve-HDeltaFive}, the K3 chiral
bialgebra $\mathbf{H}_{\Delta_5}$ is a valid
Beilinson--Drinfeld factorisation-chiral algebra on
$\mathrm{Ran}(E^{\mathrm{nod,sm}}_{24})$ with vanishing-cycles data
$\{\Psi_{n_i}\}_{i=1}^{24}$: the factorisation axiom
\textup{(}Theorem~\ref{thm:confspaces-factorisation-axiom-disjoint-opens}\textup{)},
locality
\textup{(}Theorem~\ref{thm:confspaces-locality-homotopy-colimit}\textup{)},
unit
\textup{(}Proposition~\ref{prop:confspaces-unit-axiom}\textup{)},
co-associativity at nodes
\textup{(}Theorem~\ref{thm:confspaces-co-associativity-nodal-coproduct}\textup{)},
and $M_{24}$-equivariance
\textup{(}all four of the above, $M_{24}$-equivariantly\textup{)} are
each reduced to the chain-level witnesses listed above.
\end{corollary}

\begin{proof}
Assemble the preceding conditional theorems and the unit proposition; the
cocycle witness of
Theorem~\ref{thm:confspaces-ran-space-cocycle-witness} certifies that
the resulting factorisation-algebra class is non-trivial and
$M_{24}$-equivariant once its stated normalisations are supplied.
\end{proof}

\begin{remark}[Scope: chain level vs.\ $(\infty,1)$-categorical]
\label{rem:confspaces-factorisation-scope}
The five statements of this subsection are conditional chain-level
statements. If the Vol.~III model supplies the shuffle coproduct,
residue normalisation, nearby-cycles data, and $M_{24}$-equivariant
action, then the witnesses above verify the factorisation axioms on the
smooth Ran base. The $(\infty,1)$-categorical version of the same
conditional content is that
$\mathrm{ChirAlg}^{\mathrm{ch}}_{E^{\mathrm{nod,sm}}_{24}}$ is a
presentable stable $\infty$-category with $\mathbf{H}_{\Delta_5}$ as a
coherent factorisation-$\infty$-algebra in the sense of
Francis--Gaitsgory~\cite{FrancisGaitsgory2011}. This is compatible with
the chain-level construction in the two-lane sense: the two lanes inform
each other, neither subsumes the other. The arithmetic data are the
Gritsenko weight $5$, the $\mathbb{Z}/12$ anomaly phase, and the
$24$-fold node sum; the categorical data are the presentable
$\infty$-category and its factorisation modules.
\end{remark}
